\chapter{由相容的有限观测生成随机过程}\label{ch:06}

在任意有限多个时刻写出联合分布，并不等于已经构造出一条随机路径。有限维分布必须在删除、
重排和重复坐标时彼此相容；扩张后的典范坐标过程也未必具有连续路径或 \cadlagPath{}。即使路径已经
选好，信息随时间增长以后，条件平均、停时和“从随机时刻重新开始”仍会提出新的可测性与一致可积
问题。过程理论的第一课，因而不是记住一串过程名称，而是分清这些存在机制各自解决什么。

相容的有限维分布首先经 Kolmogorov 扩张成为典范空间上的概率；增量矩估计随后把这个只有
坐标可测性的过程换成路径正则的版本。一旦信息被组织成滤过，鞅将时间一致的条件平均与停时联系起来。
Markov 核则以另一种方式压缩过去：当前状态经转移核迭代为半群，在有限链和 Poisson 过程中可以完全展开计算。
最后，协方差核 $s\wedge t$ 产生 Brownian motion；它的二次变差促使积分理论从经典的有限变差路径转向
$L^2$ 等距，并由此得到 \Ito{} 公式、随机微分方程与热方程表示。

这些主题的共同点是对结论层级的区分。有限维分布不自动给出连续路径；确定时刻的条件期望恒等式不自动
适用于无界停时；局部鞅项也只有在具备相应可积性时才能取期望。后文的构造、逼近与收敛定理正是在这些层级之间
建立可验证的过渡。

\section{过程的存在、版本与路径空间}\label{sec:06-01}

\subsection{随机过程、有限维分布与典范路径空间}\label{subsec:6-1-1}
\index{随机过程}\index{有限维分布}\index{典范路径空间}

一次观测由一个随机变量描述；一族随时间或空间位置变化的观测则要求一族随机变量。困难并不在于
给这族变量添一个下标，而在于同时回答所有有限组观测的联合问题。实验和计算机都只能读取有限多个
坐标，所以我们通常先知道
\[
  (X_{t_1},\ldots,X_{t_n})
\]
的分布，而不是一条已经放在手边的随机函数。于是首先要分清两件事：有限观测怎样组织成一个
路径空间上的分布，以及这个分布是否保证路径连续。前者是本小节与下一小节的问题；后者必须再等到
增量估计介入。

早期的过程常由随机游走、递推式或随机级数逐个构造。Kolmogorov 的公理化改变了起点：先把所有
有限维分布及其相容关系写清，再把“过程存在”化成一个测度扩张问题。典范空间观点则把坐标映射
本身当作过程，使样本点不再承担任何未经定义的物理含义。

\begin{definition}[随机过程]\label{def:6-1-1-1}
设 $(\Omega,\mathcal F,\Prob)$ 为概率空间，$T$ 为非空指标集，$(E,\mathcal E)$ 为可测空间。
一个取值于 $E$、以 $T$ 为指标的\emph{随机过程}是随机元族
\[
  X=(X_t)_{t\in T},\qquad X_t:(\Omega,\mathcal F)\longrightarrow(E,\mathcal E).
\]
固定 $\omega\in\Omega$ 后，函数 $t\mapsto X_t(\omega)$ 称为 $X$ 在 $\omega$ 处的\emph{样本路径}。
\end{definition}

定义只要求每个固定时刻的坐标可测。它没有要求 $(t,\omega)\mapsto X_t(\omega)$ 联合可测，也没有
要求样本路径关于 $t$ 可测或连续。这不是疏忽：这些正则性都需要额外论证。

\begin{definition}[有限维分布]\label{def:6-1-1-2}
设 $I$ 为非空有限集，$\mathbf t:I\to T$ 为一个有限有标号的时刻族。过程 $X$ 在这些时刻的
\emph{有限维分布}为
\[
  \mu_{\mathbf t}
  =\Law\bigl((X_{\mathbf t(i)})_{i\in I}\bigr).
\]
这是 $(E^I,\mathcal E^{\otimes I})$ 上的概率测度。
特别地，当 $I=\{1,\ldots,n\}$ 时也写作 $\mu_{t_1,\ldots,t_n}$。允许 $\mathbf t$ 取重复值；
这样，坐标的重排、删除和复制可由同一个公式处理。
\end{definition}

单时刻边缘远不足以描述一个过程。下面两个过程甚至在每个时刻都有同一个 Bernoulli 分布。

\begin{example}[随机常函数与独立坐标]\label{ex:6-1-1-2}
令 $T=\mathbb N_0=\{0,1,2,\ldots\}$，并令 $Y$ 为参数 $1/2$ 的 Bernoulli 变量。置 $X_n=Y$。则
\[
  \Prob(X_0=X_1)=1,
  \qquad
  \Law(X_{n_1},\ldots,X_{n_k})
  \text{ 集中在 }\{(0,\ldots,0),(1,\ldots,1)\}.
\]
另一方面，在乘积概率空间 $\{0,1\}^{\mathbb N_0}$ 上令 $Z_n(\omega)=\omega_n$，各坐标独立且具有同一
Bernoulli 分布。此时
\[
  \Prob(Z_0=Z_1)=\frac12,
  \qquad
  \Law(Z_{n_1},\ldots,Z_{n_k})
  =\operatorname{Bernoulli}(1/2)^{\otimes k}
\]
（这里 $n_1,\ldots,n_k$ 互异）。两者的所有一维分布相同，二维分布已经把它们分开。
\end{example}

\begin{definition}[典范路径空间与过程分布]\label{def:6-1-1-3}
令
\[
  \Omega_0=E^T,\qquad \pi_t(\omega)=\omega(t).
\]
使全部坐标映射 $\pi_t$ 可测的最小 $\sigma$-代数记为 $\mathcal E^{\otimes T}$，称为\emph{柱
$\sigma$-代数}。若 $F\Subset T$ 为有限集，记 $\pi_F:E^T\to E^F$ 为限制映射；形如
\[
  \pi_F^{-1}(A),\qquad A\in\mathcal E^{\otimes F},
\]
的集合称为有限坐标柱集。

对过程 $X$，定义路径映射
\[
  \mathbf X:\Omega\longrightarrow E^T,
  \qquad \mathbf X(\omega)(t)=X_t(\omega).
\]
因为每个 $X_t$ 可测，而柱 $\sigma$-代数由坐标映射生成，$\mathbf X$ 对
$\mathcal E^{\otimes T}$ 可测；命题~\ref{proposition:6-1-1-1} 随后给出这一点的完整等价证明。
它的推前概率 $\mathbf X_*\Prob$ 称为 $X$ 的\emph{过程分布}，记作 $\Law(X)$。
\end{definition}

这里的“路径空间”首先只是所有函数的集合。即便 $T$ 有拓扑，柱 $\sigma$-代数的定义也只记录有限
坐标观测；路径正则性还没有加入样本空间。

\begin{example}[离散时间柱事件]\label{ex:6-1-1-1}
当 $T=\mathbb N_0$ 时，路径空间为 $E^{\mathbb N_0}$。若 $A_0,A_3\in\mathcal E$，则
\[
  C=\{\omega:\omega(0)\in A_0,\ \omega(3)\in A_3\}
   =\pi_{\{0,3\}}^{-1}(A_0\times A_3)
\]
是柱事件。它询问两次可以实际执行的观测，而不读取其余无穷多个坐标。
\end{example}

\begin{proposition}[过程映射的可测性]\label{proposition:6-1-1-1}
路径映射 $\mathbf X:(\Omega,\mathcal F)\to(E^T,\mathcal E^{\otimes T})$ 可测，当且仅当每个
$X_t$ 可测。
\end{proposition}

\begin{proof}
若 $\mathbf X$ 可测，则 $X_t=\pi_t\circ\mathbf X$ 是可测映射的复合。反之，设每个 $X_t$ 可测，并令
\[
  \mathcal G=\{A\subset E^T:\mathbf X^{-1}(A)\in\mathcal F\}.
\]
逆像保持补集和可数并，所以 $\mathcal G$ 是 $E^T$ 上的 $\sigma$-代数。对每个 $t\in T$ 和
$B\in\mathcal E$，
\[
  \mathbf X^{-1}(\pi_t^{-1}(B))=X_t^{-1}(B)\in\mathcal F.
\]
因此 $\mathcal G$ 包含生成 $\mathcal E^{\otimes T}$ 的全部单坐标柱集，故
$\mathcal E^{\otimes T}\subset\mathcal G$。
\end{proof}

过程分布的有限坐标推前正是有限维分布。下一定理说明，在柱 $\sigma$-代数上，没有比全部有限维
分布更多的概率信息。

\begin{theorem}[有限维分布的唯一性]\label{theorem:6-1-1-2}
设 $P,Q$ 为 $(E^T,\mathcal E^{\otimes T})$ 上的概率测度。若对每个有限 $F\Subset T$ 都有
\[
  (\pi_F)_*P=(\pi_F)_*Q,
\]
则 $P=Q$。
\end{theorem}

\begin{proof}
令 $\mathcal C$ 为全部有限坐标柱集。若
$C=\pi_F^{-1}(A)$、$D=\pi_G^{-1}(B)$，并置 $H=F\cup G$，则
\[
 C\cap D
 =\pi_H^{-1}\!\left((\pi_F^H)^{-1}(A)\cap(\pi_G^H)^{-1}(B)\right)\in\mathcal C,
\]
其中 $\pi_F^H:E^H\to E^F$ 为坐标限制。故 $\mathcal C$ 是 $\pi$-系统，并且
$\sigma(\mathcal C)=\mathcal E^{\otimes T}$。

定义
\[
 \mathcal D=\{A\in\mathcal E^{\otimes T}:P(A)=Q(A)\}.
\]
因为 $P(E^T)=Q(E^T)=1$，取补集保持 $\mathcal D$；若 $(A_n)$ 两两不交且都属于
$\mathcal D$，则可数可加性给
\[
 P\!\left(\bigcup_nA_n\right)=\sum_nP(A_n)
 =\sum_nQ(A_n)=Q\!\left(\bigcup_nA_n\right).
\]
所以 $\mathcal D$ 是 $\lambda$-系统。题设说明 $\mathcal C\subset\mathcal D$，由
$\pi$--$\lambda$ 定理~\ref{theorem:2-2-1-1}，
$\mathcal E^{\otimes T}=\sigma(\mathcal C)\subset\mathcal D$。
\end{proof}

证明只使用有限柱集生成、可数可加性和 $\pi$--$\lambda$ 定理；它没有使用路径空间的紧性。
这种“有限观测唯一识别全局分布”的力量来自可测生成类，而不是来自一条路径可由有限数据重建。

\begin{proposition}[有限维分布的相容性]\label{proposition:6-1-1-3}
设 $X$ 为过程，$I,J$ 为非空有限集，$\mathbf t:I\to T$，且 $\phi:J\to I$。定义
\[
 q_\phi:E^I\longrightarrow E^J,
 \qquad q_\phi((x_i)_{i\in I})=(x_{\phi(j)})_{j\in J}.
\]
则
\[
  \mu_{\mathbf t\circ\phi}=(q_\phi)_*\mu_{\mathbf t}.
  \tag{6.1.1}\label{eq:6-1-1-functorial-consistency}
\]
双射 $\phi$ 给出重排，单射给出删坐标，一般映射还包括复制坐标。
\end{proposition}

\begin{proof}
对每个 $\omega\in\Omega$，坐标恒等式
\[
 (X_{\mathbf t(\phi(j))}(\omega))_{j\in J}
 =q_\phi((X_{\mathbf t(i)}(\omega))_{i\in I})
\]
成立。对两边的随机元取分布，并使用推前映射的复合律，得到
\eqref{eq:6-1-1-functorial-consistency}。
\end{proof}

相容性是过程存在的必要条件。下一小节的主要任务是证明：在适当的状态空间上，它也足够。

\begin{proposition}[典范实现]\label{proposition:6-1-1-4}
任一过程 $X$ 与概率空间
\[
  (E^T,\mathcal E^{\otimes T},\Law(X))
\]
上的坐标过程 $(\pi_t)_{t\in T}$ 具有相同的全部有限维分布。因此，只依赖过程分布的问题可以先在
典范路径空间上研究。
\end{proposition}

\begin{proof}
对任意有限 $F\Subset T$，有 $\pi_F\circ\mathbf X=(X_t)_{t\in F}$。于是
\[
 (\pi_F)_*\Law(X)
 =(\pi_F)_*(\mathbf X_*\Prob)
 =(\pi_F\circ\mathbf X)_*\Prob
 =\Law((X_t)_{t\in F}).
\]
\end{proof}

同一框架容纳多种索引：$T=\mathbb N_0$ 给离散时间序列，$T\subset\R^d$ 给随机场。Markov 性、Gaussian 性与
Brownian 增量结构都是过程分布的附加条件，而不是“随机过程”定义的不同版本。这个区分使我们可以
先统一解决存在问题，再逐项研究各类分布的特殊结构。

\begin{example}[同一分布的不同实现]\label{ex:6-1-1-3}
分别在两个概率空间上取
\[
 \begin{aligned}
 &([0,1],\Borel([0,1]),m), &&U(u)=u,\\
 &([0,2],\Borel([0,2]),\tfrac12m), &&V(v)=v/2.
 \end{aligned}
\]
二者都服从 $[0,1]$ 上的均匀分布，
却定义在不同样本空间。令 $X_t=U$、$Y_t=V$ 对所有 $t\in T$，便得到两个底层样本点毫无对应关系、
过程分布却相同的过程。若要比较 $X_t(\omega)$ 与 $Y_t(\omega)$，必须另行给出共同实现或耦合；
“同分布”本身不提供这样的逐点比较。
\end{example}

当 $T$ 不可数时，柱 $\sigma$-代数还有一条必须公开的边界。有限坐标生成不允许悄悄执行不可数次
交并。

\begin{proposition}[柱事件只依赖可数多个坐标]\label{proposition:6-1-1-countable-support}
若 $A\in\mathcal E^{\otimes T}$，则存在至多可数的 $J\subset T$，使得只要
$\omega|_J=\eta|_J$，便有
\[
  \omega\in A\quad\Longleftrightarrow\quad\eta\in A.
\]
\end{proposition}

\begin{proof}
令 $\mathcal H$ 为满足结论的全部子集 $A\subset E^T$。若 $A$ 由 $J$ 决定，则 $A^c$ 也由 $J$
决定。若 $A_n$ 由至多可数的 $J_n$ 决定，则 $\bigcup_nA_n$ 由至多可数的
$\bigcup_nJ_n$ 决定。因此 $\mathcal H$ 是 $\sigma$-代数。每个有限柱集由其有限坐标集决定，故
$\mathcal E^{\otimes T}=\sigma(\mathcal C)\subset\mathcal H$。
\end{proof}

\begin{remark}[不可数乘积中的 Borel 边界]\label{remark:6-1-1-countable-coordinates}
取不可数集 $T$ 与离散空间 $E=\{0,1\}$。全零路径 $\mathbf0$ 的单点集
\[
 \{\mathbf0\}=\bigcap_{t\in T}\{\omega:\omega(t)=0\}
\]
在乘积拓扑中是闭集，因而属于该拓扑的 Borel $\sigma$-代数。它却不属于柱
$\sigma$-代数。事实上，若它由某个至多可数 $J\subset T$ 决定，取 $t_0\in T\setminus J$，再把
$\mathbf0$ 的 $t_0$ 坐标改成 $1$，所得路径在 $J$ 上与 $\mathbf0$ 相同，却不属于该单点集，矛盾。

所以当 $T$ 不可数时，不能默认
$\mathcal E^{\otimes T}=\Borel(E^T)$。有限维分布唯一决定柱 $\sigma$-代数上的概率，却不自动决定
一个另行指定的更大 Borel $\sigma$-代数上的取值。已确定概率在柱 $\sigma$-代数上的自身完成化当然也
随之确定；问题在于不能把这个完成化与其他未经构造的路径事件 $\sigma$-代数混为一谈。连续路径空间
必须另行建立，而不能把“所有路径都连续”写成一个未经核验的不可数交。
\end{remark}

\begin{figure}[htbp]
\centering
\begin{tikzpicture}[node distance=13mm and 18mm]
  \node[book strong node] (sample) {原概率空间\\$(\Omega,\mathcal F,\Prob)$};
  \node[book strong node,right=of sample] (path) {典范路径空间\\$(E^T,\mathcal E^{\otimes T})$};
  \node[book warm node,right=of path] (finite) {有限坐标空间\\$E^F$};
  \draw[book accent arrow] (sample)--node[above,book label]{$\mathbf X$}(path);
  \draw[book arrow] (path)--node[above,book label]{$\pi_F$}(finite);
  \draw[book dashed arrow,bend right=22] (sample) to
    node[pos=.50,below=1mm,book label]{$(X_t)_{t\in F}=\pi_F\circ\mathbf X$}(finite);
  \node[book label,above=5mm of path] {$\mathbf X_*\Prob=\Law(X)$};
  \node[book label,above=5mm of finite] {$(\pi_F)_*\Law(X)=\mu_F$};
\end{tikzpicture}
\caption{过程分布位于路径空间；有限维分布是它经有限坐标投影得到的影子。}
\label{fig:6-1-1-law-and-fdd}
\end{figure}

\begin{remark}[四种不同的比较]
同一概率空间上的逐点相等、几乎处处相等、同一有限维分布以及同一过程分布是不同层次的陈述。
在柱 $\sigma$-代数上，全部有限维分布相同等价于过程分布相同；它仍不推出两个过程定义在同一空间，
更不推出它们几乎处处相等。数值模拟检验的通常也是有限坐标统计量，所以即使这些检验全部吻合，
路径连续性仍未得到证明。
\end{remark}

\begin{exercise}
\begin{enumerate}
  \item 直接由定义证明：若两个过程互为逐时刻的 a.s. 版本，则它们的有限维分布相同；指出这里只能对
        有限多个零集取并。
  \item 用式~\eqref{eq:6-1-1-functorial-consistency} 分别写出三坐标分布在交换前两坐标、删除第二坐标
        和复制第一坐标时的公式。
  \item 在 $\{0,1\}^{\N}$ 上再构造一对所有一维分布相同、二时刻分布不同的过程，并明确计算
        $\Prob(X_0=X_1)$。
  \item 证明典范实现的过程分布正是原过程分布，而不仅是具有相同的一维边缘。
\end{enumerate}
\end{exercise}

相容性回答了局部分布之间不能互相矛盾，却还没有构造全局概率。下一步要把柱集上的形式赋值证明为
真正的可数可加测度。


\subsection{Kolmogorov 扩张：由相容有限观测生成过程}\label{subsec:6-1-2}
\index{Kolmogorov 扩张定理}\index{投影相容族}

设想我们给每个有限时刻集 $F\Subset T$ 都指定了一个概率 $\mu_F$。把
$\pi_F^{-1}(A)$ 的“概率”写成 $\mu_F(A)$ 并不困难；困难在于同一个柱集可由不同的 $F$ 表示，
而且有限可加还不等于可数可加。扩张定理恰好处理这两个缺口。

历史上也可以从正泛函出发，沿 Daniell 路线得到扩张。本书采用柱代数路线：在紧状态空间中用
有限维 Radon 内正则性与 Tychonoff 紧性证明从上连续，再调用已经建立的 \Caratheodory{} 扩张；
一般标准 Borel 空间随后编码进紧区间。因而柱代数上的预测度构造不依赖正泛函表示定理。

\begin{definition}[投影相容族]\label{def:6-1-2-1}
对每个非空有限集 $F\Subset T$，设 $\mu_F$ 为
$(E^F,\mathcal E^{\otimes F})$ 上的概率测度。若 $F\subset G$ 时总有
\[
  (\pi_F^G)_*\mu_G=\mu_F,
  \tag{6.1.2}\label{eq:6-1-2-projective-consistency}
\]
则称 $(\mu_F)_{F\Subset T}$ 为\emph{投影相容族}。这里
$\pi_F^G:E^G\to E^F$ 是坐标限制映射。
\end{definition}

如果从有序时刻组开始，命题~\ref{proposition:6-1-1-3} 还要求置换和重复坐标相容。改用有限子集
$F$ 后，坐标名称已经固定，式~\eqref{eq:6-1-2-projective-consistency} 只需记录删除坐标；两种说法包含
同一份有限观测信息。

\begin{definition}[柱集候选赋值]\label{def:6-1-2-2}
令 $\mathcal A_T$ 为 $E^T$ 上全部有限坐标柱集构成的代数。对
$C=\pi_F^{-1}(A)$，考虑候选赋值
\[
  P_0(C)=\mu_F(A),
  \tag{6.1.3}\label{eq:6-1-2-cylinder-content}
\]
并置 $P_0(E^T)=1$。
\end{definition}

同一柱集可有不同的有限坐标表示，因而式
\eqref{eq:6-1-2-cylinder-content} 是否真正定义集函数，需要先由投影相容性验证。

\begin{lemma}[柱集赋值的良定义与有限可加性]\label{lemma:6-1-2-cylinder-content}
若 $E\ne\varnothing$ 且 $(\mu_F)$ 投影相容，则式
\eqref{eq:6-1-2-cylinder-content} 良定义，并在 $\mathcal A_T$ 上有限可加。
\end{lemma}

\begin{proof}
设
$\pi_F^{-1}(A)=\pi_G^{-1}(B)$，并置 $H=F\cup G$。因为 $E\ne\varnothing$，限制映射
$\pi_H:E^T\to E^H$ 为满射。又
\[
 \pi_F^{-1}(A)=\pi_H^{-1}((\pi_F^H)^{-1}(A)),\qquad
 \pi_G^{-1}(B)=\pi_H^{-1}((\pi_G^H)^{-1}(B)),
\]
所以满射性给出
\[
  (\pi_F^H)^{-1}(A)=(\pi_G^H)^{-1}(B).
\]
由相容性，
\[
 \mu_F(A)
 =\mu_H((\pi_F^H)^{-1}(A))
 =\mu_H((\pi_G^H)^{-1}(B))
 =\mu_G(B),
\]
故赋值与表示无关。

对有限多个两两不交柱集 $C_1,\ldots,C_n$，取包含它们所用全部坐标的有限集 $H$，写成
$C_j=\pi_H^{-1}(A_j)$。$\pi_H$ 满射，所以 $A_1,\ldots,A_n$ 两两不交，并且
\[
 P_0\!\left(\bigcup_{j=1}^nC_j\right)
 =\mu_H\!\left(\bigcup_{j=1}^nA_j\right)
 =\sum_{j=1}^n\mu_H(A_j)
 =\sum_{j=1}^nP_0(C_j).
\]
\end{proof}

由引理~\ref{lemma:6-1-2-cylinder-content}，$P_0$ 现在是 $\mathcal A_T$ 上良定义的有限可加概率；称它为相容族给出的
\emph{柱集预分布}。

若只停在这里，得到的仍是有限可加概率。一般代数上的有限可加赋值未必能扩张为测度。紧状态空间
提供了所需的从上连续性，而且它是在整个乘积的拓扑紧性与每个有限维分布的内正则性之间搭桥。

\begin{lemma}[紧状态空间上的柱预测度]\label{lemma:6-1-2-2}
设 $E$ 为非空紧度量空间并取 $\mathcal E=\Borel(E)$，
$(\mu_F)_{F\Subset T}$ 为投影相容概率族。则柱集预分布
$P_0$ 是 $\mathcal A_T$ 上的预测度，并且它唯一扩张为
$(E^T,\mathcal E^{\otimes T})$ 上的概率测度。
\end{lemma}

\begin{proof}
若 $T=\varnothing$，则 $E^T$ 为单点，$\mathcal A_T=\{\varnothing,E^T\}$，而 $P_0$ 是该点上的
单位质量；结论直接成立。以下设 $T\ne\varnothing$。

由引理~\ref{lemma:6-1-2-cylinder-content}，$P_0$ 已经良定义且有限可加。我们先证明：若
\[
  C_n\in\mathcal A_T,\qquad C_n\downarrow\varnothing,
\]
则 $P_0(C_n)\downarrow0$。

固定 $\varepsilon>0$。写 $C_n=\pi_{F_n}^{-1}(A_n)$。有限乘积 $E^{F_n}$ 仍是紧度量空间；
$\mu_{F_n}$ 是其上的 Borel 概率，因而由定理~\ref{theorem:2-2-2-2} 具有 Radon 内正则性。
可取紧集 $L_n\subset A_n$，使
\[
  \mu_{F_n}(A_n\setminus L_n)<\varepsilon2^{-n}.
\]
令 $K_n=\pi_{F_n}^{-1}(L_n)$。坐标投影关于乘积拓扑连续，所以 $K_n$ 在 $E^T$ 中闭，且
\[
      K_n\subset C_n,\qquad
 P_0(C_n\setminus K_n)<\varepsilon2^{-n}.
 \tag{6.1.4}\label{eq:6-1-2-closed-cylinder-approx}
\]

由 Tychonoff 定理~\ref{theorem:1-2-2-2}，$E^T$ 紧。若每个有限交
$\bigcap_{n=1}^NK_n$ 都非空，则闭集族 $(K_n)$ 具有有限交性质，从而
$\bigcap_{n\ge1}K_n\ne\varnothing$。但 $K_n\subset C_n$，这会给出
$\bigcap_nC_n\ne\varnothing$，与假设矛盾。因此存在 $N$ 使
$\bigcap_{n=1}^NK_n=\varnothing$。

任取 $x\in C_N$。因为 $(C_n)$ 递减，$x\in C_n$ 对 $1\le n\le N$ 成立；有限个
$K_n$ 的交为空，所以至少有一个 $n\le N$ 满足 $x\notin K_n$。于是
\[
 C_N\subset\bigcup_{n=1}^N(C_n\setminus K_n).
\]
有限可加非负集合函数在代数上有限次可次加，结合
\eqref{eq:6-1-2-closed-cylinder-approx} 得
\[
  P_0(C_N)
  \le\sum_{n=1}^NP_0(C_n\setminus K_n)
  <\varepsilon.
\]
由于 $P_0(C_n)$ 单调不增，且 $\varepsilon$ 任意，故 $P_0(C_n)\downarrow0$。

现在设 $(D_n)$ 为 $\mathcal A_T$ 中两两不交的集合，且
$D=\bigcup_nD_n\in\mathcal A_T$。置
\[
 R_N=D\setminus\bigcup_{n=1}^ND_n.
\]
则 $R_N\downarrow\varnothing$，有限可加性给
\[
 P_0(D)=\sum_{n=1}^NP_0(D_n)+P_0(R_N).
\]
令 $N\to\infty$，刚证的从上连续性给
$P_0(D)=\sum_nP_0(D_n)$。所以 $P_0$ 是预测度。

最后对 $P_0$ 应用 \Caratheodory{} 扩张定理~\ref{theorem:2-2-1-0}。因为
$E^T\in\mathcal A_T$ 且 $P_0(E^T)=1$，有限测度唯一性也适用；所得扩张在
$\sigma(\mathcal A_T)=\mathcal E^{\otimes T}$ 上唯一。
\end{proof}

紧性证明中没有把“每个坐标分别紧”相加成一个不可数乘积估计。真正使用的是闭柱集位于同一个
紧空间 $E^T$ 中，从而有限交性质可用。下一步把一般标准 Borel 空间编码进紧区间；编码后的路径
不必声称逐坐标都落在编码集合中。

\begin{definition}[标准 Borel 状态空间]\label{def:6-1-2-3}
一个可测空间若与某个 Polish 空间的 Borel 可测空间可测同构，称为\emph{标准 Borel 空间}。
这里保留的是可测结构，而不是某个指定的完备度量。
\end{definition}

\begin{proposition}[标准 Borel 编码约化]\label{proposition:6-1-2-3}
设 $E$ 为非空标准 Borel 空间。任一投影相容概率族
$(\mu_F)_{F\Subset T}$ 都是某个
$(E^T,\mathcal E^{\otimes T})$ 上概率的全部有限维边缘。
\end{proposition}

\begin{proof}
若 $T=\varnothing$，取单点空间 $E^T$ 上的单位质量即可。以下设 $T\ne\varnothing$。

由定理~\ref{theorem:1-4-2-4}，存在 Borel 集 $B\subset[0,1]$ 与 Borel 同构
$h:E\to B$。记 $i:B\hookrightarrow[0,1]$ 为包含映射，并对有限 $F$ 定义
\[
  \widetilde\mu_F=((i\circ h)^F)_*\mu_F.
\]
推前的复合律与式~\eqref{eq:6-1-2-projective-consistency} 说明
$(\widetilde\mu_F)$ 是 $[0,1]^F$ 上的投影相容族。由
引理~\ref{lemma:6-1-2-2}，存在 $[0,1]^T$ 的柱 $\sigma$-代数上的概率 $Q$，满足
$(\pi_F)_*Q=\widetilde\mu_F$。

取定 $e_0\in E$，并定义
\[
 r:[0,1]\longrightarrow E,
 \qquad
 r(x)=
 \begin{cases}
   h^{-1}(x),&x\in B,\\
   e_0,&x\notin B.
 \end{cases}
 \tag{6.1.5}\label{eq:6-1-2-retraction}
\]
对 $A\in\mathcal E$，集合 $h(A)$ 是 $B$ 的 Borel 子集，从而也是 $[0,1]$ 的 Borel 子集；
而
\[
 r^{-1}(A)=
 \begin{cases}
   h(A),&e_0\notin A,\\
   h(A)\cup B^c,&e_0\in A.
 \end{cases}
\]
故 $r$ 可测。逐坐标映射
$r^T:[0,1]^T\to E^T$ 对柱 $\sigma$-代数可测，因为
$\pi_t\circ r^T=r\circ\pi_t$ 可测。令 $P=(r^T)_*Q$。对有限 $F$，
\[
 (\pi_F)_*P
 =(r^F)_*(\pi_F)_*Q
 =(r^F)_*\widetilde\mu_F
 =\bigl(r\circ i\circ h\bigr)^F_*\mu_F
 =\mu_F,
\]
其中最后一步使用 $r\circ i\circ h=\id_E$。

这个构造没有使用事件
$\{x\in[0,1]^T:x(t)\in B\text{ 对所有 }t\in T\}$。当 $T$ 不可数时，该事件可能不是柱
$\sigma$-代数中的集合；坐标映射 $r^T$ 正是为了避开这一不可数交。
\end{proof}

\begin{theorem}[Kolmogorov 扩张定理]\label{theorem:6-1-2-1}
设 $(E,\mathcal E)$ 为标准 Borel 空间，$T$ 为任意指标集，
$(\mu_F)_{F\Subset T}$ 为投影相容概率族。则在
$(E^T,\mathcal E^{\otimes T})$ 上存在唯一概率测度 $P$，使
\[
  (\pi_F)_*P=\mu_F\qquad (F\Subset T).
\]
\end{theorem}

\begin{proof}
若 $T=\varnothing$，则 $E^T$ 为单点空间，结论由该点上的单位质量给出。以下设 $T\ne\varnothing$；
有限维概率的存在迫使 $E\ne\varnothing$。存在性由
命题~\ref{proposition:6-1-2-3} 得到。若 $P,Q$ 都具有指定有限维边缘，则它们在所有有限柱集上
相同；定理~\ref{theorem:6-1-1-2} 给出 $P=Q$。
\end{proof}

\begin{corollary}[过程存在与分布的对应]\label{corollary:6-1-2-4}
标准 Borel 状态空间上，投影相容有限维分布族与典范路径空间上的过程分布一一对应。具体地，
Kolmogorov 扩张所得概率空间上的坐标过程具有给定有限维分布；反之，每个过程的有限维分布投影相容。
\end{corollary}

\begin{proof}
由定理~\ref{theorem:6-1-2-1}，相容族给出唯一概率 $P$；在
$(E^T,\mathcal E^{\otimes T},P)$ 上取坐标过程即可。反方向是
命题~\ref{proposition:6-1-1-3} 的删坐标情形。两次对应互逆由有限维分布唯一性定理
\ref{theorem:6-1-1-2} 保证。
\end{proof}

下面几项计算展示相容性如何核验。

\begin{example}[任意指标族的独立坐标]\label{ex:6-1-2-1}
给每个 $t\in T$ 指定 $(E,\mathcal E)$ 上的概率 $\nu_t$。对有限 $F$ 令
\[
 \mu_F=\bigotimes_{t\in F}\nu_t.
\]
若 $F\subset G$，对可测矩形 $A=\prod_{t\in F}A_t$ 计算
\[
 \mu_G((\pi_F^G)^{-1}(A))
 =\prod_{t\in F}\nu_t(A_t)\prod_{t\in G\setminus F}\nu_t(E)
 =\mu_F(A).
\]
有限矩形是生成的 $\pi$-系统，有限测度唯一性定理~\ref{theorem:2-2-1-2} 把等式推广到
$\mathcal E^{\otimes F}$。Kolmogorov 扩张于是给出任意指标集上的独立坐标族。
\end{example}

\begin{example}[半正定核生成 Gaussian 过程]\label{ex:6-1-2-2}
这里称一个过程为 Gaussian 过程，是指它的每个有限维向量都联合 Gaussian；等价地，任意有限
线性组合都具有一维 Gaussian 分布，允许方差为零的退化情形。第~\ref{subsec:6-4-1}~小节将
在构造 Brownian motion 前正式重述这一定义。

设 $m:T\to\R$，$K:T\times T\to\R$ 对称，并满足：对任意
$t_1,\ldots,t_n\in T$ 与 $a_1,\ldots,a_n\in\R$，
\[
 \sum_{j,k=1}^na_ja_kK(t_j,t_k)\ge0.
 \tag{6.1.6}\label{eq:6-1-2-positive-kernel}
\]
对有限 $F$，矩阵 $K_F=(K(s,t))_{s,t\in F}$ 半正定。线性代数给出矩阵 $A_F$ 使
$K_F=A_FA_F^T$；取标准正态向量 $Z$，令 $m_F+A_FZ$ 的分布为 $\mu_F$。由
例~\ref{ex:5-4-3-1}，其特征函数为
\[
 \widehat\mu_F(u)
 =\exp\!\left(i\sum_{t\in F}u_tm(t)
 -\frac12\sum_{s,t\in F}u_sK(s,t)u_t\right).
 \tag{6.1.7}\label{eq:6-1-2-gaussian-fdd-cf}
\]
若 $F\subset G$，把 $u\in\R^F$ 在 $G\setminus F$ 上补成零，式
\eqref{eq:6-1-2-gaussian-fdd-cf} 表明 $(\pi_F^G)_*\mu_G$ 与 $\mu_F$ 的特征函数相同。
特征函数唯一性定理~\ref{theorem:5-4-3-2} 给出投影相容性。故存在一个均值为 $m$、协方差核为
$K$ 的 Gaussian 过程。

例如 $T=[0,\infty)$、$K(s,t)=s\wedge t$ 满足
\[
 \sum_{j,k}a_ja_k(t_j\wedge t_k)
 =\int_0^\infty\left(\sum_ja_j\1_{[0,t_j]}(u)\right)^2\dd u\ge0.
\]
扩张定理因此已经给出 Brownian 有限维分布的一个典范过程，但尚未给出连续路径。连续版本要由下一
小节的矩估计得到；Brownian motion 的完整构造见 \S\ref{subsec:6-4-1}。
\end{example}

\begin{example}[互相冲突的二时刻要求]\label{ex:6-1-2-3}
设 $X_0,X_1$ 都应为参数 $1/2$ 的 Bernoulli 变量。若再要求 $X_0=X_1$ a.s.，则二维分布满足
\[
 \Prob((X_0,X_1)=(0,1))=\Prob((X_0,X_1)=(1,0))=0.
\]
若同时要求 $X_0,X_1$ 独立，则这两个概率都必须等于 $1/4$。因此不存在同时满足两项要求的二维
分布，更不可能存在相应过程。检查相容性必须在扩张之前完成；扩张定理不会替我们调解矛盾的局部指定。
\end{example}

\begin{example}[可数状态 Markov 链的初始构造]\label{ex:6-1-2-markov-chain}
设 $E$ 可数，$\nu$ 为 $E$ 上的概率，$p(x,y)\ge0$ 且
$\sum_{y\in E}p(x,y)=1$。在 $E^{\{0,\ldots,n\}}$ 上定义
\[
\mu_{0:n}(x_0,\ldots,x_n)
=\nu(x_0)\prod_{k=1}^np(x_{k-1},x_k).
\tag{6.1.8}\label{eq:6-1-2-markov-fdd}
\]
从 $x_n,x_{n-1},\ldots,x_1$ 依次求和，并逐次使用
$\sum_y p(x,y)=1$，可得
\[
 \sum_{x_0,\ldots,x_n}\mu_{0:n}(x_0,\ldots,x_n)
 =\sum_{x_0}\nu(x_0)=1;
\]
因此这确为概率
分布。对末坐标求和得到
\[
 \sum_{x_n\in E}\mu_{0:n}(x_0,\ldots,x_n)
 =\mu_{0:n-1}(x_0,\ldots,x_{n-1}),
\]
因为 $\sum_{x_n}p(x_{n-1},x_n)=1$。反复删除末坐标后，任意有限
$F\subset\mathbb N_0$ 的分布可定义为 $\mu_{0:N}$ 在 $F$ 上的边缘，其中 $N\ge\max F$；上式保证定义与
$N$ 无关并给出投影相容性。扩张定理产生整个离散时间链。一般状态空间的转移核版本将在本章第三节
正式展开。
\end{example}

有限体积 Gibbs 分布的无限体积极限也有相似外形：大有限体积的分布必须向小体积投影为指定分布。
但实际 Gibbs 一致性通常涉及边界条件和条件核，不是把
定理~\ref{theorem:6-1-2-1} 的符号替换一下便自动完成；这里的类比只指出“局部规格先相容、全局对象后
存在”的共同结构。

\begin{figure}[htbp]
\centering
\begin{tikzpicture}[node distance=11mm and 18mm]
  \node[book strong node] (P) {$E^T$ 上的 $P$};
  \node[book node,below left=of P] (G) {$E^G$ 上的 $\mu_G$};
  \node[book node,below right=of P] (H) {$E^H$ 上的 $\mu_H$};
  \node[book warm node,below=18mm of P] (F) {$E^F$ 上的 $\mu_F$};
  \draw[book accent arrow] (P)--node[left,book label]{$\pi_G$}(G);
  \draw[book accent arrow] (P)--node[right,book label]{$\pi_H$}(H);
  \draw[book arrow] (G)--node[left,book label]{$\pi_F^G$}(F);
  \draw[book arrow] (H)--node[right,book label]{$\pi_F^H$}(F);
  \node[book label,above=2mm of F] {$F\subset G,\ F\subset H$};
\end{tikzpicture}
\caption{相容三角形要求从任一较大有限坐标集投影到 $F$ 都得到同一个 $\mu_F$；扩张定理把整张
有限集合有向图汇聚为路径空间概率 $P$。}
\label{fig:6-1-2-projective-family}
\end{figure}

\begin{remark}[扩张只解决存在，不蕴含路径正则性]
定理~\ref{theorem:6-1-2-1} 得到的是全部函数 $E^T$ 上柱 $\sigma$-代数中的概率。即使每个
有限维分布都有光滑密度，也不能据此推出样本路径连续、可测或有界。另一方面，一般非标准 Borel 状态
空间缺少上述可测编码，扩张结论可能失败或需要额外假设；“可测空间”在定理中不能无条件替换“标准
Borel 空间”。
\end{remark}

\begin{exercise}
\begin{enumerate}
  \item 对例~\ref{ex:6-1-2-1} 完成从可测矩形到全部乘积可测集的 $\pi$--$\lambda$ 论证。
  \item 证明式~\eqref{eq:6-1-2-positive-kernel} 等价于每个有限矩阵 $K_F$ 半正定，并验证常数核
        $K(s,t)=1$ 与核 $K(s,t)=s\wedge t$ 的条件。
  \item 在 $E=\{0,1\}$ 上沿引理~\ref{lemma:6-1-2-2} 的路线构造独立公平 Bernoulli 坐标，指出
        闭柱逼近在有限离散空间中为何可直接取 $K_n=C_n$。
  \item 只给出全部一维边缘 $\Law(X_t)$ 时，构造至少两个不同的相容有限维族，说明它们扩张出的过程
        分布不同。
\end{enumerate}
\end{exercise}

扩张把相容有限观测拼成了过程，却可能拼出极不规则的路径。接下来要从增量的定量控制中选择连续
版本；这是另一个存在问题，不能由有限维唯一性代替。


\subsection{版本、可分性与 Kolmogorov 连续性定理}\label{subsec:6-1-3}
\index{版本}\index{不可分辨}\index{Kolmogorov--Chentsov 定理}

假设对每个固定 $t$，两个过程 $X_t$ 与 $Y_t$ 几乎处处相等。若 $T$ 不可数，把这些概率一事件
全部相交不是测度论允许的可数运算；共同零集可能根本不存在。连续版本理论的作用，正是用可数稠密
网格和一致增量控制制造一个对所有时刻同时有效的事件。

Kolmogorov--Chentsov 方法是“离散估计生成连续路径”的原型：先在逐层加密的二进网格上同时
控制所有相邻增量，再沿格点链求和，最后借目标空间的完备性延拓。它补上的正是扩张定理看不见的
样本函数正则性。

\begin{definition}[版本与不可分辨]\label{def:6-1-3-1}
同一概率空间上、具有同一指标集并取值于同一标准 Borel 空间的两个过程 $X,Y$ 若满足
\[
  \Prob(X_t=Y_t)=1\qquad\text{对每个固定 }t\in T,
\]
则称它们互为\emph{版本}（modifications）。若存在 $N\in\mathcal F$，
$\Prob(N)=0$，使得
\[
  X_t(\omega)=Y_t(\omega)\qquad
  (\omega\notin N,\ t\in T),
\]
则称它们\emph{不可分辨}。
\end{definition}

标准 Borel 空间的对角线可测，所以定义中的逐时刻等式确实给出事件。

后一表述刻意使用一个可测共同零集，而不把不可数交
$\bigcap_{t\in T}\{X_t=Y_t\}$ 未经核验地称作事件。不可分辨蕴含互为版本；反向一般失败。

\begin{example}[单点异常版本]\label{ex:6-1-3-1}
在 $([0,1],\Borel([0,1]),m)$ 上令 $U(\omega)=\omega$，并对 $t\in[0,1]$ 定义
\[
  X_t=\1_{\{U=t\}},\qquad Y_t=0.
\]
对固定 $t$，单点集 $\{t\}$ 的 Lebesgue 测度为零，故 $X_t=Y_t$ a.s.；所以 $X,Y$ 互为版本。
但对每个 $\omega$，取 $t=U(\omega)$ 就有 $X_t(\omega)=1\ne0=Y_t(\omega)$。不存在使两条路径
处处相同的样本点，更不存在概率一的此类样本点。因此它们并非不可分辨。
\end{example}

互为版本的过程具有相同有限维分布：对固定有限时刻组，只需把有限多个坐标等式的零集取并。
因此版本改变不影响柱 $\sigma$-代数上的过程分布，却可能彻底改变逐路径性质。上例中 $Y$ 的每条路径
连续，$X$ 的每条路径都在一个点发生尖峰。

\begin{definition}[可分版本]\label{def:6-1-3-3}
设 $T$ 为度量空间，$X$ 为实过程。若存在可数稠密集 $D\subset T$ 与零集 $N$，使对每个开球
$G\subset T$ 及每个 $\omega\notin N$ 都有
\[
 \sup_{t\in G}X_t(\omega)=\sup_{t\in G\cap D}X_t(\omega),
 \qquad
 \inf_{t\in G}X_t(\omega)=\inf_{t\in G\cap D}X_t(\omega),
 \tag{6.1.9}\label{eq:6-1-3-separable-version}
\]
其中允许取扩展实值，则称该版本相对于 $D$ \emph{可分}。
\end{definition}

连续路径自动满足式~\eqref{eq:6-1-3-separable-version}：对 $t\in G$，从 $G\cap D$ 中取收敛到
$t$ 的序列，再用连续性。单点异常过程则说明任意版本未必可分。事实上，固定可数稠密集 $D$ 后，
$\Prob(U\in D)=0$；当 $U(\omega)\notin D$ 时，任何包含 $U(\omega)$ 的开球 $G$ 都满足
\[
  \sup_{t\in G}X_t(\omega)=1,
  \qquad \sup_{t\in G\cap D}X_t(\omega)=0.
\]

\begin{definition}[路径 \Holder{} 正则性]\label{def:6-1-3-2}
设 $T\subset\R^d$。若过程的一条 $\R^m$ 值路径 $x:T\to\R^m$ 满足
\[
  \|x(t)-x(s)\|\le M\|t-s\|^\gamma\qquad(s,t\in T)
  \tag{6.1.10}\label{eq:6-1-3-holder-path}
\]
对某个有限常数 $M$ 成立，则称该路径为全局 $\gamma$-\Holder{}。若式
\eqref{eq:6-1-3-holder-path} 在每个紧子集上成立而常数可依赖紧集与路径，则称其局部
$\gamma$-\Holder{}。
\end{definition}

连续性定理的证明先在二进网格上进行。令
\[
 D_n=\{0,2^{-n},2\cdot2^{-n},\ldots,1\}^d,
 \qquad D=\bigcup_{n\ge0}D_n.
\]
若 $u,v\in D_n$ 只在一个坐标上相差 $2^{-n}$，称 $\{u,v\}$ 为第 $n$ 层的一条相邻边。

\begin{lemma}[二进相邻边估计]\label{lemma:6-1-3-2}
设 $X_t\in\R^m$，$t\in[0,1]^d$，并且存在 $\alpha,\beta,C>0$ 使
\[
 \E\|X_t-X_s\|^\alpha\le C\|t-s\|^{d+\beta}
 \qquad(s,t\in[0,1]^d).
 \tag{6.1.11}\label{eq:6-1-3-moment-condition}
\]
若 $0<\gamma<\beta/\alpha$，则以概率一存在随机层数 $N$，使每个 $n\ge N$ 及第 $n$ 层的
每条相邻边 $\{u,v\}$ 都满足
\[
  \|X_u-X_v\|\le2^{-n\gamma}.
  \tag{6.1.12}\label{eq:6-1-3-eventual-edge-control}
\]
\end{lemma}

\begin{proof}
第 $n$ 层在第 $j$ 个坐标方向上的边数是
$2^n(2^n+1)^{d-1}$，故全部相邻边不超过
\[
 d\,2^n(2^n+1)^{d-1}\le d\,2^{d-1}2^{nd}.
\]
固定一条边 $\{u,v\}$。由 Markov 不等式~\ref{theorem:5-1-1-2} 和
$\|u-v\|=2^{-n}$，
\[
 \Prob(\|X_u-X_v\|>2^{-n\gamma})
 \le2^{n\alpha\gamma}\E\|X_u-X_v\|^\alpha
 \le C2^{-n(d+\beta-\alpha\gamma)}.
\]
令 $B_n$ 表示第 $n$ 层至少有一条边违反
\eqref{eq:6-1-3-eventual-edge-control}。并集界给
\[
 \Prob(B_n)
 \le Cd2^{d-1}2^{-n(\beta-\alpha\gamma)}.
\]
因为 $\beta-\alpha\gamma>0$，级数 $\sum_n\Prob(B_n)$ 收敛。Borel--Cantelli 引理
\ref{theorem:5-1-3-1} 给出 $\Prob(B_n\ \mathrm{i.o.})=0$，这正是所需的最终逐层控制。
\end{proof}

只控制相邻边还需一次链式计算，才能比较任意两个稠密格点。

\begin{lemma}[二进链估计]\label{lemma:6-1-3-dyadic-chain}
设 $\gamma>0$，函数 $x:D\to\R^m$ 满足：存在 $M<\infty$，使每个 $n$ 及第 $n$ 层相邻边
$\{u,v\}$ 都有
\[
  \|x(u)-x(v)\|\le M2^{-n\gamma}.
\]
则存在只依赖于 $d,\gamma$ 的常数 $A_{d,\gamma}$，使
\[
  \|x(u)-x(v)\|
  \le A_{d,\gamma}M\|u-v\|_\infty^\gamma
  \qquad(u,v\in D).
  \tag{6.1.13}\label{eq:6-1-3-dyadic-chain}
\]
\end{lemma}

\begin{proof}
取不同的 $u,v\in D$，并先设 $\|u-v\|_\infty\le1/2$。选 $n\ge0$ 使
\[
  2^{-(n+1)}<\|u-v\|_\infty\le2^{-n}.
  \tag{6.1.14}\label{eq:6-1-3-scale-choice}
\]
再取 $N\ge n$ 使 $u,v\in D_N$。对每个第 $k$ 层格点，逐坐标舍入到最近的第 $k-1$ 层格点；
所得父点与原点可由至多 $d$ 条第 $k$ 层相邻边连接。由此分别从 $u,v$ 构造父点链
\[
 u=u_N,u_{N-1},\ldots,u_n,
 \qquad
 v=v_N,v_{N-1},\ldots,v_n.
\]
舍入误差和式~\eqref{eq:6-1-3-scale-choice} 表明 $u_n,v_n$ 的每个坐标至多相差三个第 $n$ 层
网格步长，所以它们可由至多 $3d$ 条第 $n$ 层相邻边连接。三角不等式于是给
\begin{align*}
 \|x(u)-x(v)\|
 &\le3dM2^{-n\gamma}
   +2dM\sum_{k=n+1}^N2^{-k\gamma}\\
 &\le dM\left(3+\frac{2}{2^\gamma-1}\right)2^{-n\gamma}\\
 &\le 2^\gamma d\left(3+\frac{2}{2^\gamma-1}\right)
 M\|u-v\|_\infty^\gamma.
\end{align*}
若 $\|u-v\|_\infty>1/2$，在 $D_0$ 处汇合两条父点链，并用同一个几何级数估计；此时
$\|u-v\|_\infty^\gamma\ge2^{-\gamma}$，把所得固定常数再乘 $2^\gamma$ 即覆盖该情形。
\end{proof}

\begin{theorem}[Kolmogorov--Chentsov 连续性定理]\label{theorem:6-1-3-1}
设 $X_t\in\R^m$，$t\in[0,1]^d$，并满足矩条件
\eqref{eq:6-1-3-moment-condition}。则 $X$ 有连续版本。更精确地，对每个
\[
  0<\gamma<\frac{\beta}{\alpha},
\]
$X$ 都有全局 $\gamma$-\Holder{} 版本。

若过程以 $\R^d$（或 $[0,\infty)^d$）为指标，且对每个 $R>0$，式
\eqref{eq:6-1-3-moment-condition} 在参数立方体 $[-R,R]^d$（或 $[0,R]^d$）上成立，指数
$\alpha,\beta$ 固定而常数可依赖 $R$，则存在局部 $\gamma$-\Holder{} 版本，其中仍可取任意
$0<\gamma<\beta/\alpha$。
\end{theorem}

\begin{proof}
先处理参数集 $[0,1]^d$。固定
$0<\gamma<\beta/\alpha$。由
引理~\ref{lemma:6-1-3-2}，存在概率一事件 $\Omega_*$，使
\eqref{eq:6-1-3-eventual-edge-control} 从某个随机层数起成立。对
$\omega\in\Omega_*$，把此前有限多层中所有边的量
\[
  2^{n\gamma}\|X_u(\omega)-X_v(\omega)\|
\]
取最大，并再与 $1$ 取最大，得到有限随机常数 $M(\omega)$。于是每一层的每条相邻边都满足
\[
 \|X_u(\omega)-X_v(\omega)\|\le M(\omega)2^{-n\gamma}.
\]
引理~\ref{lemma:6-1-3-dyadic-chain} 给出
\[
 \|X_u(\omega)-X_v(\omega)\|
 \le A_{d,\gamma}M(\omega)\|u-v\|_\infty^\gamma
 \qquad(u,v\in D).
 \tag{6.1.15}\label{eq:6-1-3-holder-on-D}
\]

对每个 $t\in[0,1]^d$，固定一列 $q_n(t)\in D_n$，使
$q_n(t)\to t$；若 $t\in D$，约定从足够大的 $n$ 起 $q_n(t)=t$。式
\eqref{eq:6-1-3-holder-on-D} 说明 $(X_{q_n(t)}(\omega))_n$ 在完备空间 $\R^m$ 中为 Cauchy 列。
定义
\[
 Y_t(\omega)=
 \begin{cases}
  \displaystyle\lim_{n\to\infty}X_{q_n(t)}(\omega),&\omega\in\Omega_*,\\
  0,&\omega\notin\Omega_*.
 \end{cases}
 \tag{6.1.16}\label{eq:6-1-3-extension-definition}
\]
把 $\Omega_*$ 换成其中可测的概率一事件（由 Borel--Cantelli 证明中的可数格点事件直接给出），
可见 $Y_t$ 是可测随机向量。令二进格点趋近 $s,t$，式
\eqref{eq:6-1-3-holder-on-D} 给出
\[
 \|Y_t(\omega)-Y_s(\omega)\|
 \le A_{d,\gamma}M(\omega)\|t-s\|_\infty^\gamma
 \qquad(\omega\in\Omega_*).
\]
在 $\Omega_*^c$ 上路径为零，所以 $Y$ 的每条路径都是 $\gamma$-\Holder{}。

还需核验 $Y$ 是 $X$ 的版本。固定 $t$。矩条件给
\[
 \E\|X_{q_n(t)}-X_t\|^\alpha
 \le C\|q_n(t)-t\|^{d+\beta}\longrightarrow0.
\]
由 Markov 不等式，$X_{q_n(t)}\to X_t$ 依概率；定义
\eqref{eq:6-1-3-extension-definition} 又给
$X_{q_n(t)}\to Y_t$ a.s.，因而也依概率收敛。依概率极限的唯一性给
$\Prob(Y_t=X_t)=1$。这里的零集可以依赖 $t$，这正符合“版本”的定义。紧参数集情形得证。

最后固定任意 $0<\rho<\beta/\alpha$ 并处理 $\R^d$；半空间情形同理。令
$Q_n=[-n,n]^d$。经过仿射缩放，把已证结论以指数 $\rho$ 应用于
$X|_{Q_n}$，得到 $Q_n$ 上的 $\rho$-\Holder{} 版本 $Y^{(n)}$。对可数稠密集
$D_n=Q_n\cap\Q^d$ 中每个 $t$，$Y_t^{(n)}$ 与 $Y_t^{(n+1)}$ 都是 $X_t$ 的版本。对
$t\in D_n$ 的零集取可数并，再用两条路径在 $Q_n$ 上的连续性，得到一个零集 $N_n$，使
\[
 Y_t^{(n)}(\omega)=Y_t^{(n+1)}(\omega)
 \qquad(\omega\notin N_n,\ t\in Q_n).
\]
再对 $n$ 取可数并，所有局部版本在嵌套立方体上同时吻合。于该概率一事件上将它们拼接，并在
异常集上取零路径，便得到全空间上的版本；每个 $Q_n$ 上的估计给出局部
$\rho$-\Holder{} 性。
\end{proof}

上述论证中，每层 $O(2^{nd})$ 条边正好说明参数维数 $d$ 为何出现在矩条件中；严格不等式
$\alpha\gamma<\beta$ 保证坏事件概率可和。Borel--Cantelli 定理先在可数二进格点集上给出共同零集，
再由目标空间的完备性将估计延拓到全部参数点。因而这个几何级数论证不包含端点
$\gamma=\beta/\alpha$。
若允许的指数 $\gamma>1$，所得 \Holder{} 路径在连通的立方体上必为常函数：把任意线段等分为 $n$
段并求和，差值至多为 $Mn^{1-\gamma}\|t-s\|^\gamma\to0$。这与定理并不冲突，而是矩条件过强时
过程退化的结论。

\begin{proposition}[连续版本的唯一性]\label{proposition:6-1-3-3}
设 $T$ 为可分度量空间，$X,Y$ 取值于同一度量空间。若 $X,Y$ 互为版本，且两者在各自一个概率一
事件上具有连续路径，则
$X,Y$ 不可分辨。
\end{proposition}

\begin{proof}
取可数稠密集 $D\subset T$。对每个 $t\in D$，事件 $\{X_t=Y_t\}$ 有概率一；对这些事件以及
$X,Y$ 路径连续的两个概率一事件取可数交，得到概率一事件 $A$。固定 $\omega\in A$。连续函数
$t\mapsto X_t(\omega)$ 与 $t\mapsto Y_t(\omega)$ 在稠密集 $D$ 上相等，故在 $T$ 上处处相等。
取 $N=A^c$ 即得不可分辨性。
\end{proof}

这个命题也解释了上一证明的局部拼接为什么不能“在每个紧集上随便选一个版本”以后便停止：必须先在
可数稠密集上协调重叠，再由连续性获得同一个零集上的逐点一致。

\begin{proposition}[连续路径空间上的分布]\label{proposition:6-1-3-4}
令 $K=[0,1]^d$，$\mathcal C=C(K;\R^m)$ 配以上确界范数。若一个 $\R^m$ 值过程有连续版本，则可取
一个每条路径都连续的版本 $Y$，使路径映射
\[
  \mathbf Y:\Omega\longrightarrow\mathcal C,
  \qquad \mathbf Y(\omega)=Y_\bullet(\omega)
\]
为 Borel 随机元。并且对任一可数稠密集 $D\subset K$，
\[
  \Borel(\mathcal C)=\sigma(e_t:t\in D),
  \qquad e_t(f)=f(t).
  \tag{6.1.17}\label{eq:6-1-3-path-borel}
\]
因此 $\Law(\mathbf Y)$ 是 $\mathcal C$ 上的 Borel 概率，其有限维分布由连续评价映射推前恢复。
\end{proposition}

\begin{proof}
先在原连续版本的异常零集上把整条路径改成零路径，便可假设每条路径连续，同时不改变任何固定时刻的
版本关系。空间 $\mathcal C$ 是 Banach 空间。它也是可分的：在每一级二进网格上取有理向量值，
再作逐小立方体的多线性插值，所得函数全体可数；连续函数的一致连续性表明这些函数在上确界范数下
稠密。

每个评价映射 $e_t$ 连续，所以
$\sigma(e_t:t\in D)\subset\Borel(\mathcal C)$。反过来，对 $f,g\in\mathcal C$，连续性与 $D$ 稠密给
\[
  \|f-g\|_\infty=\sup_{t\in D}\|f(t)-g(t)\|.
\]
于是对 $r>0$，
\[
 B(f,r)
 =\bigcup_{\substack{q\in\Q\\0<q<r}}
   \bigcap_{t\in D}\{g\in\mathcal C:\|g(t)-f(t)\|<q\}
 \in\sigma(e_t:t\in D).
\]
可分度量空间的开集是可数个基开球之并，故得到
\eqref{eq:6-1-3-path-borel}。

对 $t\in D$，$e_t\circ\mathbf Y=Y_t$ 可测。式
\eqref{eq:6-1-3-path-borel} 因而给出 $\mathbf Y$ 的 Borel 可测性。最后，有限评价映射
$(e_{t_1},\ldots,e_{t_n})$ 连续，且
\[
 (e_{t_1},\ldots,e_{t_n})_*\Law(\mathbf Y)
 =\Law(Y_{t_1},\ldots,Y_{t_n}),
\]
所以路径空间上的分布恢复全部有限维分布。
\end{proof}

这个结论并未声称连续函数集合在 $\R^K$ 的柱 $\sigma$-代数中可测。事实上，当 $K$ 不可数时，任意
可数坐标集之外都可修改一个函数而不改变已读坐标，连续性因此不是由柱坐标中的某个可数子族决定的。
正确做法是先构造连续版本，再把它直接视作可分 Banach 空间 $C(K;\R^m)$ 中的随机元。
这也把过程分布放进 Prokhorov 紧性可以作用的 Polish 空间；但要证明一族过程分布是紧的，仍需另行给出
一致的路径模量与点值控制，单个过程有连续版本并不自动提供族的一致紧性。

\begin{remark}[两类增量模型的路径正则性]
以下两个例子从给定的增量分布出发，对比连续性判据对 Gaussian 型与 Poisson 型增量给出的
不同结论。
\end{remark}

\begin{example}[Brownian 型增量矩]\label{ex:6-1-3-2}
假定某过程满足 $B_t-B_s\sim N(0,|t-s|)$。若 $Z\sim N(0,1)$，则对每个 $p>0$，
$C_p=\E|Z|^p<\infty$，而尺度换元给
\[
  \E|B_t-B_s|^p=C_p|t-s|^{p/2}.
\]
给定 $\gamma<1/2$，可取足够大的 $p>2$ 使
$\gamma<1/2-1/p$。在定理~\ref{theorem:6-1-3-1} 中取
$\alpha=p$、$d=1$、$\beta=p/2-1$，便得到 $\gamma$-\Holder{} 版本。固定一个 $p>2$ 只能推出
$\gamma<1/2-1/p$；要覆盖每个 $\gamma<1/2$，矩阶必须随目标指数增大。定理也不包含端点
$\gamma=1/2$。
\end{example}

\begin{example}[Gaussian 核的连续性判据]\label{ex:6-1-3-gaussian-kernel-continuity}
设 $X$ 是例~\ref{ex:6-1-2-2} 构造的中心 Gaussian 过程，参数集为 $[0,1]^d$。若存在
$L,H>0$ 使
\[
 \E|X_t-X_s|^2
 =K(t,t)+K(s,s)-2K(s,t)
 \le L\|t-s\|^{2H},
\]
则对 $Z\sim N(0,1)$ 及任意 $p>0$，Gaussian 尺度关系给
\[
 \E|X_t-X_s|^p
 =\E|Z|^p\bigl(\E|X_t-X_s|^2\bigr)^{p/2}
 \le C_pL^{p/2}\|t-s\|^{pH}.
\]
给定 $\gamma<H$，取 $p$ 充分大，使 $p(H-\gamma)>d$。在
定理~\ref{theorem:6-1-3-1} 中置 $\alpha=p$、$\beta=pH-d>0$，便有
$\gamma<\beta/\alpha=H-d/p$，从而得到 $\gamma$-\Holder{} 版本。这里路径正则性来自协方差核对角线
附近的定量控制，不是“Gaussian”这个名称本身。
\end{example}

\begin{example}[Poisson 型路径不能连续]\label{ex:6-1-3-3}
假定 $\lambda>0$、$N_0=0$，且 $N_t-N_s$ 在 $0\le s<t$ 时服从参数
$\lambda(t-s)$ 的 Poisson 分布。由于增量
非负，
\[
 \E|N_t-N_s|=\lambda|t-s|,
\]
这只有参数维数 $d=1$ 的临界一次幂，没有定理所需的额外 $\beta>0$。更高阶矩在短时间内也仍以
$|t-s|$ 为主阶。

它不可能有连续版本。若 $Y$ 是连续版本，则对每个有理 $q\in[0,1]$，$Y_q$ 取整数值 a.s.；
对可数多个有理数取交并用连续性，得到几乎每条 $Y$ 路径都取值于离散集 $\Z$。从连通区间
$[0,1]$ 到 $\Z$ 的连续映射只能为常函数，所以 $Y_1=Y_0=0$ a.s.。但
$\Prob(N_1>0)=1-e^{-\lambda}>0$，与 $Y_1$ 是 $N_1$ 的版本矛盾。此类过程应放入 \cadlagPathSpace{}，
而不是连续路径空间。
\end{example}

\begin{figure}[htbp]
\centering
\begin{tikzpicture}[scale=0.92]
  \begin{scope}[xshift=0cm]
    \draw[book line,step=0.5] (0,0) grid (2,2);
    \draw[BookAccent,very thick] (0.5,1)--(1,1)--(1.5,1)--(1.5,1.5);
    \fill[BookAccent] (0.5,1) circle (1.5pt);
    \fill[BookAccent] (1,1) circle (1.5pt);
    \fill[BookAccent] (1.5,1) circle (1.5pt);
    \fill[BookAccent] (1.5,1.5) circle (1.5pt);
    \node[book label] at (1,-0.35) {相邻边尾界};
  \end{scope}
  \draw[book accent arrow] (2.35,1)--(3.25,1);
  \begin{scope}[xshift=3.6cm]
    \draw[book line,step=0.25] (0,0) grid (2,2);
    \draw[BookWarm,very thick]
      (0.25,0.5)--(0.5,0.5)--(0.5,0.75)--(0.75,0.75)--(0.75,1)--(1,1)--(1,1.25)--(1.25,1.25)--(1.25,1.5)--(1.5,1.5);
    \fill[BookWarm] (0.25,0.5) circle (1.7pt);
    \fill[BookWarm] (1.5,1.5) circle (1.7pt);
    \node[book label] at (1,-0.35) {多层链与几何级数};
  \end{scope}
  \draw[book accent arrow] (5.95,1)--(6.85,1);
  \node[book strong node,minimum width=2.5cm,minimum height=1.25cm] at (8.35,1)
    {稠密集上的 \Holder{} 控制\\$\Downarrow$\\完备性延拓};
\end{tikzpicture}
\caption{Kolmogorov--Chentsov 方法先同时控制每层有限多条相邻边，再以二进链求和，最后从
可数稠密集延拓到连续路径。}
\label{fig:6-1-3-dyadic-extension}
\end{figure}

\begin{remark}[版本、路径事件与分布]
版本改变保持所有有限维分布，却可能改变某些逐路径陈述的代表。若已经选定连续版本，则
命题~\ref{proposition:6-1-3-3} 保证该连续版本在不可分辨意义下唯一，
命题~\ref{proposition:6-1-3-4} 又把它的分布放在一个可分 Banach 路径空间上。以后讨论路径上确界、
停时、随机积分或 SDE 时，所用版本及其共同零集必须明确；“对每个 $t$ 都 a.s.”不能替代这项工作。
\end{remark}

\begin{exercise}
\begin{enumerate}
  \item 逐项验证例~\ref{ex:6-1-3-1} 中 $X,Y$ 互为版本但并非不可分辨，并证明 $X$ 相对于任一固定
        可数稠密集都不是可分版本。
  \item 重做引理~\ref{lemma:6-1-3-2} 的边数计算，明确指出参数维数 $d$ 在何处被矩指数
        $d+\beta$ 抵消。
  \item 从式~\eqref{eq:6-1-3-eventual-edge-control} 出发，完整构造父点链并推导
        \eqref{eq:6-1-3-dyadic-chain} 中的常数。
  \item 假定 $B_t-B_s\sim N(0,|t-s|)$，对给定 $\gamma<1/2$ 选一个具体偶整数 $p$，并写出
        Kolmogorov--Chentsov 定理中的 $\alpha,\beta$。
  \item 证明若 $T$ 可分，两个连续版本的路径空间分布相同；再说明为什么这个结论仍不提供它们在
        不同底层概率空间上的逐点耦合。
\end{enumerate}
\end{exercise}

至此，有限维相容性产生了过程分布，增量矩估计又从中选出可分析的路径。下一节加入随时间增长的
信息 $\sigma$-代数；鞅将把条件平均在这条信息流中的稳定性编码为等式。

\section{鞅：随信息演化的条件平均}\label{sec:06-02}

\subsection{滤过、鞅与条件平均的时间一致性}\label{subsec:6-2-1}
\index{滤过}\index{鞅}\index{次鞅}\index{可预测过程}

设 $X\in L^1$ 是一天结束时才能完全看见的赔付，而
$\mathcal F_0\subset\mathcal F_1\subset\cdots$ 是逐日收到的信息。时刻 $n$ 的自然估计是
\[
  M_n=\E[X\mid\mathcal F_n].
\]
如果到了时刻 $m<n$，先计算第 $n$ 日的估计再退回第 $m$ 日，塔式性质给出
\begin{equation}\label{eq:6-2-time-consistency-opening}
  \E[M_n\mid\mathcal F_m]
  =\E[\E[X\mid\mathcal F_n]\mid\mathcal F_m]
  =\E[X\mid\mathcal F_m]=M_m.
\end{equation}
这条时间一致性比“平均不变”强得多：它要求每一层已有信息上的条件平均都不变。鞅正是把
\eqref{eq:6-2-time-consistency-opening} 从一个终端变量的逐步估计中抽出来的结构。

\begin{definition}[滤过与适应性]\label{def:6-2-1-1}
设 $(\Omega,\mathcal F,\Prob)$ 是概率空间。离散时间滤过是满足
\[
  \mathcal F_0\subset\mathcal F_1\subset\cdots\subset\mathcal F
\]
的子 $\sigma$-代数族；连续时间滤过 $(\mathcal F_t)_{t\ge0}$ 满足
$\mathcal F_s\subset\mathcal F_t$ 对所有 $s\le t$ 成立。

过程 $X=(X_t)$ 若对每个 $t$ 都有 $X_t$ 对 $\mathcal F_t$ 可测，称它对该滤过
\emph{适应}（adapted）。离散时间过程 $(H_n)_{n\ge1}$ 若 $H_n$ 对
$\mathcal F_{n-1}$ 可测，称它可预测（predictable）。换言之，$H_n$ 必须在第 $n$ 个增量
出现以前确定。
\end{definition}

\begin{definition}[鞅、次鞅与上鞅]\label{def:6-2-1-2}
设 $(X_n)$ 是实值、可积且适应的过程。若
\[
  \E[X_{n+1}\mid\mathcal F_n]=X_n\quad\text{a.s.},
\]
则称 $X$ 为鞅；若左边不小于 $X_n$，称为次鞅；若左边不大于 $X_n$，称为上鞅。
复值鞅仍用等式定义；次鞅与上鞅只用于实值过程。

连续时间过程 $X=(X_t)_{t\ge0}$ 称为鞅，是指每个 $X_t\in L^1$、过程适应，并且
\[
  \E[X_t\mid\mathcal F_s]=X_s\quad\text{a.s.}\qquad(0\le s\le t).
\]
连续时间次鞅和上鞅相应把等号改为 $\ge$ 和 $\le$。这里必须直接要求任意 $s\le t$；
连续时间没有可以迭代的“相邻一步”。
\end{definition}

\begin{definition}[自然滤过]\label{def:6-2-1-3}
过程 $X$ 的自然滤过为
\[
  \mathcal F_n^X=\sigma(X_0,\ldots,X_n)
  \quad\text{或}\quad
  \mathcal F_t^X=\sigma(X_s:0\le s\le t).
\]
它记录截至当前时刻可由过程本身观察到的全部事件。给过程加入额外信息会得到更大的滤过；
鞅性质依赖所选滤过，因而不能把滤过从记号中无声删去。
\end{definition}

\begin{remark}[连续时间滤过的通常条件]\label{remark:6-2-1-usual}
连续时间中常要求滤过完备且右连续：前者表示它包含终端 $\sigma$-代数中所有零集的子集，后者表示
\[
  \mathcal F_t=\bigcap_{s>t}\mathcal F_s.
\]
这两项合称\emph{通常条件}（usual conditions）。它们使几乎处处修改与停时极限具有合适的可测性。离散时间中不需要
右连续性；Brownian 自然滤过的完备右连续化见 \S\ref{subsec:6-4-1}。
\end{remark}

“martingale”原是赌博术语；Doob 把它从公平游戏的比喻提升为条件期望理论的中心对象。独立增量只是一种
产生鞅的机制，而“未来增量在现有信息下条件均值为零”才是保留下来的本质，因此似然比、调和函数和
后来的随机积分都能进入同一框架。

相邻一步的条件可由塔式性质迭代，从而恢复任意两个时刻之间的时间一致性。

\begin{proposition}[时间一致性]\label{proposition:6-2-1-1}
设 $(M_n)$ 可积且适应。它是鞅，当且仅当对所有整数 $0\le m\le n$，
\[
  \E[M_n\mid\mathcal F_m]=M_m\quad\text{a.s.}
\]
次鞅情形也有相应结论：一步条件等价于
$\E[X_n\mid\mathcal F_m]\ge X_m$ 对所有 $m\le n$ 成立；把次鞅结论用于 $-X$，便得到
上鞅版本。
\end{proposition}

\begin{proof}
任意时刻条件取 $n=m+1$ 便给出相邻一步条件。反过来，设一步条件成立。对 $m<n$，塔式性质与归纳假设给出
\[
 \E[M_n\mid\mathcal F_m]
 =\E\bigl[\E[M_n\mid\mathcal F_{n-1}]\mid\mathcal F_m\bigr]
 =\E[M_{n-1}\mid\mathcal F_m]=M_m.
\]
归纳从 $n=m+1$ 开始，故鞅结论成立。若 $X$ 是次鞅，则条件期望的单调性给
\[
 \E[X_n\mid\mathcal F_m]
 =\E\bigl[\E[X_n\mid\mathcal F_{n-1}]\mid\mathcal F_m\bigr]
 \ge \E[X_{n-1}\mid\mathcal F_m]\ge X_m,
\]
同样完成归纳。上鞅情形对 $-X$ 使用次鞅结论即可。
\end{proof}

\begin{example}[独立中心和]\label{ex:6-2-1-1}
设 $\xi_1,\xi_2,\ldots$ 独立、可积且 $\E\xi_k=0$，令
$S_0=0$、$S_n=\sum_{k=1}^n\xi_k$ 及
$\mathcal F_n=\sigma(\xi_1,\ldots,\xi_n)$。由于 $\xi_{n+1}$ 独立于
$\mathcal F_n$，
\[
 \E[S_{n+1}\mid\mathcal F_n]
 =S_n+\E[\xi_{n+1}\mid\mathcal F_n]
 =S_n+\E\xi_{n+1}=S_n.
\]
所以中心独立增量的部分和是鞅。若均值为 $\mu_k$，则
$S_n-\sum_{k\le n}\mu_k$ 才是鞅；适应性本身不会消去漂移。
\end{example}

\begin{example}[闭鞅]\label{ex:6-2-1-2}
给定 $X\in L^1$ 和递增信息 $(\mathcal F_n)$，令
$M_n=\E[X\mid\mathcal F_n]$。每个 $M_n$ 可积并对 $\mathcal F_n$ 可测，且
\[
  \E[M_{n+1}\mid\mathcal F_n]
  =\E[\E[X\mid\mathcal F_{n+1}]\mid\mathcal F_n]
  =\E[X\mid\mathcal F_n]=M_n.
\]
因此 $(M_n)$ 是鞅。它在每一时刻都是同一终端变量 $X$ 的条件平均，故称为由 $X$ 闭合的鞅。
\end{example}

\begin{example}[似然比鞅]\label{ex:6-2-1-3}
设 $P,Q$ 是同一可测空间上的概率，并且对每个 $n$，
$Q|_{\mathcal F_n}\ll P|_{\mathcal F_n}$。取有限阶段的 Radon--Nikodym 密度
\[
  L_n=\frac{\dd Q|_{\mathcal F_n}}{\dd P|_{\mathcal F_n}}.
\]
$L_n$ 非负、$\mathcal F_n$-可测且 $\E_P L_n=1$。对任意
$A\in\mathcal F_n$，
\[
 \int_A L_{n+1}\dd P=Q(A)=\int_A L_n\dd P.
\]
条件期望的刻画因而给出
$\E_P[L_{n+1}\mid\mathcal F_n]=L_n$。所以 $(L_n)$ 是非负 $P$-鞅。
这里的逐阶段绝对连续是假设，不是“观察同一份数据”自动带来的结论。
\end{example}

凸函数把平均守恒变成单向漂移。

\begin{theorem}[条件 Jensen 产生次鞅]\label{theorem:6-2-1-2}
若 $(M_n)$ 是实值鞅，$\varphi:\R\to\R$ 凸，并且每个
$\varphi(M_n)$ 可积，则 $(\varphi(M_n))$ 是次鞅。特别地，只要相应幂次可积，
$(|M_n|^p)$ 对 $p\ge1$ 是次鞅，$(M_n^2)$ 也是次鞅。
\end{theorem}

\begin{proof}
条件 Jensen 不等式（定理~\ref{theorem:5-2-3-1}）及鞅等式给出
\[
 \E[\varphi(M_{n+1})\mid\mathcal F_n]
 \ge \varphi\bigl(\E[M_{n+1}\mid\mathcal F_n]\bigr)
 =\varphi(M_n).
\]
$\varphi(M_n)$ 对 $\mathcal F_n$ 可测，故满足次鞅定义。函数
$x\mapsto|x|^p$（$p\ge1$）和 $x\mapsto x^2$ 均凸，代入即得两个特例。
\end{proof}

鞅也可由逐步产生的“新信息”来描述。

\begin{proposition}[鞅差分分解]\label{proposition:6-2-1-3}
令 $M=(M_n)_{n\ge0}$ 为鞅，并置 $D_n=M_n-M_{n-1}$（$n\ge1$）。则
$D_n$ 对 $\mathcal F_n$ 可测且
\[
  \E[D_n\mid\mathcal F_{n-1}]=0.
\]
反之，若可积变量 $D_n$ 对 $\mathcal F_n$ 可测并满足上述条件，则对任意
$M_0\in L^1(\mathcal F_0)$，
$M_n=M_0+\sum_{k=1}^nD_k$ 是鞅。若 $M_n\in L^2$，则不同差分在 $L^2$ 中正交：
$\E[D_j\overline{D_k}]=0$ 对 $j\ne k$ 成立。
\end{proposition}

\begin{proof}
适应性给出 $D_n$ 的 $\mathcal F_n$-可测性，而
\[
 \E[D_n\mid\mathcal F_{n-1}]
 =\E[M_n\mid\mathcal F_{n-1}]-M_{n-1}=0.
\]
反过来，有限和可积且适应，并且
\[
 \E[M_n\mid\mathcal F_{n-1}]
 =M_{n-1}+\E[D_n\mid\mathcal F_{n-1}]=M_{n-1}.
\]
最后设 $j<k$。此时 $D_j$ 对 $\mathcal F_j\subset\mathcal F_{k-1}$ 可测，条件期望取因子给出
\[
 \E[D_j\overline{D_k}]
 =\E\bigl[D_j\overline{\E[D_k\mid\mathcal F_{k-1}]}\bigr]=0.
\]
平方可积保证乘积可积，因此计算合法。
\end{proof}

“下注策略”只能使用下注前的信息；这正是可预测性的数学内容。

\begin{proposition}[可预测变换]\label{proposition:6-2-1-4}
设 $M$ 是鞅，$(H_k)_{k\ge1}$ 是有界可预测过程。定义
\[
 (H\mathbin{\boldsymbol\cdot}M)_0=0,
 \qquad
 (H\mathbin{\boldsymbol\cdot}M)_n
 =\sum_{k=1}^nH_k(M_k-M_{k-1}).
\]
则 $H\mathbin{\boldsymbol\cdot}M$ 是鞅。若 $X$ 是次鞅且 $H_k\ge0$，同一变换是次鞅。
\end{proposition}

\begin{proof}
令 $D_k=M_k-M_{k-1}$。有界性保证 $H_kD_k\in L^1$，而
$H_k$ 对 $\mathcal F_{k-1}$ 可测，故
\[
 \E[H_kD_k\mid\mathcal F_{k-1}]
 =H_k\E[D_k\mid\mathcal F_{k-1}]=0.
\]
命题~\ref{proposition:6-2-1-3} 的反向即说明变换为鞅。若 $X$ 是次鞅，
$\E[X_k-X_{k-1}\mid\mathcal F_{k-1}]\ge0$；再由 $H_k\ge0$ 得每个变换差分的条件
期望非负，所以所得过程是次鞅。
\end{proof}

可预测性是结论的必要结构条件。令 $M$ 为简单对称随机游走，若先看见
$D_k=M_k-M_{k-1}\in\{-1,1\}$ 再取 $H_k=\sgn(D_k)$，则
$H_kD_k=|D_k|=1$，累计收益恒为 $n$。这里 $H_k$ 对 $\mathcal F_k$ 可测，却不对
$\mathcal F_{k-1}$ 可测。这个系数使用了下注时尚不可得的结果，因而不是允许的可预测策略。

次鞅比鞅多出的系统漂移可以被唯一分离出来。

\begin{proposition}[离散 Doob 分解]\label{proposition:6-2-1-5}
设 $X=(X_n)_{n\ge0}$ 是可积次鞅。存在分解
\[
  X_n=M_n+A_n,
\]
其中 $M$ 是鞅，$A_0=0$，并且 $A$ 可积、可预测且逐点递增；该分解在不可分辨意义下唯一。具体地，
\begin{equation}\label{eq:6-2-doob-decomposition}
 A_n=\sum_{k=1}^n
 \E[X_k-X_{k-1}\mid\mathcal F_{k-1}],
 \qquad M_n=X_n-A_n.
\end{equation}
\end{proposition}

\begin{proof}
记
$\widetilde a_k=\E[X_k-X_{k-1}\mid\mathcal F_{k-1}]$。次鞅性给出
$\widetilde a_k\ge0$ a.s.。逐项置 $a_k=(\widetilde a_k)^+$；它仍是同一个条件期望的
$\mathcal F_{k-1}$-可测可积版本，并且对每个样本点都有 $a_k\ge0$。于是 $A_0=0$、
$A_n=\sum_{k\le n}a_k$ 可积、递增，并且对 $n\ge1$，$A_n$ 对
$\mathcal F_{n-1}$ 可测。再计算
\[
 \E[M_n-M_{n-1}\mid\mathcal F_{n-1}]
 =\E[X_n-X_{n-1}\mid\mathcal F_{n-1}]-a_n=0,
\]
故 $M$ 是鞅。

若还有 $X=M'+A'$，其中 $M'$ 为鞅且 $A'$ 满足所列条件，则
\[
 A'_n-A'_{n-1}
 =\E[X_n-X_{n-1}\mid\mathcal F_{n-1}],
\]
因为左边已对 $\mathcal F_{n-1}$ 可测，而 $M'$ 的增量条件均值为零。由 $A'_0=0$
逐项求和得到 $A'_n=A_n$ a.s.，继而 $M'_n=M_n$ a.s.。时间指标可数，故可把这些零集取并，
得到两组过程在同一个概率一事件上对所有 $n$ 同时相等。
\end{proof}

\begin{example}[平方随机游走的补偿]
设例~\ref{ex:6-2-1-1} 中的增量还满足
$\E\xi_k^2=\sigma^2<\infty$。由于 $S_{k-1}$ 对 $\mathcal F_{k-1}$ 可测，且
$\xi_k$ 独立于过去，
\begin{align*}
 \E[S_k^2-S_{k-1}^2\mid\mathcal F_{k-1}]
 &=\E[2S_{k-1}\xi_k+\xi_k^2\mid\mathcal F_{k-1}]\\
 &=\sigma^2.
\end{align*}
因此 $S_n^2$ 是次鞅，其 Doob 分解为
\[
 S_n^2=\underbrace{(S_n^2-n\sigma^2)}_{\text{鞅}}
       +\underbrace{n\sigma^2}_{\text{可预测补偿}}.
\]
同一计算将在 Poisson 过程中变成补偿计数，在 Brownian 理论中变成二次变差。
\end{example}

\begin{figure}[htbp]
\centering
\begin{tikzpicture}[node distance=12mm and 17mm]
  \node[book strong node] (f0) {$\mathcal F_0$};
  \node[book strong node,right=of f0] (f1) {$\mathcal F_1$};
  \node[book strong node,right=of f1] (f2) {$\mathcal F_2$};
  \node[book strong node,right=of f2] (fn) {$\cdots\subset\mathcal F_n$};
  \node[book warm node,above=15mm of f1] (x) {终端变量 $X$};
  \draw[book accent arrow] (f0)--node[below,book label] {信息增加}(f1);
  \draw[book accent arrow] (f1)--(f2);
  \draw[book accent arrow] (f2)--(fn);
  \draw[book dashed arrow] (x)--node[left,book label] {$\E[\,\cdot\mid\mathcal F_0]$}(f0);
  \draw[book dashed arrow] (x)--node[right,book label] {$M_1$}(f1);
  \draw[book dashed arrow] (x)--node[right,book label] {$M_2$}(f2);
  \draw[decorate,decoration={brace,mirror,amplitude=4pt},BookWarm]
    ($(f1.south west)+(0,-5mm)$)--($(f2.south east)+(0,-5mm)$)
    node[midway,below=5pt,book label] {差分只携带相邻层的新信息};
\end{tikzpicture}
\caption{终端变量向递增信息层作条件平均，形成闭鞅；塔式性质保证投影的时间一致性。}
\label{fig:6-2-1-filtration}
\end{figure}

适应只说明“目前能够读到 $X_n$”，并不说明下一步条件均值为零；可预测只说明策略在增量出现前
确定，也不保证其有界或可积。鞅定义中的可积性、适应性和条件平均三项缺一不可。连续时间还要另行
处理滤过的完备化与右连续化；图~\ref{fig:6-2-1-filtration} 只表达离散层级，不能替代这些条件。

\begin{exercise}
设 $X\in L^1$，$M_n=\E[X\mid\mathcal F_n]$。直接用条件期望的积分刻画而非塔式性质的结论，
验证 $M$ 的鞅等式。
\end{exercise}

\begin{exercise}
设 $M\in L^2$ 是实值离散鞅。证明
$\E M_n^2=\E M_0^2+\sum_{k=1}^n\E(M_k-M_{k-1})^2$，并指出每一步使用了
命题~\ref{proposition:6-2-1-3} 的哪一部分。
\end{exercise}

\begin{exercise}
在简单对称随机游走中取 $X_n=S_n+n$。验证 $X$ 适应且可积，但不是鞅；判断它是次鞅还是上鞅。
\end{exercise}

最大不等式将把终端时刻的矩控制转成整段路径的控制。为此必须允许“第一次越界”这样的随机时刻
进入证明。

\subsection{Doob 最大不等式、上穿与鞅收敛}\label{subsec:6-2-2}
\index{Doob 最大不等式}\index{上穿不等式}\index{鞅收敛定理}

逐个时刻估计 $\Prob(|M_k|\ge\lambda)$，再对 $k\le n$ 作并集估计，通常会多付出一个
随 $n$ 增长的因子。鞅结构允许在路径\emph{第一次}越界的时刻使用条件平均，把所有越界路径只计算
一次。另一方面，最大值有界仍不直接排除序列在两个水平之间永远振荡；上穿次数正是记录这种振荡的
离散计数器。

Doob 的最大不等式与收敛定理遵循同一个方法：先建立对所有有限时刻统一的估计，再让时间上界趋于
无穷。上穿论证则把一个可预测交易策略转化为路径不等式
\eqref{eq:6-2-upcross-pathwise}，从而将振荡次数与期望估计联系起来。

\begin{definition}[最大过程]\label{def:6-2-2-1}
对离散过程 $M$，定义
\[
  M_n^*=\max_{0\le k\le n}|M_k|.
\]
连续时间相应的表达式是 $\sup_{0\le s\le t}|M_s|$；在使用它以前还需保证路径可测性或右连续性，
本小节只讨论离散时间。
\end{definition}

\begin{definition}[上穿次数]\label{def:6-2-2-2}
给定实数列 $x_0,\ldots,x_n$ 与 $a<b$，区间 $[a,b]$ 的上穿次数
$U_n[a,b]$ 是满足
\[
  0\le s_1<t_1<s_2<t_2<\cdots<s_m<t_m\le n,
  \qquad x_{s_j}\le a,\quad x_{t_j}\ge b
\]
的最大整数 $m$。对随机过程 $X$，把 $x_k$ 换成 $X_k(\omega)$，便得到随机变量
$U_n[a,b]$。
\end{definition}

\begin{definition}[闭鞅与一致可积鞅]\label{def:6-2-2-3}
若存在 $Y\in L^1$ 使
\[
  M_n=\E[Y\mid\mathcal F_n]\quad\text{a.s.}
\]
对每个 $n$ 成立，则称 $M$ 为闭鞅。若随机变量族
$\{M_n:n\ge0\}$ 按定义~\ref{def:5-3-2-1} 一致可积，称 $M$ 为一致可积鞅。
\end{definition}

\begin{theorem}[Doob $L^1$ 最大不等式]\label{theorem:6-2-2-1}
设 $(X_k)_{0\le k\le n}$ 是非负次鞅。对每个 $\lambda>0$，
\begin{equation}\label{eq:6-2-doob-weak-one}
 \lambda\Prob\!\left(\max_{0\le k\le n}X_k\ge\lambda\right)
 \le \E\!\left[X_n\1_{\{\max_{k\le n}X_k\ge\lambda\}}\right]
 \le \E X_n.
\end{equation}
\end{theorem}

\begin{proof}
令
\[
 A_k=\{X_0<\lambda,\ldots,X_{k-1}<\lambda,\ X_k\ge\lambda\},
 \qquad 0\le k\le n,
\]
其中 $A_0=\{X_0\ge\lambda\}$。这些事件两两不交，$A_k\in\mathcal F_k$，且其并为
$A=\{\max_{k\le n}X_k\ge\lambda\}$。由命题~\ref{proposition:6-2-1-1} 的次鞅版本，
\begin{align*}
 \lambda\Prob(A_k)
 &\le \E[X_k\1_{A_k}]\\
 &\le \E\!\left[\E[X_n\mid\mathcal F_k]\1_{A_k}\right]
 =\E[X_n\1_{A_k}].
\end{align*}
对 $k$ 求和便得到第一个不等式。第二个不等式来自 $X_n\ge0$ 和
$\1_A\le1$。
\end{proof}

\begin{figure}[htbp]
\centering
\begin{tikzpicture}[node distance=7mm and 8mm]
  \node[book node] (a0) {$A_0$\\初始即越界};
  \node[book node,right=of a0] (a1) {$A_1$\\首次在 $1$ 越界};
  \node[right=of a1] (dots) {$\cdots$};
  \node[book node,right=of dots] (an) {$A_n$\\首次在 $n$ 越界};
  \node[book strong node,below=10mm of dots] (u)
    {$\displaystyle\bigsqcup_{k=0}^n A_k
      =\left\{\max_{j\le n}X_j\ge\lambda\right\}$};
  \draw[book arrow] (a0)--(u);
  \draw[book arrow] (a1)--(u);
  \draw[book arrow] (an)--(u);
\end{tikzpicture}
\caption{首次越界时刻把最大事件分成互不相交的 $\mathcal F_k$-可测区域；每条越界路径只被计算一次。}
\label{fig:6-2-2-first-crossing}
\end{figure}

弱型 $L^1$ 估计控制每一个尾概率；当 $p>1$ 时，层蛋糕公式可以把全部尾概率重新积成强
$L^p$ 估计。

\begin{theorem}[Doob $L^p$ 不等式]\label{theorem:6-2-2-2}
令 $p>1$，$q=p/(p-1)$。若 $M=(M_k)_{0\le k\le n}$ 是鞅且
$M_n\in L^p$，则
\[
  \|M_n^*\|_p\le q\|M_n\|_p.
\]
若 $M=(M_k)_{0\le k\le n}$ 是非负次鞅且 $M_n\in L^p$，同一结论成立，此时
$M_n^*=\max_{k\le n}M_k$。
\end{theorem}

\begin{proof}
在鞅情形，时间一致性与条件模长不等式给
\[
 |M_k|^p
 \le\bigl(\E[|M_n|\mid\mathcal F_k]\bigr)^p
 \le\E[|M_n|^p\mid\mathcal F_k],
\]
故 $M_k\in L^p$。非负次鞅情形由
$0\le M_k\le\E[M_n\mid\mathcal F_k]$ 和同一标量条件 Jensen 估计得到相同结论。
鞅情形令 $X_k=|M_k|$；条件模长不等式给
\[
 \E[X_{k+1}\mid\mathcal F_k]
 \ge \left|\E[M_{k+1}\mid\mathcal F_k]\right|=|M_k|,
\]
所以 $X$ 是非负次鞅。非负次鞅情形直接令 $X=M$。统一记
$X_n^*=\max_{k\le n}X_k$。对 $R>0$，尾积分公式与
定理~\ref{theorem:6-2-2-1} 给出
\begin{align*}
 \E[(X_n^*\wedge R)^p]
 &=p\int_0^R\lambda^{p-1}\Prob(X_n^*\ge\lambda)\dd\lambda\\
 &\le p\int_0^R\lambda^{p-2}
       \E[X_n\1_{\{X_n^*\ge\lambda\}}]\dd\lambda\\
 &=\frac{p}{p-1}\E\!\left[X_n(X_n^*\wedge R)^{p-1}\right].
\end{align*}
最后一式由 Tonelli 定理逐点积分得到。\Holder{} 不等式于是给出
\[
 \E[(X_n^*\wedge R)^p]
 \le q\|X_n\|_p
       \bigl\|(X_n^*\wedge R)^{p-1}\bigr\|_q
 =q\|X_n\|_p
       \|X_n^*\wedge R\|_p^{p-1}.
\]
若左边为零结论成立；否则约去其 $(p-1)/p$ 次幂，得到
$\|X_n^*\wedge R\|_p\le q\|X_n\|_p$。令 $R\uparrow\infty$，由单调收敛定理
得到结论。
\end{proof}

\begin{example}[随机游走最大值]\label{ex:6-2-2-1}
对方差为一的中心随机游走 $S_n=\sum_{k=1}^n\xi_k$，取 $p=2$ 得
\[
 \Prob\!\left(\max_{k\le n}|S_k|\ge\lambda\right)
 \le \frac{\E(S_n^*)^2}{\lambda^2}
 \le \frac{4\E S_n^2}{\lambda^2}
 =\frac{4n}{\lambda^2}.
\]
逐时刻使用 Chebyshev 不等式再作并集估计只给出
$\lambda^{-2}\sum_{k=1}^n\E S_k^2=n(n+1)/(2\lambda^2)$。最大不等式利用了路径的共同历史，
因此把二次量级降到一次量级。
\end{example}

接下来把振荡翻译成一项可预测交易。策略在 $X\le a$ 时买入一单位，随后一直持有，直到
$X\ge b$ 时卖出；完成一次交易便完成一次上穿。

\begin{theorem}[Doob 上穿不等式]\label{theorem:6-2-2-3}
若 $(X_k)_{0\le k\le n}$ 是实值次鞅，$a<b$，则
\begin{equation}\label{eq:6-2-upcrossing}
 (b-a)\E U_n[a,b]
 \le \E[(X_n-a)^-]+\E[X_n-X_0],
\end{equation}
其中 $x^- =\max\{-x,0\}$。
\end{theorem}

\begin{proof}
逐路径递归定义 $H_k\in\{0,1\}$。起初不持有；若在时刻 $k-1$ 不持有且
$X_{k-1}\le a$，则令 $H_k=1$ 并开始持有；持有期间令 $H_k=1$，直到某个时刻
$k-1$ 已满足 $X_{k-1}\ge b$，从下一增量起令 $H_k=0$；随后重复。这个规则只使用
$X_0,\ldots,X_{k-1}$，所以 $H_k$ 对 $\mathcal F_{k-1}$ 可测。

记 $G_n=\sum_{k=1}^nH_k(X_k-X_{k-1})$。每个已经卖出的单位至少赚 $b-a$。若时刻
$n$ 尚有一笔未完成持仓，买入价不高于 $a$，其账面收益至少为 $X_n-a$。因此逐路径有
\begin{equation}\label{eq:6-2-upcross-pathwise}
  (b-a)U_n[a,b]\le G_n+(X_n-a)^-.
\end{equation}
上述贪心规则完成的交易数确实等于定义
\ref{def:6-2-2-2} 中的最大上穿数。设它依次选出
\[
 s_1=\min\{k\le n:X_k\le a\},\qquad
 t_1=\min\{k>s_1:X_k\ge b\},
\]
并在 $t_j<n$ 时递归地取
\[
 s_{j+1}=\min\{k>t_j:X_k\le a\},\qquad
 t_{j+1}=\min\{k>s_{j+1}:X_k\ge b\},
\]
空集的最小值视为 $+\infty$。若这个规则完成了 $q$ 次，而
$s'_1<t'_1<\cdots<s'_m<t'_m$ 是任意一组上穿，则归纳得
\[
 s_j\le s'_j,\qquad t_j\le t'_j\qquad(1\le j\le m).
\]
其中 $j=1$ 来自最小值的定义；若对 $j$ 成立，则
$t_j\le t'_j<s'_{j+1}$，所以 $s_{j+1}\le s'_{j+1}$，进而
$t_{j+1}\le t'_{j+1}$。因此 $m\le q$；反过来，贪心规则自身给出长度为 $q$ 的上穿族，故
$q=U_n[a,b]$。于是前述收益分解正好给出
式~\eqref{eq:6-2-upcross-pathwise}。
由于 $H$ 非负且可预测，次鞅差分给 $\E G_n\ge0$，但这不是
\eqref{eq:6-2-upcrossing} 所需的方向。注意互补策略 $1-H$ 也非负且可预测，并且
\[
  X_n-X_0=G_n+
  \sum_{k=1}^n(1-H_k)(X_k-X_{k-1}).
\]
第二项期望非负，故 $\E G_n\le\E(X_n-X_0)$。对
\eqref{eq:6-2-upcross-pathwise} 取期望即得结论。
\end{proof}

末项 $\E[X_n-X_0]$ 不能随意换成只含 $X_0$ 正部的表达式。例如取确定性递增次鞅
$X_k=k$ 及 $0<a<b<n$，上穿次数为一；若丢掉终端漂移，右端可能变成零。公式的方向与
定义~\ref{def:6-2-2-2} 所数的“低到高”上穿严格配套。

\begin{theorem}[鞅／次鞅收敛定理]\label{theorem:6-2-2-4}
若实值次鞅 $(X_n)$ 满足
\[
  \sup_{n\ge0}\E X_n^+<\infty,
\]
则存在有限的 $X_\infty\in L^1$，使 $X_n\to X_\infty$ a.s.；这里不声称
$L^1$ 收敛。

若鞅 $(M_n)$ 对某个 $p>1$ 满足 $\sup_n\|M_n\|_p<\infty$，则存在
$M_\infty\in L^p$，使 $M_n\to M_\infty$ a.s. 且在 $L^p$ 中收敛。
\end{theorem}

\begin{proof}
令 $C=\sup_n\E X_n^+<\infty$。对固定有理数 $a<b$，
定理~\ref{theorem:6-2-2-3} 的右端满足
\begin{align*}
 \E[(X_n-a)^-]+\E[X_n-X_0]
 &=\E[\max\{X_n,a\}]-\E X_0\\
 &\le C+a^+-\E X_0.
\end{align*}
因此 $\sup_n\E U_n[a,b]<\infty$。由于 $U_n[a,b]\uparrow U_\infty[a,b]$，
单调收敛定理给 $\E U_\infty[a,b]<\infty$，从而
$U_\infty[a,b]<\infty$ a.s. 对每一对有理数 $a<b$ 成立。这样的有理数对只有可数多个，
故可取一个共同的概率一事件。

若在该事件上 $\liminf_nX_n<\limsup_nX_n$，可在两者之间选有理数 $a<b$；数列随后会完成
无穷多次从 $a$ 以下到 $b$ 以上的上穿，与 $U_\infty[a,b]<\infty$ 矛盾。因此
$X_n$ 在扩展实数中有极限。又因次鞅的均值非降，
\[
 \E X_n^- =\E X_n^+-\E X_n\le C-\E X_0,
\]
所以 $\sup_n\E|X_n|\le2C-\E X_0<\infty$。Fatou 引理给
\[
 \E\!\left[\liminf_n|X_n|\right]
 \le\liminf_n\E|X_n|<\infty.
\]
这排除了极限为 $\pm\infty$ 的正概率事件，并且给出
$X_\infty\in L^1$。

现在设 $M$ 在 $L^p$ 中有界。若 $M$ 为实值，它也在 $L^1$ 中有界，故第一部分给出 a.s. 极限；
若 $M$ 为复值，则 $\Re M$ 与 $\Im M$ 都是 $L^1$ 有界的实鞅，分别应用第一部分后仍得到
$M_n\to M_\infty$ a.s.。对每个 $N$，定理~\ref{theorem:6-2-2-2} 给
\[
 \E\!\left[\max_{k\le N}|M_k|^p\right]
 \le q^p\sup_n\E|M_n|^p,
 \qquad q=\frac{p}{p-1}.
\]
令 $N\to\infty$，得到 $M^*=\sup_k|M_k|\in L^p$。于是
$|M_n-M_\infty|^p\le(2M^*)^p$，且左边 a.s. 趋于零；支配收敛定理给
$\|M_n-M_\infty\|_p\to0$，并同时说明 $M_\infty\in L^p$。
\end{proof}

上一定理的第一部分没有给 $L^1$ 收敛；缺失的条件正是一致可积。

\begin{proposition}[UI、$L^1$ 收敛与闭合]\label{proposition:6-2-2-5}
对离散时间鞅 $(M_n)$，下列三项等价：
\begin{enumerate}
 \item 随机变量族 $\{M_n:n\ge0\}$ 一致可积；
 \item 存在 $M_\infty\in L^1$，使 $M_n\to M_\infty$ 于 $L^1$；
 \item 存在 $Y\in L^1$，使 $M_n=\E[Y\mid\mathcal F_n]$ 对所有 $n$ 成立。
\end{enumerate}
在第 (2) 项成立时，可在第 (3) 项取 $Y=M_\infty$。
\end{proposition}

\begin{proof}
设 (1) 成立。一致可积蕴含 $\sup_n\E|M_n|<\infty$。实值情形由
定理~\ref{theorem:6-2-2-4} 得到 a.s. 收敛；复值情形分别把该定理用于实部和虚部。
因此 $M_n$ a.s. 收敛到某个 $M_\infty\in L^1$。
Vitali 定理~\ref{theorem:5-3-2-1} 把 a.s. 收敛与 UI 合并，得到
$M_n\to M_\infty$ 于 $L^1$，所以 (2) 成立。

设 (2) 成立。固定 $n$。对每个 $m\ge n$，时间一致性给
$M_n=\E[M_m\mid\mathcal F_n]$。条件期望的 $L^1$ 收缩性给
\[
 \bigl\|\E[M_m-M_\infty\mid\mathcal F_n]\bigr\|_1
 \le\|M_m-M_\infty\|_1\longrightarrow0.
\]
因此 $M_n=\E[M_\infty\mid\mathcal F_n]$ a.s.，得到 (3)。

最后设 (3) 成立，并令 $A_{n,K}=\{|M_n|>K\}\in\mathcal F_n$。由条件模长
不等式，
\[
 \E[|M_n|\1_{A_{n,K}}]
 \le\E[|Y|\1_{A_{n,K}}].
\]
而 $\|M_n\|_1\le\|Y\|_1$，故
$\Prob(A_{n,K})\le\|Y\|_1/K$。对任意 $R>0$，
\[
 \E[|Y|\1_{A_{n,K}}]
 \le \E[|Y|\1_{\{|Y|>R\}}]
      +R\frac{\|Y\|_1}{K}.
\]
先取 $R$ 使第一项小，再令 $K\to\infty$，所得界与 $n$ 无关，故
$\{M_n\}$ UI。
\end{proof}

\begin{example}[非 UI 非负鞅]\label{ex:6-2-2-2}
考虑 Galton--Watson 过程：$Z_0=1$，每个个体独立地产生 $0$ 或 $2$ 个后代，概率各为
$1/2$。给定第 $n$ 代的信息 $\mathcal F_n$，每个现存个体的平均后代数为一，所以
\[
  \E[Z_{n+1}\mid\mathcal F_n]=Z_n,
  \qquad \E Z_n=1.
\]
令 $q_n=\Prob(Z_n=0)$。从一个祖先出发，下一代要么立即灭绝，要么由两个独立子过程组成，故
\[
  q_0=0,\qquad q_{n+1}=\frac{1+q_n^2}{2}.
\]
由 $0\le q_n\le1$ 及
\[
 q_{n+1}-q_n=\frac{(1-q_n)^2}{2}\ge0,
\]
归纳得序列 $(q_n)$ 递增且不超过一；其极限 $q$ 满足
$q=(1+q^2)/2$，即 $(q-1)^2=0$，所以 $q_n\uparrow1$。灭绝后过程保持为零，因而
\[
 \Prob(\text{最终灭绝})=\Prob\Bigl(\bigcup_n\{Z_n=0\}\Bigr)=\lim_nq_n=1,
\]
从而 $Z_n\to0$ a.s.。但 $\E Z_n=1$ 对所有 $n$ 成立，所以不可能在 $L^1$ 收敛到零；
命题~\ref{proposition:6-2-2-5} 说明 $(Z_n)$ 不是 UI。稀有而庞大的幸存族群携带了没有出现在
点态极限中的全部期望。
\end{example}

\begin{example}[条件期望逼近]\label{ex:6-2-2-3}
设 $X\in L^p$、$p>1$，并令 $M_n=\E[X\mid\mathcal F_n]$。条件 Jensen 给
$\|M_n\|_p\le\|X\|_p$，故定理~\ref{theorem:6-2-2-4} 给出 a.s. 及 $L^p$ 极限
$M_\infty$。若
\[
  \mathcal F_\infty=\sigma\Bigl(\bigcup_{n\ge0}\mathcal F_n\Bigr),
\]
则 $M_\infty$ 对 $\mathcal F_\infty$ 可测。对 $A\in\bigcup_n\mathcal F_n$，从某一时刻起
$\E[M_k\1_A]=\E[X\1_A]$；令 $k\to\infty$ 得同一等式对 $M_\infty$ 成立。
包含这些 $A$ 且保持该等式的事件类是 Dynkin 系统，而
$\bigcup_n\mathcal F_n$ 是代数，故 $\pi$--$\lambda$ 定理推出
\[
  M_\infty=\E[X\mid\mathcal F_\infty].
\]
若 $X$ 本身对 $\mathcal F_\infty$ 可测，极限才等于 $X$；“信息递增”不等于“最终揭示一切”。
\end{example}

\begin{figure}[htbp]
\centering
\begin{tikzpicture}[x=0.72cm,y=0.72cm]
  \draw[->,BookInk!75] (0,0)--(12.2,0) node[right] {$k$};
  \draw[->,BookInk!75] (0,-1.2)--(0,5.2) node[above] {$X_k$};
  \draw[BookAccent,semithick] (0,1.2)--(12,1.2) node[right] {$a$};
  \draw[BookAccent,semithick] (0,3.6)--(12,3.6) node[right] {$b$};
  \draw[BookInk,thick]
    plot coordinates {(0,2.4) (1,0.8) (2,1.0) (3,2.0) (4,3.9)
      (5,3.0) (6,0.7) (7,1.1) (8,2.5) (9,4.1) (10,2.7) (11,0.9) (12,2.0)};
  \fill[BookWarm] (1,0.8) circle (2pt) node[below left] {$s_1$};
  \fill[BookWarm] (4,3.9) circle (2pt) node[above] {$t_1$};
  \fill[BookWarm] (6,0.7) circle (2pt) node[below] {$s_2$};
  \fill[BookWarm] (9,4.1) circle (2pt) node[above] {$t_2$};
  \draw[decorate,decoration={brace,amplitude=4pt},BookWarm]
    (1.1,0.25)--(3.9,0.25) node[midway,below=5pt] {持有};
  \draw[decorate,decoration={brace,amplitude=4pt},BookWarm]
    (6.1,0.25)--(8.9,0.25) node[midway,below=5pt] {持有};
\end{tikzpicture}
\caption{可预测策略只在看到 $X\le a$ 后买入，在首次达到 $X\ge b$ 后卖出；两笔完成交易对应两次上穿。}
\label{fig:6-2-2-upcrossings}
\end{figure}

上穿次数有限排除的是实序列的无限振荡；极限为 $\pm\infty$ 仍需用 $L^1$ 界排除，这正是
定理~\ref{theorem:6-2-2-4} 证明的最后一步。另一方面，$L^1$ 有界鞅未必在 $L^1$ 中收敛，
例~\ref{ex:6-2-2-2} 已给出完全可计算的反例。$p=1$ 处最大不等式只剩
\eqref{eq:6-2-doob-weak-one} 的弱型估计；强型常数 $p/(p-1)$ 在 $p\downarrow1$ 时发散并非偶然。

\begin{exercise}
从定理~\ref{theorem:6-2-2-1} 出发，逐行重做定理~\ref{theorem:6-2-2-2} 的尾积分证明，
并解释为何 $p=1$ 时出现不可积的因子 $\lambda^{-1}$。
\end{exercise}

\begin{exercise}
对一个确定性实数列写出生成 $H_k\in\{0,1\}$ 的递归规则，证明
\eqref{eq:6-2-upcross-pathwise}，并验证 $H_k$ 只依赖 $x_0,\ldots,x_{k-1}$。
\end{exercise}

\begin{exercise}
证明若 $\sup_n\|M_n\|_p<\infty$、$p>1$，则 $\{M_n:n\ge0\}$ 一致可积。
比较这一短证与定理~\ref{theorem:6-2-2-4} 中得到 $L^p$ 收敛的证明。
\end{exercise}

有了路径最大值与极限控制，下一步可以在由路径自身决定的时刻读取鞅。但“看见赢了就停”只有在
停下的决定不借用未来，并且尾部质量不会逃到无穷时才安全。

\subsection{停时、可选抽样与一致可积性}\label{subsec:6-2-3}
\index{停时}\index{可选抽样}\index{停止过程}

“第一次达到十元”可由目前为止的记录判断；“最后一次达到十元”通常要等未来全部结束以后才知道。
这个差别决定了随机时刻能否与条件期望相容。

\begin{definition}[停时]\label{def:6-2-3-1}
取值于 $\{0,1,2,\ldots\}\cup\{\infty\}$ 的随机变量 $\tau$ 若满足
\[
  \{\tau\le n\}\in\mathcal F_n\qquad(n\ge0),
\]
则称 $\tau$ 为关于滤过 $(\mathcal F_n)$ 的停时。等价地，
$\{\tau=n\}\in\mathcal F_n$ 对每个有限 $n$ 成立。
\end{definition}

等价性来自
\[
 \{\tau=n\}=\{\tau\le n\}\setminus\{\tau\le n-1\},
 \qquad
 \{\tau\le n\}=\bigcup_{k=0}^n\{\tau=k\},
\]
其中 $n=0$ 时把第二个被减集合视为空集。特别地，
$\{\tau\ge k\}=\{\tau\le k-1\}^{\mathrm c}\in\mathcal F_{k-1}$：在第 $k$ 个增量出现前，
我们知道过程是否还没有停。

\begin{definition}[停止过程与停止信息]\label{def:6-2-3-2}
过程 $X$ 在停时 $\tau$ 处的停止过程为
\[
  X_n^\tau=X_{n\wedge\tau}.
\]
停下时可知的事件构成
\[
 \mathcal F_\tau
 =\{A\in\mathcal F:A\cap\{\tau\le n\}\in\mathcal F_n
                  \text{ 对每个 }n\}.
\]
这通常远大于 $\sigma(\tau)$：它不仅知道何时停，也保留停下以前观察到的路径信息。
\end{definition}

\begin{lemma}[停止信息的基本性质]\label{lemma:6-2-3-stopped-sigma}
$\mathcal F_\tau$ 是 $\sigma$-代数。若 $\sigma\le\tau$ 是两个停时，则
$\mathcal F_\sigma\subset\mathcal F_\tau$。若 $X$ 适应且 $\tau<\infty$，则
$X_\tau$ 对 $\mathcal F_\tau$ 可测；若 $\tau\le N$，则 $X_\tau$ 可积，只要
$X_0,\ldots,X_N$ 均可积。
\end{lemma}

\begin{proof}
空集属于 $\mathcal F_\tau$。若 $A\in\mathcal F_\tau$，则
\[
 A^{\mathrm c}\cap\{\tau\le n\}
 =\{\tau\le n\}\setminus(A\cap\{\tau\le n\})\in\mathcal F_n.
\]
可数并的条件由交与并的分配律逐项得到，故 $\mathcal F_\tau$ 是 $\sigma$-代数。

若 $A\in\mathcal F_\sigma$ 且 $\sigma\le\tau$，则
\[
 A\cap\{\tau\le n\}
 =(A\cap\{\sigma\le n\})\cap\{\tau\le n\}\in\mathcal F_n,
\]
所以 $A\in\mathcal F_\tau$。最后，对 Borel 集 $B$，
\[
 \{X_\tau\in B\}\cap\{\tau\le n\}
 =\bigcup_{k=0}^n\bigl(\{\tau=k\}\cap\{X_k\in B\}\bigr)\in\mathcal F_n,
\]
这证明停止值可测。若 $\tau\le N$，则
$|X_\tau|\le\sum_{k=0}^N|X_k|$，故可积。
\end{proof}

\begin{example}[首次命中]\label{ex:6-2-3-1}
若 $X$ 适应，$A$ 是状态空间中的可测集，则
\[
  \tau_A=\inf\{n\ge0:X_n\in A\}
\]
是停时，因为
\[
  \{\tau_A\le n\}=\bigcup_{k=0}^n\{X_k\in A\}\in\mathcal F_n.
\]
约定空集的下确界为 $\infty$，所以从未命中也已包含在定义中。
\end{example}

\begin{example}[最后命中不是停时]\label{ex:6-2-3-2}
令 $S_0=0$，$S_n$ 为简单对称随机游走，并只观察到时刻 $2$。定义
\[
  L=\max\{k\in\{0,1,2\}:S_k=0\}.
\]
$S_1$ 必为 $\pm1$；若第二步与第一步反向，则 $L=2$，否则 $L=0$。因此
$\Prob(L\le0)=1/2$。自然滤过的 $\mathcal F_0$ 是平凡 $\sigma$-代数，而
$\{L\le0\}$ 概率既非零也非一，所以该事件不在 $\mathcal F_0$，$L$ 不是停时。
失败的不是“最大值”这个运算，而是判定最后一次命中必须查看未来。
\end{example}

\begin{definition}[可选抽样问题]\label{def:6-2-3-3}
对停时 $\sigma\le\tau$，可选抽样问题询问何时有
\[
  \E[M_\tau\mid\mathcal F_\sigma]=M_\sigma.
\]
该式首先要求停止值可测且可积，还要求从有界截断
$\sigma\wedge n,\tau\wedge n$ 放开时能够交换极限与期望。
\end{definition}

\begin{proposition}[停止保持鞅性]\label{proposition:6-2-3-1}
若 $M$ 是鞅、$\tau$ 是停时，则停止过程 $(M_{n\wedge\tau})_{n\ge0}$ 仍是鞅。
\end{proposition}

\begin{proof}
由 $\{\tau\ge k\}\in\mathcal F_{k-1}$，过程
$H_k=\1_{\{\tau\ge k\}}$ 有界且可预测。逐路径有
\[
 M_{n\wedge\tau}
 =M_0+\sum_{k=1}^n\1_{\{\tau\ge k\}}(M_k-M_{k-1}).
\]
命题~\ref{proposition:6-2-1-4} 因而说明右边是鞅。
\end{proof}

\begin{theorem}[有界可选抽样]\label{theorem:6-2-3-2}
若 $\sigma\le\tau\le N$ 是有界停时，$M$ 是鞅，则
\[
  \E[M_\tau\mid\mathcal F_\sigma]=M_\sigma\quad\text{a.s.}
\]
特别地，$\E M_\tau=\E M_\sigma=\E M_0$。
\end{theorem}

\begin{proof}
停止值可积且 $M_\sigma$ 对 $\mathcal F_\sigma$ 可测。固定
$A\in\mathcal F_\sigma$。有限和恒等式
\[
 M_\tau-M_\sigma
 =\sum_{k=1}^N\1_{\{\sigma<k\le\tau\}}(M_k-M_{k-1})
\]
逐路径成立。对每个 $k$，
\[
 A\cap\{\sigma<k\le\tau\}
 =(A\cap\{\sigma\le k-1\})\cap\{\tau\ge k\}
 \in\mathcal F_{k-1}.
\]
因此
\begin{align*}
 \E[\1_A\1_{\{\sigma<k\le\tau\}}(M_k-M_{k-1})]
 &=\E\!\left[\1_A\1_{\{\sigma<k\le\tau\}}
     \E[M_k-M_{k-1}\mid\mathcal F_{k-1}]\right]\\
 &=0.
\end{align*}
对有限个 $k$ 求和，得到 $\E[M_\tau\1_A]=\E[M_\sigma\1_A]$。这正是所列条件
期望等式的积分刻画。取 $A=\Omega$ 得 $\E M_\tau=\E M_\sigma$；再把已经证明的条件式
用于停时对 $0\le\sigma$，得到 $\E M_\sigma=\E M_0$。
\end{proof}

有界性允许逐项处理有限和。对无界停时，正确的替代条件不是“几乎处处有限”本身，而是保证全部
截断停止值的尾部不会逃逸。

\begin{theorem}[UI 可选抽样]\label{theorem:6-2-3-3}
设 $M$ 是一致可积鞅，$\sigma\le\tau<\infty$ a.s. 是停时。则
$M_\sigma,M_\tau\in L^1$，并且
\[
  \E[M_\tau\mid\mathcal F_\sigma]=M_\sigma\quad\text{a.s.}
\]
\end{theorem}

\begin{proof}
由命题~\ref{proposition:6-2-2-5}，存在 $M_\infty\in L^1$ 使
\[
  M_n=\E[M_\infty\mid\mathcal F_n],
  \qquad M_n\longrightarrow M_\infty\quad\text{于 }L^1\text{ 及 a.s.}
\]
先取任意有界停时 $\rho\le N$。对 $A\in\mathcal F_\rho$，
定理~\ref{theorem:6-2-3-2} 与
$M_N=\E[M_\infty\mid\mathcal F_N]$ 给出
\[
 \E[M_\rho\1_A]=\E[M_N\1_A]=\E[M_\infty\1_A].
\]
故 $M_\rho=\E[M_\infty\mid\mathcal F_\rho]$。令
$B_{\rho,K}=\{|M_\rho|>K\}$。条件模长不等式、Markov 不等式和任意 $R>0$ 给出
\[
 \E[|M_\rho|\1_{B_{\rho,K}}]
 \le \E[|M_\infty|\1_{\{|M_\infty|>R\}}]
      +R\frac{\|M_\infty\|_1}{K}.
\]
先取 $R$ 使第一项任意小，再令 $K\to\infty$；该界与 $\rho$ 无关。所以
\[
 \{M_\rho:\rho\text{ 为有界停时}\}
\]
是一致可积族。

令 $\rho_n=\tau\wedge n$。因为 $\tau<\infty$ a.s.，
$M_{\rho_n}\to M_\tau$ a.s.；上面的 UI 与 Vitali 定理给
$M_{\rho_n}\to M_\tau$ 于 $L^1$，特别 $M_\tau$ 可积。若
$A\in\mathcal F_\tau$，置 $A_n=A\cap\{\tau\le n\}$。若 $m<n$，则
\[
 A_n\cap\{\rho_n\le m\}=A\cap\{\tau\le m\}\in\mathcal F_m;
\]
若 $m\ge n$，则该交集为 $A_n\in\mathcal F_n\subset\mathcal F_m$。故
$A_n\in\mathcal F_{\rho_n}$，于是
\[
 \E[M_{\rho_n}\1_{A_n}]=\E[M_\infty\1_{A_n}].
\]
令 $n\to\infty$：左边用 $L^1$ 收敛，右边用 DCT，得到
$\E[M_\tau\1_A]=\E[M_\infty\1_A]$，即
\[
 M_\tau=\E[M_\infty\mid\mathcal F_\tau].
\]
令 $\eta_n=\sigma\wedge n$。有界停时公式给
$M_{\eta_n}=\E[M_\infty\mid\mathcal F_{\eta_n}]$。因为 $\sigma<\infty$ a.s.，
$M_{\eta_n}\to M_\sigma$ a.s.；有界停时停止值族的统一可积性与 Vitali 定理再给
$L^1$ 收敛。固定 $B\in\mathcal F_\sigma$，置 $B_n=B\cap\{\sigma\le n\}$。与 $A_n$ 的核验
相同，$B_n\in\mathcal F_{\eta_n}$，从而
\[
 \E[M_{\eta_n}\1_{B_n}]=\E[M_\infty\1_{B_n}].
\]
当 $n\to\infty$ 时，左边由 $L^1$ 收敛及
$\E[|M_\sigma|\1_{B\setminus B_n}]\to0$ 收敛到
$\E[M_\sigma\1_B]$；右边由 DCT 收敛到 $\E[M_\infty\1_B]$。故
$M_\sigma=\E[M_\infty\mid\mathcal F_\sigma]$。由
引理~\ref{lemma:6-2-3-stopped-sigma}，
$\mathcal F_\sigma\subset\mathcal F_\tau$；塔式性质因此给出所求等式。
\end{proof}

若只允许停时在一个零集上取 $\infty$，可在该零集上定义
$M_\tau=M_\infty$，上述结论不变。若 $\Prob(\tau=\infty)>0$，则还必须规定“无穷时刻的信息”。
一种安全约定是令
$\mathcal F_\infty=\sigma(\bigcup_n\mathcal F_n)$，并在
定义~\ref{def:6-2-3-2} 的条件之外要求
$A\cap\{\tau=\infty\}\in\mathcal F_\infty$，同时在该事件上置
$M_\tau=M_\infty$。若不加这项终端条件，原定义对 $\{\tau=\infty\}$ 上的集合没有任何约束，
不能无条件声称 $M_\tau$ 对停止信息可测。这是“允许取 $\infty$”时不能省略的边界。

\begin{corollary}[可积停时与有界增量]\label{corollary:6-2-3-bounded-increments}
设 $M$ 是鞅，并且存在常数 $C<\infty$ 使
$|M_k-M_{k-1}|\le C$ a.s. 对所有 $k$ 成立。若停时
$\sigma\le\tau$ 满足 $\E\tau<\infty$，则 $M_\sigma,M_\tau\in L^1$，且
\[
  \E[M_\tau\mid\mathcal F_\sigma]=M_\sigma.
\]
\end{corollary}

\begin{proof}
由 $\sigma\le\tau$ 与 $\E\tau<\infty$，两个停时都 a.s. 有限，而且
\[
 |M_\sigma|\le|M_0|+C\sigma\le|M_0|+C\tau,
 \qquad |M_\tau|\le|M_0|+C\tau,
\]
故两个停止值可积。
由命题~\ref{proposition:6-2-3-1}，$N_n=M_{n\wedge\tau}$ 是鞅。又因
\[
 |N_n-M_\tau|
 \le C(\tau-n)^+,
\]
而 $(\tau-n)^+\downarrow0$ 且被可积变量 $\tau$ 支配，DCT 给
$N_n\to M_\tau$ 于 $L^1$。命题~\ref{proposition:6-2-2-5} 因而说明 $N$ 是 UI 鞅。
由于 $\sigma\le\tau$，有 $N_\sigma=M_\sigma$、$N_\tau=M_\tau$；对 $N$ 应用
定理~\ref{theorem:6-2-3-3} 即得结论。
\end{proof}

\begin{proposition}[随机游走赌徒破产]\label{proposition:6-2-3-4}
令 $(S_n)$ 为从 $i\in\{0,\ldots,N\}$ 出发的简单对称随机游走，并令
\[
  \tau=\inf\{n\ge0:S_n\in\{0,N\}\}.
\]
则
\[
  \Prob_i(S_\tau=N)=\frac{i}{N},
  \qquad
  \E_i\tau=i(N-i).
\]
\end{proposition}

\begin{proof}
边界起点的结论直接成立，以下设 $0<i<N$。令 $\tau_m=\tau\wedge m$。过程
$S_n$ 与 $S_n^2-n$ 都是鞅：后一结论是命题~\ref{proposition:6-2-1-5} 后平方随机游走计算中
$\sigma^2=1$ 的情形。对有界停时 $\tau_m$ 使用
定理~\ref{theorem:6-2-3-2}，得到
\begin{equation}\label{eq:6-2-gambler-truncation}
  \E_i S_{\tau_m}=i,
  \qquad
  \E_i S_{\tau_m}^2-\E_i\tau_m=i^2.
\end{equation}
在命中以前 $S_n\in\{1,\ldots,N-1\}$，命中时位于 $0$ 或 $N$，故
$0\le S_{\tau_m}\le N$。由第二式，
\[
  0\le\E_i\tau_m=\E_iS_{\tau_m}^2-i^2\le N^2.
\]
单调收敛定理先给 $\E_i\tau\le N^2<\infty$，因而 $\tau<\infty$ a.s.。现在才放开截断。
由 $S_{\tau_m}\to S_\tau$ a.s. 及统一界 $N$，DCT 和
\eqref{eq:6-2-gambler-truncation} 给 $\E_iS_\tau=i$。若
$p=\Prob_i(S_\tau=N)$，则 $S_\tau$ 只取 $0,N$，所以 $Np=i$，即 $p=i/N$。

最后，对第二式令 $m\to\infty$；左边的停时项用 MCT，平方项用 DCT，得到
\[
 \E_i\tau=\E_iS_\tau^2-i^2=N^2\frac{i}{N}-i^2=i(N-i).
\]
证明次序很重要：$\E\tau<\infty$ 是先由有界截断推出的，不是预先借用待证公式。
\end{proof}

\begin{proposition}[可积停时下的 Wald 第一恒等式]\label{proposition:6-2-3-5}
设 $\xi_1,\xi_2,\ldots$ i.i.d.，$\E|\xi_1|<\infty$，
$\E\xi_1=\mu$，并令
$\mathcal F_n=\sigma(\xi_1,\ldots,\xi_n)$、
$S_n=\sum_{k=1}^n\xi_k$。若停时 $\tau$ 满足 $\E\tau<\infty$，则
$S_\tau=\sum_{k=1}^{\tau}\xi_k$ 可积，且
\[
  \E S_\tau=\mu\E\tau.
\]
\end{proposition}

\begin{proof}
事件 $\{\tau\ge k\}$ 属于 $\mathcal F_{k-1}$，而 $\xi_k$ 独立于
$\mathcal F_{k-1}$。Tonelli 定理给出
\begin{align*}
 \E\sum_{k=1}^{\tau}|\xi_k|
 &=\sum_{k=1}^\infty\E[\1_{\{\tau\ge k\}}|\xi_k|]\\
 &=\E|\xi_1|\sum_{k=1}^\infty\Prob(\tau\ge k)
 =\E|\xi_1|\E\tau<\infty.
\end{align*}
所以随机级数绝对可积，可以交换求和与期望。再次利用独立性，
\begin{align*}
 \E S_\tau
 &=\sum_{k=1}^\infty
   \E[\1_{\{\tau\ge k\}}\xi_k]
 =\sum_{k=1}^\infty \Prob(\tau\ge k)\E\xi_k
 =\mu\E\tau.
\end{align*}
\end{proof}

\begin{example}[第 $r$ 次 Bernoulli 成功]
令 $\xi_k\sim\mathrm{Bernoulli}(p)$，$0<p<1$，取整数 $r\ge1$，并令
$\tau=\inf\{n:S_n=r\}$。使用 Wald 恒等式以前先验证 $\E\tau<\infty$。
把试验按长度 $r$ 分块；每一整块全部成功的概率为 $p^r$，各块相互独立。若
$\tau>mr$，则前 $m$ 块中没有一块全部成功，故
\[
  \Prob(\tau>mr)\le(1-p^r)^m.
\]
把尾和按这些长度为 $r$ 的块分组，得到
\[
 \E\tau=\sum_{k\ge1}\Prob(\tau\ge k)
 \le r\sum_{m\ge0}(1-p^r)^m=\frac{r}{p^r}<\infty.
\]
命中时 $S_\tau=r$，命题~\ref{proposition:6-2-3-5} 因而给
$r=p\E\tau$，即 $\E\tau=r/p$。粗尾界只负责合法化 Wald；最后的精确答案来自恒等式。
当 $p=1$ 时 $\tau=r$，同一答案可直接读出。
\end{example}

最后看一个专门揭示假设缺口的策略。

\begin{example}[加倍下注策略的失效]\label{ex:6-2-3-3}
令公平硬币结果 $\varepsilon_k$ 在赢时为 $1$、输时为 $-1$，并令
$\tau=\inf\{k\ge1:\varepsilon_k=1\}$。在第 $k$ 局尚未赢时下注
$2^{k-1}$，即取
\[
 H_k=2^{k-1}\1_{\{\tau\ge k\}},
 \qquad W_n=\sum_{k=1}^nH_k\varepsilon_k.
\]
$H_k$ 在第 $k$ 次抛掷前可知；对每个固定 $n$，差分条件均值为零，所以 $(W_n)$ 是鞅。
又因 $\Prob(\tau>n)=2^{-n}$，$\tau<\infty$ a.s.。若首次在第 $k$ 局赢，过去总损失为
$1+2+\cdots+2^{k-2}=2^{k-1}-1$，本局赢得 $2^{k-1}$，故
$W_\tau=1$ a.s.

这不与可选抽样矛盾。事实上
\[
 W_n=
 \begin{cases}
 1,&\tau\le n,\\
 -(2^n-1),&\tau>n,
 \end{cases}
 \qquad \E W_n=0.
\]
在概率 $2^{-n}$ 的尾事件上，损失大小接近 $2^n$；因此对任意 $K$，取足够大的 $n$ 有
\[
 \E[|W_n|\1_{\{|W_n|>K\}}]
 \ge (2^n-1)2^{-n}\longrightarrow1.
\]
所以 $(W_n)$ 不一致可积。a.s. 有限的停时并未阻止期望质量沿越来越稀有的失败路径逃到无穷。
\end{example}

\begin{figure}[htbp]
\centering
\begin{tikzpicture}[x=1.15cm,y=0.72cm]
  \draw[->,BookInk!75] (0,0)--(8.4,0) node[right] {时间};
  \foreach \x in {0,...,7}{
    \draw[BookInk!45] (\x,-0.12)--(\x,0.12);
    \node[below] at (\x,-0.12) {\scriptsize $\x$};
  }
  \fill[BookAccent!12] (2.05,0.35) rectangle (6.95,1.1);
  \node[book label] at (4.5,0.72) {$\{\sigma<k\le\tau\}$：第 $k$ 步前即可判定是否仍在抽样};
  \draw[BookAccent,thick] (2,0.25)--(2,1.25) node[above] {$\sigma$};
  \draw[BookAccent,thick] (7,0.25)--(7,1.25) node[above] {$\tau$};
  \begin{scope}[yshift=-1.65cm]
    \draw[BookWarm,thick] (0,0)--(8,0) node[right] {$0$};
    \draw[BookWarm,thick] (0,2)--(8,2) node[right] {$N$};
    \draw[BookInk,thick] plot coordinates
      {(0,0.8) (1,1.2) (2,0.8) (3,1.2) (4,1.6) (5,1.2) (6,1.6) (7,2)};
    \fill[BookWarm] (7,2) circle (2pt) node[above left] {首次命中};
  \end{scope}
\end{tikzpicture}
\caption{上图的随机区间给出可选抽样证明中的可预测系数；下图是停于 $0,N$ 的赌徒破产路径。}
\label{fig:6-2-3-stopping}
\end{figure}

有界停时、UI 鞅和 Wald 定理中的“可积停时加 i.i.d. 可积增量”分别是三组不同的充分条件。
它们虽然都导出停止值的期望恒等式，却不可相互替代。尤其，$\mathcal F_\tau$ 不是
$\sigma(\tau)$，停止值可积也不能由把 $n$ 换成 $\tau$ 的符号操作得到；每次放开截断都必须明确
使用 DCT、MCT 或 UI/Vitali 中的一种。

\begin{exercise}
证明命题~\ref{proposition:6-2-3-1} 时，不引用可预测变换，直接分别计算
$\{\tau\le n\}$ 与 $\{\tau>n\}$ 上的
$\E[M_{(n+1)\wedge\tau}\mid\mathcal F_n]$。
\end{exercise}

\begin{exercise}
设 $\sigma\le\tau\le N$。证明事件
$\{\sigma<k\le\tau\}$ 属于 $\mathcal F_{k-1}$，并用这一事实独立重证有界可选抽样。
\end{exercise}

\begin{exercise}
逐项复核命题~\ref{proposition:6-2-3-4}：先由平方鞅得到
$\sup_m\E\tau_m<\infty$，再说明两个极限分别使用 MCT 与 DCT。若随机游走每步向右概率为
$p\ne1/2$，指出 $S_n$ 的哪一个指数函数可能替代线性鞅。
\end{exercise}

\begin{exercise}
在例~\ref{ex:6-2-3-3} 中直接计算 $\E|W_n-W_\tau|$，说明 a.s. 收敛为什么不升级为
$L^1$ 收敛，并指出定理~\ref{theorem:6-2-3-3} 的哪项假设失败。
\end{exercise}

鞅把完整过去压缩成一个条件平均。下一节将研究另一种压缩：对 Markov 过程，给定当前状态以后，
完整过去不再改变未来的条件分布。转移核和半群正是把这种状态压缩组织成可迭代算子的语言。

\section{Markov 演化、半群与跳过程}\label{sec:06-03}

\subsection{Markov 核、Chapman--Kolmogorov 与转移半群}\label{subsec:6-3-1}
\index{Markov 核}\index{Chapman--Kolmogorov 方程}\index{转移半群}

知道过程的当前状态，并不总能预测下一步。若
\[
 X_{n+1}=X_n+X_{n-1}+\xi_{n+1},
\]
其中 $\xi_{n+1}$ 是新的独立噪声，那么给定 $X_n$ 仍不足以确定下一步的条件分布；至少还要记住
$X_{n-1}$。当然可以把状态扩成 $(X_{n-1},X_n)$，但这也说明了 Markov 性的真正内容：所选的
“当前状态”必须已经包含预测未来所需的全部过去信息。正则条件分布只保证条件分布可以写成核，
却不会替我们证明这一充分性。

先看三个核，它们依次把确定演化、离散噪声与扩散放进同一格式。

\begin{example}[确定性动力系统]\label{ex:6-3-1-1}
设 $(E,\mathcal E)$ 为可测空间，$T:E\to E$ 可测。定义
\[
 K(x,A)=\delta_{T(x)}(A)=\1_A(T(x)).
\]
固定 $x$ 时这是概率测度；固定 $A$ 时它关于 $x$ 可测。对有界可测 $f$，
\[
 Kf(x)=\int_E f(y)K(x,\dd y)=f(T(x)).
\]
因此确定性轨道 $x,T(x),T^2(x),\ldots$ 是 Markov 演化的退化情形，且
$K^nf=f\circ T^n$。
\end{example}

\begin{example}[加性噪声]\label{ex:6-3-1-2}
令 $E=\R^d$，$F:E\to E$ 为 Borel 映射，$\xi_1,\xi_2,\ldots$ 独立同分布，共同分布为
$\nu$，并与 $X_0$ 独立。递推式
\[
 X_{n+1}=F(X_n)+\xi_{n+1}
\]
给出
\[
 K(x,A)=\nu(A-F(x)),\qquad
 Kf(x)=\int_{\R^d}f(F(x)+z)\nu(\dd z).
\]
可测性先对 $f=\1_A$ 由 $(x,z)\mapsto F(x)+z$ 的联合可测性和 Tonelli 得到；简单函数与
单调收敛再给一般非负 $f$。独立性与取出已知因子说明
$\E[f(X_{n+1})\mid\mathcal F_n]=Kf(X_n)$，其中
$\mathcal F_n=\sigma(X_0,\xi_1,\ldots,\xi_n)$。
\end{example}

\begin{example}[Gaussian 热核]\label{ex:6-3-1-3}
沿用定义~\ref{def:4-4-2-2} 的热核 $p_t$，并令
\[
 q_t(x)=p_{t/2}(x)=(2\pi t)^{-d/2}e^{-|x|^2/(2t)}.
\]
它是 $N(0,tI_d)$ 的密度。由
$p_s*p_t=p_{s+t}$ 得 $q_s*q_t=q_{s+t}$。于是
\[
 P_tf(x)=\int_{\R^d}f(x+y)q_t(y)\dd y
\]
满足 $P_sP_t=P_{s+t}$，保持正函数和常数，并且
$\|P_tf\|_\infty\leq\|f\|_\infty$。按照本小节稍后在
评注~\ref{proposition:6-3-1-4} 中给出的生成元定义，这里的时间规范使生成元在光滑函数上为
$\tfrac12\Delta$；第 4 章的 $p_t$ 规范则给生成元 $\Delta$。
\end{example}

\begin{definition}[Markov 核]\label{def:6-3-1-1}
从可测空间 $(E,\mathcal E)$ 到 $(F,\mathcal F)$ 的 \emph{Markov 核}是函数
$K:E\times\mathcal F\to[0,1]$，满足：
\begin{enumerate}
 \item 对每个 $x\in E$，$A\mapsto K(x,A)$ 是 $(F,\mathcal F)$ 上的概率测度；
 \item 对每个 $A\in\mathcal F$，$x\mapsto K(x,A)$ 是 $\mathcal E$-可测函数。
\end{enumerate}
若 $f:F\to\R$ 非负或有界可测，定义
\[
 Kf(x)=\int_F f(y)K(x,\dd y).
\]
若 $\mu$ 是 $E$ 上的概率，定义
\[
 \mu K(A)=\int_EK(x,A)\mu(\dd x).
\]
前一个运算把函数从 $F$ 拉回 $E$，后一个运算把测度从 $E$ 推向 $F$；两者方向相反。
\end{definition}

\begin{definition}[Markov 性]\label{def:6-3-1-2}
设 $(X_n)$ 对滤过 $(\mathcal F_n)$ 适应并取值于 $(E,\mathcal E)$。若存在核
$K:E\leadsto E$，使每个有界可测 $f:E\to\R$ 都满足
\begin{equation}\label{eq:6-3-1-markov-property}
 \E[f(X_{n+1})\mid\mathcal F_n]=Kf(X_n)\qquad\text{a.s.},
\end{equation}
则称 $X$ 是相对于该滤过、转移核为 $K$ 的时间齐次 Markov 链。时间非齐次情形以
单步核 $K_n$ 表示从时刻 $n$ 到 $n+1$ 的转移，即把右端改为 $K_nf(X_n)$。
\end{definition}

多步核不能仅凭某个给定初始分布下的条件期望版本，在未到达状态上逐点识别；因此多步核取为单步核的复合。

若已经指定一个从 $x$ 出发、满足 $\Prob_x(X_0=x)=1$ 的 Markov 分布，就把相应的概率与期望
记为 $\Prob_x$、$\E_x$；初始分布为 $\mu$ 时相应记为 $\Prob_\mu$、$\E_\mu$。这套下标记号只记录
所选初始分布，不自行断言该分布的存在或唯一性。

\begin{remark}[条件分布与标准 Borel 假设]\label{remark:6-3-1-rcd}
若 $E$ 是标准 Borel 空间，定理~\ref{theorem:5-2-2-1} 可为 $X_{n+1}$ 给定 $X_n$ 选择
正则条件分布；这说明可以选出一个以当前状态为参数的可测核。式~\eqref{eq:6-3-1-markov-property} 还要求给定整个
$\mathcal F_n$ 后没有额外变化，这一要求不是正则条件分布存在性的推论。若扩大滤过并泄露未来噪声，
原来的 Markov 等式也可能失效。
若 $E$ 不是标准 Borel 空间，条件期望等式仍可作为 Markov 性的定义，但把逐个集合得到的条件概率
同时选成一个关于状态联合可测的核，可能做不到。例~\ref{ex:6-3-1-2} 的显式公式
$K(x,A)=\nu(A-F(x))$ 同时绕过了这两层选择问题。
\end{remark}

核可以逐步拼接。下面的证明把“对集合可测”和“对函数积分”两层都写出来。

\begin{proposition}[核的复合]\label{proposition:6-3-1-1}
若 $K:E\leadsto F$、$L:F\leadsto G$ 是 Markov 核，则
\[
 (KL)(x,A)=\int_F L(y,A)K(x,\dd y),\qquad A\in\mathcal G,
\]
定义一个核 $KL:E\leadsto G$。对每个非负或有界可测 $f:G\to\R$，
\[
 (KL)f=K(Lf).
\]
\end{proposition}

\begin{proof}
固定 $A\in\mathcal G$。函数 $y\mapsto L(y,A)$ 可测且非负；核积分的简单函数逼近说明
$x\mapsto(KL)(x,A)$ 可测。固定 $x$，若 $(A_j)$ 两两不交，则 Tonelli 定理给
\begin{align*}
 (KL)\left(x,\bigcup_jA_j\right)
 &=\int_F\sum_jL(y,A_j)K(x,\dd y)\\
 &=\sum_j\int_FL(y,A_j)K(x,\dd y)
 =\sum_j(KL)(x,A_j).
\end{align*}
此外 $(KL)(x,G)=1$，故 $A\mapsto(KL)(x,A)$ 为概率。

对 $f=\1_A$，函数恒等式就是定义。线性性把它推广到非负简单函数，单调收敛再推广到非负
可测函数；有界实值函数对正负部应用该结论，复值函数再对实部、虚部应用。
\end{proof}

\begin{theorem}[Chapman--Kolmogorov 方程]\label{theorem:6-3-1-2}
设时间非齐次 Markov 链的单步核为 $(K_n)_{n\ge0}$。置
\[
 K_{m,m}=I,\qquad K_{m,n}=K_mK_{m+1}\cdots K_{n-1}\quad(m<n).
\]
则对每个有界可测 $f$，
\[
 \E[f(X_n)\mid\mathcal F_m]=K_{m,n}f(X_m)\quad\text{a.s.}\qquad(m\le n),
\]
并且这些逐点定义的核满足
\[
 K_{m,n}=K_{m,k}K_{k,n}\qquad(m\leq k\leq n).
\]
时间齐次离散链的 $n$ 步核为 $K^n$；若 $X_0$ 的分布是 $\mu$，则
\[
 \Law(X_n)=\mu K^n,
 \qquad
 \E_\mu f(X_n)=\int_EK^nf(x)\mu(\dd x)
\]
对每个有界可测 $f$ 成立。
\end{theorem}

\begin{proof}
先核验核复合的结合律。对任意三个可复合核 $K,L,R$ 及任意有界可测 $f$，
命题~\ref{proposition:6-3-1-1} 给出
\[
 ((KL)R)f=(KL)(Rf)=K(L(Rf))=K((LR)f)=(K(LR))f.
\]
取 $f=\1_A$，便有 $((KL)R)(x,A)=(K(LR))(x,A)$ 对每个 $x,A$ 成立。因此上式中有限个
单步核的复合无须加括号，并且按 $m\le k\le n$ 分组立即得到
$K_{m,n}=K_{m,k}K_{k,n}$；这是逐点核恒等式，不受某个初始分布的零质量状态影响。

再证明条件期望公式。$n=m$ 时公式是适应性给出的
$\E[f(X_m)\mid\mathcal F_m]=f(X_m)$；$n=m+1$ 时就是 Markov 性。若公式对
$(m+1,n)$ 成立，则塔式性质及单步 Markov 等式给
\begin{align*}
 \E[f(X_n)\mid\mathcal F_m]
 &=\E\bigl[K_{m+1,n}f(X_{m+1})\mid\mathcal F_m\bigr]\\
 &=K_m(K_{m+1,n}f)(X_m)=K_{m,n}f(X_m).
\end{align*}
对 $n-m$ 归纳即得一般公式。

时间齐次时对 $n$ 归纳得 $K_{0,n}=K^n$。最后
\[
 \Prob_\mu(X_n\in A)
 =\int_EK^n(x,A)\mu(\dd x)=(\mu K^n)(A),
\]
函数公式由简单函数与单调收敛得到。
\end{proof}

上面的证明把多步预测压成终点函数。若还要观察中间路径，核仍可递归积分。

\begin{proposition}[Markov 条件期望的路径版本]\label{proposition:6-3-1-3}
设 $X$ 是转移核为 $K$ 的时间齐次 Markov 链。对 $m\geq1$ 和有界可测
$F:E^m\to\R$，存在有界可测函数 $G_F:E\to\R$，使
\[
 \E[F(X_{n+1},\ldots,X_{n+m})\mid\mathcal F_n]=G_F(X_n)\quad\text{a.s.}
\]
具体地，$G_F(x)$ 由从最末坐标向前逐次对 $K$ 积分得到。
\end{proposition}

\begin{proof}
先把“逐次积分”写成无歧义的公式。对有界可测
$F:E^m\to\R$，定义
\begin{equation}\label{eq:6-3-1-path-kernel}
 G_F(x_0)=
 \int_E\!\cdots\!\int_E
 F(x_1,\ldots,x_m)
 K(x_{m-1},\dd x_m)\cdots K(x_0,\dd x_1).
\end{equation}
该公式可从最后一个变量向前归纳理解。具体地，参数核积分的可测性与
命题~\ref{proposition:6-3-1-1} 的证明相同：对示性函数成立，再经简单函数和单调收敛
推广。逐层使用这一事实，可知式~\eqref{eq:6-3-1-path-kernel} 的每一层积分都是有界可测
函数，特别地 $G_F$ 有界可测。

先取矩形乘积函数
$F(x_1,\ldots,x_m)=\prod_{j=1}^mf_j(x_j)$，其中每个 $f_j$ 有界可测。令
\[
 g_m=f_m,\qquad g_j(x)=f_j(x)K g_{j+1}(x)\quad(j=m-1,\ldots,1).
\]
此时式~\eqref{eq:6-3-1-path-kernel} 给出 $G_F=Kg_1$。从 $X_{n+m}$ 开始，反复使用塔式性质和
式~\eqref{eq:6-3-1-markov-property}，得到
\[
 \E\!\left[\prod_{j=1}^mf_j(X_{n+j})\middle|\mathcal F_n\right]
 =(Kg_1)(X_n)=G_F(X_n).
\]

令 $\mathcal H$ 为使所述条件等式以式~\eqref{eq:6-3-1-path-kernel} 中的 $G_F$ 成立的有界可测
$F$ 的集合。条件期望与核积分的线性性说明 $\mathcal H$ 是向量空间。若
$0\le F_j\uparrow F$ 且整列一致有界，则条件单调收敛和式
\eqref{eq:6-3-1-path-kernel} 中的逐层单调收敛同时给出 $F\in\mathcal H$。矩形示性函数已属于
$\mathcal H$，函数单调类定理遂给出全部有界可测 $F$。
\end{proof}

确定时刻的 Markov 性还可以停在一个随机时刻，但离散性在证明中承担了关键的可数分片。

\begin{theorem}[离散时间强 Markov 性]\label{theorem:6-3-1-strong-markov}
设 $X$ 是相对于 $(\mathcal F_n)$ 的时间齐次 Markov 链，$\tau<\infty$ a.s. 是同一滤过的停时。
则对 $m\geq0$ 与有界可测 $f$，
\[
 \E[f(X_{\tau+m})\mid\mathcal F_\tau]=K^mf(X_\tau)\quad\text{a.s.}
\]
更一般地，停时后的任意有界有限柱函数由命题~\ref{proposition:6-3-1-3} 的递归核从
$X_\tau$ 起计算。
\end{theorem}

\begin{proof}
记 $Y=K^mf(X_\tau)$。它对 $\mathcal F_\tau$ 可测：对 Borel 集 $B$ 与整数 $n$，
\[
 \{Y\in B\}\cap\{\tau\leq n\}
 =\bigcup_{k=0}^n\{K^mf(X_k)\in B\}\cap\{\tau=k\}\in\mathcal F_n.
\]
任取 $A\in\mathcal F_\tau$。由停止信息的定义，
$A\cap\{\tau=n\}\in\mathcal F_n$。确定时刻 Markov 性给
\begin{align*}
 \E[\1_Af(X_{\tau+m})]
 &=\sum_{n=0}^\infty
   \E[\1_{A\cap\{\tau=n\}}f(X_{n+m})]\\
 &=\sum_{n=0}^\infty
   \E[\1_{A\cap\{\tau=n\}}K^mf(X_n)]
 =\E[\1_AY].
\end{align*}
级数可交换是因为 $f$ 有界而各片两两不交。这个积分恒等式正是条件期望的定义。

有限未来柱函数先对矩形函数逐层使用同一计算，再以函数单调类推广。证明依赖的是可数值停时的分片；
连续时间强 Markov 性还需要路径连续性和滤过右连续性，不能从这里直接类推。
\end{proof}

\begin{definition}[转移半群]\label{def:6-3-1-3}
设 $(K_t)_{t\geq0}$ 是 $E$ 上的 Markov 核族，满足
\[
 K_0=I,\qquad K_{s+t}=K_sK_t\quad(s,t\geq0).
\]
若适应过程 $X=(X_t)_{t\geq0}$ 对每个有界可测 $f$ 及 $0\leq s\leq t$ 满足
\[
 \E[f(X_t)\mid\mathcal F_s]=K_{t-s}f(X_s)\quad\text{a.s.},
\]
则称 $X$ 是相对于 $(\mathcal F_t)$、转移核族为 $(K_t)$ 的连续时间齐次 Markov 过程。
定义算子 $P_tf=K_tf$。若对每个 $x$ 还指定了满足
$\Prob_x(X_0=x)=1$ 的过程分布，且
$K_t(x,A)=\Prob_x(X_t\in A)$，则同一算子也写成
$P_tf(x)=\E_xf(X_t)$。核的 Chapman--Kolmogorov 恒等式给出
$P_0=I$ 与 $P_{s+t}=P_sP_t$。每个 $P_t$ 保持非负函数和常数 $1$，并满足
$\|P_tf\|_\infty\leq\|f\|_\infty$；这样的算子族称为转移半群或 Markov 半群。
\end{definition}

Poisson 过程、一致化构造与 Brownian motion 的转移半群都可由独立增量或显式核验证这一恒等式。
连续时间的半群律是额外结构，不是离散时间定理~\ref{theorem:6-3-1-2} 的形式推论。

\begin{remark}[Feller 性与生成元的边界]\label{proposition:6-3-1-4}
若 $E$ 为局部紧 Hausdorff 空间，且 $P_tC_0(E)\subset C_0(E)$，则称 $(P_t)$ 具有 Feller 性。
若进一步对每个 $f\in C_0(E)$ 有 $\|P_tf-f\|_\infty\to0$，它便是第 4 章意义下的收缩
$C_0$ 半群，可在
\[
 D(L)=\left\{f:\lim_{t\downarrow0}\frac{P_tf-f}{t}
 \text{ 在 }C_0(E)\text{ 中存在}\right\}
\]
定义生成元。半群律本身既不推出 Feller 性，也不推出强连续性。有限链和一致有界跳率链提供了可直接验证的例子；
一般 Feller 过程的构造与生成元刻画则需要更强的假设。
\end{remark}

\begin{figure}[htbp]
\centering
\begin{tikzpicture}[node distance=13mm and 18mm]
 \node[book strong node] (x) {状态 $x$};
 \node[book node,right=of x] (y) {中间状态 $y$};
 \node[book strong node,right=of y] (z) {状态 $z$};
 \draw[book accent arrow] (x)--node[above,book label]{$K_s(x,\dd y)$}(y);
 \draw[book accent arrow] (y)--node[above,book label]{$K_t(y,\dd z)$}(z);
 \draw[book dashed arrow,bend right=28] (x) to
   node[below,book label]{$K_{s+t}(x,\dd z)=\int K_t(y,\dd z)K_s(x,\dd y)$}(z);
 \node[book warm node,below=16mm of y] (f) {函数方向\\$P_{s+t}f=P_s(P_tf)$};
 \draw[book arrow] (z)--(f);
 \draw[book arrow] (f)--(x);
\end{tikzpicture}
\caption{核沿时间正向复合；函数从终点向起点反向积分，同一恒等式写成半群律。}
\label{fig:6-3-1-kernel-semigroup}
\end{figure}

这个结构有三条容易混淆的边界。第一，$Kf$ 和 $\mu K$ 的方向相反。第二，Markov 性相对于给定滤过
表述；加入未来信息可能破坏它。第三，半群律只组织有限维转移分布，不保证样本路径连续，甚至不保证
$t\mapsto P_tf$ 强连续。

\begin{exercise}
\begin{enumerate}
 \item 逐项验证加性噪声核 $K(x,A)=\nu(A-F(x))$ 的两个核条件。
 \item 证明 $K^nf(x)=\E_xf(X_n)$，并写出 $K^2(x,A)$ 的二重积分。
 \item 对热核 $q_t$ 直接完成平方，重新证明 $q_s*q_t=q_{s+t}$。
 \item 构造一个 Markov 链及一个泄露 $X_{n+1}$ 的扩大滤过，使链不再满足相同的 Markov 条件等式。
\end{enumerate}
\end{exercise}

核可以有不变分布，链也可能从任意初态逐渐忘掉起点；这两句话并不等价。下一小节把不变、
分布收敛和时间平均分开处理。

\subsection{不变分布、可逆性与 Markov 链遍历}\label{subsec:6-3-2}
\index{不变分布}\index{平稳 Markov 链}\index{可逆性}\index{遍历定理}

若 $X_0$ 的分布为 $\mu$，则 $X_n$ 的分布是 $\mu P^n$。一个自然问题是寻找
$\pi P=\pi$；然而即使找到了这样的 $\pi$，仍有两件事没有回答：从给定初态出发的
$P^n(i,\cdot)$ 是否趋于 $\pi$，以及单条路径的时间平均是否趋于 $\pi$-平均。周期链会否定第一句
而保留第二句，多个闭沟通类则会同时破坏唯一性和全局收敛。

\begin{definition}[不变分布与平稳链]\label{def:6-3-2-1}
若状态空间上的概率 $\pi$ 满足 $\pi P=\pi$，则称 $\pi$ 是转移核 $P$ 的\emph{不变分布}。
若 $X_0\sim\pi$，则每个 $X_n\sim\pi$；更进一步，对任意
$0\leq n_1<\cdots<n_k$ 与 $r\geq0$，Markov 核的迭代公式给
\[
 \Law(X_{n_1},\ldots,X_{n_k})
 =\Law(X_{n_1+r},\ldots,X_{n_k+r}).
\]
为了看清最后一步，对任意可测集 $A_1,\ldots,A_k$，路径核公式给出
\begin{align*}
 &\Prob(X_{n_1+r}\in A_1,\ldots,X_{n_k+r}\in A_k)\\
 &\quad=\int_E\!\cdots\!\int_E\pi(\dd x_0)
   P^{n_1+r}(x_0,\dd x_1)\1_{A_1}(x_1)
   \prod_{j=2}^k P^{n_j-n_{j-1}}(x_{j-1},\dd x_j)\1_{A_j}(x_j)\\
 &\quad=\int_E\!\cdots\!\int_E\pi(\dd x_1)\1_{A_1}(x_1)
   \prod_{j=2}^k P^{n_j-n_{j-1}}(x_{j-1},\dd x_j)\1_{A_j}(x_j),
\end{align*}
其中第二个等号使用 $\pi P^{n_1+r}=\pi$；对未平移的时刻组同样使用
$\pi P^{n_1}=\pi$ 得到同一积分。矩形是乘积 $\sigma$-代数的生成 $\pi$-系统，因而上式确定整个
联合分布。这样的链称为平稳链。平稳只表示分布对时间平移不变，不表示各时刻相互独立。
\end{definition}

\begin{definition}[可逆性]\label{def:6-3-2-2}
若 $\pi$ 是状态空间上的概率，并且 $E\times E$ 上有测度恒等
\begin{equation}\label{eq:6-3-2-detailed-balance}
 \pi(\dd x)P(x,\dd y)=\pi(\dd y)P(y,\dd x),
\end{equation}
则称 $P$ 相对于 $\pi$ \emph{可逆}。在有限或可数状态空间上，这等价于细致平衡（detailed balance）关系
\[
 \pi(i)P(i,j)=\pi(j)P(j,i)\qquad(i,j\in E).
\]
对式~\eqref{eq:6-3-2-detailed-balance} 的两边关于 $x$ 积分，得到
$\pi P=\pi$；所以可逆性蕴含不变性，反向一般不成立。
\end{definition}

\begin{definition}[不可约与周期]\label{def:6-3-2-3}
有限状态链称为\emph{不可约}，若任意 $i,j$ 之间都存在 $n\geq0$ 使 $P^n(i,j)>0$。
状态 $i$ 的周期为
\[
 d(i)=\gcd\{n\geq1:P^n(i,i)>0\}.
\]
不可约链中所有状态的周期相同；共同周期为 $1$ 时称链非周期。
\end{definition}

上述“共同周期”需要由不可约性推出。任取状态 $i,j$，选择
$r,s\ge0$ 使 $P^r(i,j)>0$ 且 $P^s(j,i)>0$。若 $n\ge1$ 且 $P^n(j,j)>0$，则
Chapman--Kolmogorov 方程给出
\[
 P^{r+s}(i,i)\ge P^r(i,j)P^s(j,i)>0,
\]
以及
\[
 P^{r+n+s}(i,i)
 \ge P^r(i,j)P^n(j,j)P^s(j,i)>0.
\]
因而 $d(i)$ 同时整除 $r+s$ 和 $r+n+s$，也就整除 $n$。这对每个满足
$P^n(j,j)>0$ 的 $n$ 成立，故 $d(i)\mid d(j)$。交换 $i,j$ 得
$d(j)\mid d(i)$，从而 $d(i)=d(j)$。

\begin{example}[二状态链]\label{ex:6-3-2-1}
令
\[
 P=\begin{pmatrix}1-a&a\\ b&1-b\end{pmatrix},
 \qquad 0\leq a,b\leq1,\quad a+b>0.
\]
解 $\pi P=\pi$ 与 $\pi_0+\pi_1=1$，得到
\[
 \pi=\left(\frac b{a+b},\frac a{a+b}\right).
\]
向量 $(1,-1)$ 是特征值 $1-a-b$ 的左特征向量；任意两个概率行向量之差都与
$(1,-1)$ 成比例，因而
\[
 \mu P^n-\pi=(1-a-b)^n(\mu-\pi).
\]
当 $a,b>0$ 时链不可约。若 $a=b=1$，链在两个状态间确定性交替，周期为 $2$，所以
$P^n(0,\cdot)$ 不收敛；但其 \Cesaro{} 平均趋于 $(1/2,1/2)$。其余不可约情形满足
$|1-a-b|<1$，上式给出几何收敛率。
\end{example}

\begin{example}[有限无向图上的随机游走]\label{ex:6-3-2-2}
设 $G=(V,E_G)$ 是有限连通无向图，$d(x)$ 为顶点度。简单随机游走的转移概率为
$P(x,y)=d(x)^{-1}$（若 $x\sim y$），否则为零。令
\[
 \pi(x)=\frac{d(x)}{\sum_{z\in V}d(z)}.
\]
当 $x\sim y$ 时
\[
 \pi(x)P(x,y)=\frac1{\sum_zd(z)}=\pi(y)P(y,x),
\]
其余情形两边同为零，所以链可逆。若图为二分图，每一步都从一个部跳到另一个部，周期为 $2$；
把转移改为 $(I+P)/2$，每个状态都有正的自环概率，周期便降为 $1$，不变分布保持不变。
\end{example}

\begin{example}[多个闭类]\label{ex:6-3-2-3}
在状态集 $\{0,1\}$ 上取 $P=I$。两个状态都是吸收态，每个
$\theta\delta_0+(1-\theta)\delta_1$ 都是不变分布；从 $i$ 出发的分布恒为 $\delta_i$。
所以“不变分布存在”既不保证唯一，也不保证极限与初态无关。缺失的正是不可约性。
\end{example}

\begin{example}[无限状态中不可约仍不够]\label{ex:6-3-2-4}
在 $\Z$ 上的简单对称随机游走不可约。若 $\pi$ 是不变概率，则
\[
 \pi(k)=\frac{\pi(k-1)+\pi(k+1)}2,
\]
故差分 $\pi(k+1)-\pi(k)$ 为常数。一个在整条 $\Z$ 上非负的仿射序列只能是常数，而非零常数
不可求和，所以不存在不变概率。将一步改为以 $1/2$ 留在原地、以 $1/4$ 各向左右移动，只把上式
两边同时加入 $\pi(k)/2$，结论不变；此链已经非周期。计数测度虽然可逆，却有无限总质量。
这说明有限状态结论不能只凭“不可约、非周期”移植到一般状态空间；还必须另行证明足以控制
回返概率与回返时间的条件。
\end{example}

\begin{figure}[htbp]
\centering
\begin{tikzpicture}[node distance=12mm and 18mm]
  \node[book strong node] (a0) {$0$};
  \node[book strong node,right=of a0] (a1) {$1$};
  \draw[book accent arrow,bend left=20] (a0) to node[above,book label]{1} (a1);
  \draw[book accent arrow,bend left=20] (a1) to node[below,book label]{1} (a0);
  \node[book label,above=7mm of a0,xshift=9mm] {闭类 $A={0,1}$：周期 $2$};

  \node[book warm node,right=34mm of a1] (b) {$2$};
  \draw[book accent arrow,loop above] (b) to node[above,book label]{1} (b);
  \node[book label,above=7mm of b] {闭类 $B={2}$：周期 $1$};

  \node[book node,below=18mm of a1,xshift=19mm] (tr) {瞬时态 $3$};
  \draw[book dashed arrow] (tr) to node[left,book label]{正概率} (a1);
  \draw[book dashed arrow] (tr) to node[right,book label]{正概率} (b);
\end{tikzpicture}
\caption{两个闭沟通类使不变分布不唯一；左侧的二循环又显示，即使限制在一个不可约闭类内，
周期仍可使单时刻分布来回振荡。定理~\ref{theorem:6-3-2-1}--\ref{theorem:6-3-2-3}
分别给出不可约情形下的不变性、单时刻收敛与时间平均。}
\label{fig:6-3-2-classes-period}
\end{figure}

有限状态不可约性提供一个很实用的几何尾估计。固定状态 $i_*$。对每个 $x$ 选择一条从 $x$
到 $i_*$ 的正概率有限路径。状态有限，因而存在 $m<\infty$ 与 $\varepsilon>0$，使从任意 $x$
出发在 $m$ 步以内命中 $i_*$ 的概率至少为 $\varepsilon$。更明确地，设为 $x$ 选取的路径长度为
$m_x$、路径概率乘积为 $p_x>0$，则可取
\[
 m=\max_xm_x,\qquad \varepsilon=\min_xp_x>0.
\]
若 $\tau_{i_*}=\inf\{n\geq0:X_n=i_*\}$，则 Markov 性在确定时刻 $km$ 给出
\begin{align*}
 \Prob_x(\tau_{i_*}>(k+1)m)
 &=\E_x\!\left[
   \1_{\{\tau_{i_*}>km\}}
   \Prob_{X_{km}}(\tau_{i_*}>m)
   \right]\\
 &\le(1-\varepsilon)\Prob_x(\tau_{i_*}>km).
\end{align*}
从 $k=0$ 归纳便得
\begin{equation}\label{eq:6-3-2-hitting-geometric-tail}
 \sup_x\Prob_x(\tau_{i_*}>km)\leq(1-\varepsilon)^k.
\end{equation}
特别地，所有命中时间和从 $i_*$ 出发的首次正回返时间
$\tau_{i_*}^+=\inf\{n\geq1:X_n=i_*\}$ 都有有限期望。
由于从 $i_*$ 出发的 $\tau_{i_*}$ 等于零，式~\eqref{eq:6-3-2-hitting-geometric-tail}
本身不直接控制正回返时间。由第一步后的 Markov 性，
若 $X_1=y$，则 $\tau_{i_*}^+$ 的剩余部分与从 $y$ 出发命中 $i_*$ 的时间同分布，故
\[
 \E_{i_*}\tau_{i_*}^+
 =1+\sum_{y\in E}P(i_*,y)\E_y\tau_{i_*}<\infty.
\]
这里每一项由式~\eqref{eq:6-3-2-hitting-geometric-tail} 可积，且状态集有限。

\begin{theorem}[有限不可约链的不变分布]\label{theorem:6-3-2-1}
有限状态不可约 Markov 链存在唯一的不变分布 $\pi$，并且 $\pi(j)>0$ 对每个状态 $j$ 成立。
固定 $i\in E$ 时，它由一次回返周期的占用时间给出：
\begin{equation}\label{eq:6-3-2-cycle-invariant-law}
 \pi(j)=
 \frac{\E_i\sum_{k=0}^{\tau_i^+-1}\1_{\{X_k=j\}}}
      {\E_i\tau_i^+}.
\end{equation}
\end{theorem}

\begin{proof}
式~\eqref{eq:6-3-2-hitting-geometric-tail} 说明分母有限且为正。记
\[
 u(j)=\E_i\sum_{k=0}^{\tau_i^+-1}\1_{\{X_k=j\}},
 \qquad m_i=\E_i\tau_i^+.
\]
由于 $\sum_ju(j)=m_i$，式~\eqref{eq:6-3-2-cycle-invariant-law} 定义一个概率。对任意 $\ell$，
把随机上限写成非负指标函数后，条件期望和 Tonelli 定理给
\begin{align*}
 \sum_ju(j)P(j,\ell)
&=\sum_{k=0}^\infty
  \E_i\!\left[\1_{\{k<\tau_i^+\}}P(X_k,\ell)\right]\\
&=\sum_{k=0}^\infty
  \E_i\!\left[\1_{\{k<\tau_i^+\}}
   \E_i[\1_{\{X_{k+1}=\ell\}}\mid\mathcal F_k]\right]\\
 &=\E_i\sum_{k=1}^{\tau_i^+}\1_{\{X_k=\ell\}}.
\end{align*}
这里 $\{k<\tau_i^+\}\in\mathcal F_k$，所以第三行取出已知因子的步骤合法。
最后一个和与 $u(\ell)$ 的差为
$\1_{\{X_{\tau_i^+}=\ell\}}-\1_{\{X_0=\ell\}}=0$，因为周期首尾都等于 $i$。
故 $uP=u$，从而 $\pi P=\pi$。从 $i$ 到 $j$ 可选择一条内部不再经过 $i$ 的简单正概率路径，
所以一个回返周期以正概率访问 $j$，即 $u(j)>0$。

只剩唯一性。若 $\rho P=\rho$，定义反向转移矩阵
\[
 \widehat P(\ell,j)=\frac{\pi(j)P(j,\ell)}{\pi(\ell)}.
\]
不变性说明每一行和为一；不可约性在反向边上保持。令 $h(j)=\rho(j)/\pi(j)$，则
\[
 h(\ell)=\sum_j\widehat P(\ell,j)h(j).
\]
取 $h$ 的最大点 $\ell_0$。上式是若干不超过 $h(\ell_0)$ 的数的凸组合，所以每个满足
$\widehat P(\ell_0,j)>0$ 的 $j$ 都取同一最大值。沿反向链的正概率路径传播，不可约性迫使
$h$ 在全部状态上为常数。由 $\sum_j\rho(j)=\sum_j\pi(j)=1$，该常数为 $1$，故 $\rho=\pi$。
\end{proof}

非周期性可以转成所有状态共享的一块正质量。

\begin{lemma}[有限不可约非周期矩阵的共同正幂]\label{lemma:6-3-2-primitive}
若有限状态转移矩阵 $P$ 不可约且非周期，则存在 $m\geq1$，使
$P^m(i,j)>0$ 对所有 $i,j$ 成立。
\end{lemma}

\begin{proof}
固定状态 $r$。集合 $A=\{n\geq1:P^n(r,r)>0\}$ 对加法封闭，且其最大公因数为 $1$。
可从 $A$ 中选有限多个 $a_1,\ldots,a_q$，其最大公因数仍为 $1$。初等整数论说明由这些
$a_k$ 的非负整数线性组合组成的半群包含所有充分大的整数。具体地，其余生成元的剩余类在
$\Z/a_1\Z$ 中生成整个有限群；有限群中一个生成元剩余类的负元是该生成元剩余类的某个非负整数倍，
所以每个剩余类 $c$ 都可选到一个由 $a_2,\ldots,a_q$ 的非负组合表示的整数 $b_c\ge0$。若
$n\equiv c\pmod{a_1}$ 且 $n\ge\max_cb_c$，则
$n=b_c+\ell a_1$，其中 $\ell\ge0$。这就证明所有充分大的整数都属于该半群。

对每个 $i,j$，由不可约性选择正概率路径 $i\to r$ 和 $r\to j$，长度分别记为
$\alpha_i,\beta_j$。当 $n-\alpha_i-\beta_j$ 足够大时，可在 $r$ 处插入若干回路，使总长度恰为
$n$，因而 $P^n(i,j)>0$。状态对只有有限多个，取一个同时充分大的 $m$ 即得结论。
\end{proof}

\begin{theorem}[有限遍历收敛]\label{theorem:6-3-2-2}
若有限状态 Markov 链不可约且非周期，$\pi$ 为其不变分布，则对每个初态 $i$，
\[
 d_{\mathrm{TV}}(P^n(i,\cdot),\pi)\longrightarrow0.
\]
收敛事实上具有几何速度。
\end{theorem}

\begin{proof}
取引理~\ref{lemma:6-3-2-primitive} 中的 $m$，令
\[
 a(j)=\min_iP^m(i,j),\qquad
 \eta=\sum_ja(j)>0,
 \qquad \nu(j)=a(j)/\eta.
\]
若 $\eta=1$，所有 $P^m(i,\cdot)$ 都等于 $\nu$，结论直接成立。以下设 $0<\eta<1$，并定义
\[
 R(i,j)=\frac{P^m(i,j)-\eta\nu(j)}{1-\eta}.
\]
$R$ 是转移矩阵，且 $P^m(i,\cdot)=\eta\nu+(1-\eta)R(i,\cdot)$。因此对任意概率
$\mu,\rho$，Markov 核的全变差收缩性给
\[
 d_{\mathrm{TV}}(\mu P^m,\rho P^m)
 =(1-\eta)d_{\mathrm{TV}}(\mu R,\rho R)
 \leq(1-\eta)d_{\mathrm{TV}}(\mu,\rho).
\]
这里最后一个不等式可直接从
$\sup_A|\int R(x,A)(\mu-\rho)(\dd x)|$ 及有符号测度的正负分解得到。
迭代后
\[
 d_{\mathrm{TV}}(\delta_iP^{qm},\pi)
 \leq(1-\eta)^q d_{\mathrm{TV}}(\delta_i,\pi),
\]
因为 $\pi P^{qm}=\pi$。若 $n=qm+r$、$0\leq r<m$，再作用 $P^r$ 只会缩小全变差，故结论成立。
\end{proof}

单时刻收敛用到了非周期性；时间平均不需要它。回返时刻把一条相关路径切成独立周期。

\begin{theorem}[时间平均遍历定理]\label{theorem:6-3-2-3}
若 $(X_n)$ 是有限状态不可约 Markov 链，则对任意函数 $f:E\to\R$ 及任意初始分布，
\[
 \frac1n\sum_{k=0}^{n-1}f(X_k)\longrightarrow\sum_{x\in E}f(x)\pi(x)
 \qquad\text{a.s.}
\]
其中 $\pi$ 是定理~\ref{theorem:6-3-2-1} 的唯一不变分布。
\end{theorem}

\begin{proof}
固定状态 $i$。从任意初态出发，式~\eqref{eq:6-3-2-hitting-geometric-tail} 保证首次命中 $i$
的时刻 $\tau_0$ 有限 a.s.。递归定义
\[
 \tau_{r+1}=\inf\{n>\tau_r:X_n=i\},\qquad
 L_r=\tau_{r+1}-\tau_r,
 \qquad
 Y_r=\sum_{k=\tau_r}^{\tau_{r+1}-1}f(X_k).
\]
离散强 Markov 性说明 $(L_r,Y_r)_{r\geq0}$ 独立同分布，其分布等于从 $i$ 出发的一个回返
周期。这一点可以逐项验证：对 $\ell\ge1$，事件
\[
 \{L_r=\ell,\ (X_{\tau_r},\ldots,X_{\tau_r+\ell})\in A\}
\]
是停时后有限柱事件；强 Markov 性说明它在给定 $\mathcal F_{\tau_r}$ 后的概率只等于从
$i$ 出发的相应概率。对 $\ell$ 取可数并，再对 $r$ 归纳，便得整个周期块两两独立且同分布；
$Y_r$ 是该有限块的可测函数，所以结论传递到 $(L_r,Y_r)$。几何尾界给
$\E L_0<\infty$；状态有限使 $f$ 有界，故
$|Y_0|\leq\|f\|_\infty L_0$ 也可积。强大数律给
\[
 \frac{\sum_{r=0}^{q-1}Y_r}{\sum_{r=0}^{q-1}L_r}\longrightarrow
 \frac{\E_i\sum_{k=0}^{\tau_i^+-1}f(X_k)}{\E_i\tau_i^+}
 =\sum_xf(x)\pi(x).
\]

需要控制首尾不完整周期。首块长度 $\tau_0$ 是有限随机数，除以 $n$ 后趋零。对尾块，
可积性和尾和公式给每个 $\delta>0$
\[
 \sum_{r=1}^\infty\Prob(L_r>\delta r)
 =\sum_{r=1}^\infty\Prob(L_0>\delta r)
 \le \frac1\delta\int_0^\infty\Prob(L_0>u)\dd u+1
 =\frac{\E L_0}{\delta}+1<\infty,
\]
所以 Borel--Cantelli 说明 $L_r/r\to0$ a.s.。若
$N(n)=\max\{r:\tau_r\leq n\}$，则周期长的强大数律给
$\tau_r/r\to\E L_0$。几何尾界还说明每个回返时刻有限，所以 $N(n)\to\infty$。置
$r=N(n)$，由 $\tau_r\le n<\tau_{r+1}$ 得
\[
 \frac{r}{\tau_r}\ge\frac r n,
 \qquad
 \frac{r+1}{n}>\frac{r+1}{\tau_{r+1}}.
\]
两端分别趋于 $1/\E L_0$，而 $(r+1)/n-r/n=1/n\to0$；夹逼给
$N(n)/n\to1/\E L_0$。因而
$L_{N(n)}/n=(L_{N(n)}/N(n))(N(n)/n)\to0$。尾块奖励的绝对值不超过
$\|f\|_\infty L_{N(n)}$。更明确地，对充分大的 $n$，
\[
 \sum_{k=0}^{n-1}f(X_k)
 =\sum_{k=0}^{\tau_0-1}f(X_k)
  +\sum_{r=0}^{N(n)-1}Y_r
  +\sum_{k=\tau_{N(n)}}^{n-1}f(X_k).
\]
首项除以 $n$ 趋零，末项绝对值除以 $n$ 至多
$\|f\|_\infty L_{N(n)}/n\to0$。又
\[
 \frac{\tau_{N(n)}-\tau_0}{n}
 =1-\frac{n-\tau_{N(n)}+\tau_0}{n}\longrightarrow1,
\]
因为 $0\le n-\tau_{N(n)}<L_{N(n)}$。把完整周期的奖励和写成
\[
 \frac1n\sum_{r=0}^{N(n)-1}Y_r
 =\frac{\sum_{r=0}^{N(n)-1}Y_r}{\sum_{r=0}^{N(n)-1}L_r}
  \frac{\tau_{N(n)}-\tau_0}{n},
\]
前一因子已由强大数律确定，故全时间平均收敛到 $\sum_xf(x)\pi(x)$。
\end{proof}

\begin{proposition}[周期链的 \Cesaro{} 分布收敛]\label{proposition:6-3-2-5}
有限不可约链无论周期为何，都有
\[
 \frac1n\sum_{k=0}^{n-1}P^k(i,\cdot)\longrightarrow\pi
 \quad\text{于全变差}.
\]
\end{proposition}

\begin{proof}
对每个状态 $j$，把遍历定理用于 $f=\1_{\{j\}}$，再用有界收敛定理取期望，得到
\[
 \frac1n\sum_{k=0}^{n-1}P^k(i,j)\longrightarrow\pi(j).
\]
状态集有限，逐坐标收敛可在有限和中取极限；由
$d_{\mathrm{TV}}(\mu,\nu)=\tfrac12\sum_j|\mu(j)-\nu(j)|$ 即得全变差收敛。
二状态确定性交替链说明原来的 $P^n(i,\cdot)$ 可以不收敛。
\end{proof}

\begin{proposition}[可逆算子的 $L^2$ 结构]\label{proposition:6-3-2-4}
若 $P$ 相对于 $\pi$ 可逆，则它在实或复 Hilbert 空间 $L^2(\pi)$ 上自伴且为收缩：
\[
 \langle f,Pg\rangle_\pi=\langle Pf,g\rangle_\pi,
 \qquad \|Pf\|_{L^2(\pi)}\leq\|f\|_{L^2(\pi)}.
\]
\end{proposition}

\begin{proof}
采用 $\langle f,g\rangle_\pi=\int\overline f g\,\dd\pi$ 的复内积约定。对有界 $f,g$，
细致平衡恒等式与 Fubini 定理给出
\begin{align*}
 \langle f,Pg\rangle_\pi
 &=\int\!\int \overline{f(x)}g(y)P(x,\dd y)\pi(\dd x)\\
 &=\int\!\int \overline{f(x)}g(y)P(y,\dd x)\pi(\dd y)
 =\langle Pf,g\rangle_\pi.
\end{align*}
截断后令截断高度趋于无穷，可推广到 $L^2$。另一方面，Jensen 不等式和不变性给
\[
 \|Pf\|_2^2
 \leq\int P(|f|^2)(x)\pi(\dd x)
 =\int|f|^2\dd\pi.
\]
\end{proof}

自伴性为谱方法打开了入口，却不能单独消除周期。二状态确定性交替链在零均值子空间上的特征值为
$-1$；只控制“第二大特征值”而忽略负端，不能推出 $P^n$ 收敛。惰性化把它改成 $0$，或者必须控制
零均值子空间上的绝对谱半径。一般谱隙理论留待后文，这里不以它重新证明有限遍历结论。

\begin{remark}[MCMC 中的三种长时间问题]
在有限状态 MCMC 中，通常先设计核 $P$ 使目标分布 $\pi$ 满足细致平衡，由
定义~\ref{def:6-3-2-2} 得到 $\pi P=\pi$。这只解决“目标分布是否不变”。若还有不可约性与
非周期性，定理~\ref{theorem:6-3-2-2} 才给出从固定初态出发的全变差收敛，这是讨论预热期（burn-in）的
分布依据。无论周期如何，定理~\ref{theorem:6-3-2-3} 仍保证单条轨道上的样本平均收敛。
因此，“不变”、“预热后近似目标分布”和“时间平均可用”是三个需要不同结论支撑的陈述。
\end{remark}

\begin{figure}[htbp]
\centering
\begin{tikzpicture}[node distance=9mm and 12mm]
 \node[book strong node] (i0) {回返 $i$\ $\tau_0$};
 \node[book node,right=of i0] (b1) {周期 1\\长度 $L_0$\\奖励 $Y_0$};
 \node[book strong node,right=of b1] (i1) {回返 $i$\ $\tau_1$};
 \node[book node,right=of i1] (b2) {周期 2\\长度 $L_1$\\奖励 $Y_1$};
 \node[book strong node,right=of b2] (i2) {回返 $i$\ $\tau_2$};
 \draw[book accent arrow] (i0)--(b1);
 \draw[book accent arrow] (b1)--(i1);
 \draw[book accent arrow] (i1)--(b2);
 \draw[book accent arrow] (b2)--(i2);
 \node[book warm node,below=15mm of i1] (law) {强 Markov 性 $\Longrightarrow (L_r,Y_r)$ i.i.d.};
 \draw[book dashed arrow] (b1)--(law);
 \draw[book dashed arrow] (b2)--(law);
\end{tikzpicture}
\caption{回返把相关的 Markov 路径切成独立同分布的再生块；时间平均成为奖励和与长度和的比。}
\label{fig:6-3-2-regeneration}
\end{figure}

\begin{exercise}
\begin{enumerate}
 \item 求二状态链的 $P^n$，并分别验证周期情形的单时刻失败和 \Cesaro{} 收敛。
 \item 证明细致平衡蕴含不变性，并给出一个不变但不可逆的有限链。
 \item 对有限无向连通图核验度分布；判断三角形、偶环和加入自环后的周期。
 \item 在定理~\ref{theorem:6-3-2-3} 的证明中，逐步证明
 $N(n)/n\to1/\E_i\tau_i^+$。
\end{enumerate}
\end{exercise}

有限链的长期结构由回返组织。连续时间最基本的回返模型从随机等待开始：若等待时间没有记忆，
计数过程就会产生独立平稳增量。

\subsection{Poisson 过程、跳链与生成元}\label{subsec:6-3-3}
\index{Poisson 过程}\index{指数等待时间}\index{补偿鞅}\index{$Q$-矩阵}

\begin{example}[稀有 Bernoulli 极限]\label{ex:6-3-3-1}
把长度为 $t$ 的时间段切成 $n$ 个小格，并假设每格以概率 $\lambda t/n$ 发生一次事件，格间相互
独立。总事件数 $B_n$ 服从 $\operatorname{Bin}(n,\lambda t/n)$，所以对固定 $k$，
\[
 \Prob(B_n=k)
 =\binom nk\left(\frac{\lambda t}{n}\right)^k
  \left(1-\frac{\lambda t}{n}\right)^{n-k}
 \longrightarrow e^{-\lambda t}\frac{(\lambda t)^k}{k!}.
\]
这个计算给出了 Poisson 边缘，却没有说明不同时间段的计数怎样耦合。过程的独立增量、等待时间和
路径阶梯结构还必须另外建立。
\end{example}

\begin{definition}[Poisson 过程]\label{def:6-3-3-1}
速率 $\lambda>0$ 的 \emph{Poisson 过程}是过程 $(N_t)_{t\geq0}$，满足：
\begin{enumerate}
 \item $N_0=0$ a.s.，路径右连续有左极限、非减、取非负整数值且每次跳跃大小为 $1$；
 \item 不相交时间区间上的增量独立，且增量平稳；
 \item 对 $0\leq s<t$，
 \[
   N_t-N_s\sim\operatorname{Poisson}(\lambda(t-s)).
 \]
\end{enumerate}
\end{definition}

若先从这些有限维分布在典范乘积空间上扩张，随后构造的阶梯路径过程只保证与之有相同的过程分布；
除非两者已置于同一概率空间并逐时刻比较，不能把“同分布的正则实现”称为修正。

\begin{definition}[到达时刻与等待时间]\label{def:6-3-3-2}
对单位跳计数过程，令
\[
 T_0=0,\qquad T_n=\inf\{t\geq0:N_t\geq n\},\qquad
 W_n=T_n-T_{n-1}.
\]
$T_n$ 是第 $n$ 次到达时刻，$W_n$ 是相邻到达间隔。若所有 $T_n$ 有限且趋于无穷，则
\[
 N_t=\max\{n\geq0:T_n\leq t\}.
\]
\end{definition}

等待时间的无记忆性是 Markov 性在单个时钟上的原型。

\begin{proposition}[指数分布的无记忆性]\label{proposition:6-3-3-1}
若 $W\sim\operatorname{Exp}(\lambda)$，即
$\Prob(W>t)=e^{-\lambda t}$（$t\geq0$），则
\[
 \Prob(W>s+t\mid W>s)=\Prob(W>t).
\]
反过来，设 $W$ 取值于扩展半轴 $[0,\infty]$，生存函数 $G(t)=\Prob(W>t)$ 连续，并满足
$G(s+t)=G(s)G(t)$。若 $W$ 既不 a.s. 等于 $0$，也不 a.s. 等于 $+\infty$，则存在
$\lambda>0$ 使 $G(t)=e^{-\lambda t}$。
\end{proposition}

\begin{proof}
指数分布的正向计算为
\[
 \frac{\Prob(W>s+t)}{\Prob(W>s)}
 =\frac{e^{-\lambda(s+t)}}{e^{-\lambda s}}=e^{-\lambda t}.
\]
反向中，乘法式在 $s=t=0$ 时给 $G(0)=G(0)^2$。若 $G(0)=0$，则
$\Prob(W>0)=0$，即 $W=0$ a.s.，这是被排除的退化情形；故 $G(0)=1$。非退化性与连续性还
排除某个有限时刻取零：若 $G(t_0)=0$，则
$G(t_0/n)^n=0$，所以 $G(t_0/n)=0$ 对每个 $n$，与 $G(0)=1$ 处连续矛盾。因此
$H(t)=-\log G(t)$ 定义良好、连续、非降并满足 $H(s+t)=H(s)+H(t)$。对有理
$q\geq0$，加法性给 $H(q)=qH(1)$；连续性把它推广到全部 $t\geq0$。令
$\lambda=H(1)$。若 $\lambda=0$，则 $G\equiv1$，对应 $W=+\infty$ a.s.，已被排除；故
$\lambda>0$。
\end{proof}

\begin{theorem}[由等待时间构造 Poisson 过程]\label{theorem:6-3-3-2}
设 $W_1,W_2,\ldots$ 独立同分布为 $\operatorname{Exp}(\lambda)$，令
\[
 T_n=\sum_{k=1}^nW_k,\qquad
 N_t=\max\{n\geq0:T_n\leq t\}.
\]
则 $T_n\to\infty$ a.s.；在这个概率一事件上，上式对每个有限 $t$ 都定义一个有限整数，
并在其补集上统一取零路径。这样得到的阶梯过程是速率
$\lambda$ 的 Poisson 过程。
\end{theorem}

\begin{proof}
事件 $A_n=\{W_n\geq1\}$ 独立且概率为 $e^{-\lambda}>0$，所以
$\sum_n\Prob(A_n)=\infty$。Borel--Cantelli II 说明 $A_n$ 发生无穷多次，故
$T_n\to\infty$ a.s.。把例外零集上的整条路径统一定义为零以后，每个 $N_t$ 都处处有定义，
而且后续全部路径性质可在同一个概率一事件上同时核验。

卷积归纳给 $T_n$ 的 Gamma 密度
\begin{equation}\label{eq:6-3-3-gamma-density}
 f_{T_n}(u)=\frac{\lambda^n u^{n-1}}{(n-1)!}e^{-\lambda u}\1_{(0,\infty)}(u).
\end{equation}
确切地说，$n=1$ 是指数密度；若公式对 $n$ 成立，则
\begin{align*}
 f_{T_{n+1}}(u)
 &=\int_0^u\frac{\lambda^nv^{n-1}}{(n-1)!}e^{-\lambda v}
     \lambda e^{-\lambda(u-v)}\dd v\\
 &=\frac{\lambda^{n+1}u^n}{n!}e^{-\lambda u}.
\end{align*}
于是对 $n\geq0$，
\begin{align*}
 \Prob(N_t=n)
 &=\Prob(T_n\leq t<T_{n+1})\\
 &=\int_0^tf_{T_n}(u)\Prob(W_{n+1}>t-u)\dd u
 =e^{-\lambda t}\frac{(\lambda t)^n}{n!},
\end{align*}
其中 $n=0$ 直接由 $\Prob(W_1>t)$ 计算。

还要证明增量独立。固定确定时刻 $s$，记
$\mathcal F_s^N=\sigma(N_u:0\le u\le s)$。在 $\{N_s=n\}$ 上，$s$ 以前的整条计数路径由
$(T_1,\ldots,T_n)$ 决定；反过来，每个 $T_j\wedge s$ 都是该路径的可测函数。因此
$\mathcal F_s^N$ 在 $\{N_s=n\}$ 上的迹由事件
\[
 \{N_s=n\}\cap\{(T_1,\ldots,T_n)\in C\},
 \qquad C\in\Borel(\R^n),
\]
生成（$n=0$ 时只剩 $\{N_s=0\}$）。

在 $\{N_s=n\}$ 上置剩余等待时间 $R_s=T_{n+1}-s$。设 $r\ge0$，
$C\in\Borel(\R^n)$，并取 Borel 集 $A_2,\ldots,A_m\subset(0,\infty)$。对
$(W_1,\ldots,W_n)$ 条件化，并使用后续等待时间的独立性，得到
\begin{align*}
&\Prob\bigl((T_1,\ldots,T_n)\in C,\ N_s=n,\ R_s>r,
             \ W_{n+j}\in A_j\ (2\le j\le m)\bigr)\\
&=\E\!\left[
 \1_{\{(T_1,\ldots,T_n)\in C,\ T_n\le s\}}
 \Prob(W_{n+1}>s-T_n+r\mid W_1,\ldots,W_n)
 \right]
 \prod_{j=2}^m\Prob(W_1\in A_j)\\
&=e^{-\lambda r}
 \E\!\left[
 \1_{\{(T_1,\ldots,T_n)\in C,\ T_n\le s\}}
 e^{-\lambda(s-T_n)}\right]
 \prod_{j=2}^m\Prob(W_1\in A_j)\\
&=e^{-\lambda r}\prod_{j=2}^m\Prob(W_1\in A_j)\,
 \Prob\bigl((T_1,\ldots,T_n)\in C,\ N_s=n\bigr).
\end{align*}
第一行中的 $N_s=n$ 等价于 $T_n\le s<T_{n+1}$；最后一行使用
$\Prob(W_{n+1}>s-T_n\mid W_1,\ldots,W_n)=e^{-\lambda(s-T_n)}$。
对矩形 $\prod_{j=2}^mA_j$ 先用有限交，再用 $\pi$--$\lambda$ 定理，以上分解说明：条件于
$\mathcal F_s^N$，$R_s,W_{N_s+2},W_{N_s+3},\ldots$ 相互独立、都服从
$\operatorname{Exp}(\lambda)$，而且其联合条件分布不依赖过去。于是从 $s$ 起的到达过程由一列
新的 i.i.d. 指数等待时间构造，并与 $\mathcal F_s^N$ 独立。已经证明的边缘计算
应用于这列新时钟，得到
$N_{s+t}-N_s\sim\operatorname{Poisson}(\lambda t)$，并由递归得到任意有限个不相交区间的
增量独立。阶梯路径的右连续、左极限、非减和单位跳性质由定义直接读出。
\end{proof}

反方向不能只由第一到达时间的尾公式猜测全部等待时间；必须同时读出不同到达之间的联合结构。

\begin{proposition}[Poisson 增量反推出指数等待]\label{proposition:6-3-3-arrivals}
对定义~\ref{def:6-3-3-1} 的 Poisson 过程，等待时间
$W_1,W_2,\ldots$ 独立同分布为 $\operatorname{Exp}(\lambda)$。更精确地，
$(T_1,\ldots,T_n)$ 在区域 $0<t_1<\cdots<t_n$ 上具有联合密度
\begin{equation}\label{eq:6-3-3-arrival-joint-density}
 \lambda^ne^{-\lambda t_n}\1_{\{0<t_1<\cdots<t_n\}}.
\end{equation}
\end{proposition}

\begin{proof}
先固定 $t>0$ 和分割 $0=s_0<s_1<\cdots<s_m=t$。独立增量给出，在
$\sum_jn_j=n$ 时
\begin{align*}
 &\Prob\bigl(N_{s_j}-N_{s_{j-1}}=n_j,\ 1\leq j\leq m\mid N_t=n\bigr)\\
 &\qquad=
 \frac{n!}{\prod_jn_j!}\prod_{j=1}^m
 \left(\frac{s_j-s_{j-1}}{t}\right)^{n_j}.
\end{align*}
这正是 $n$ 个独立 $\operatorname{Unif}(0,t)$ 点落入各小区间的多项分布。半开区间组成的有限
分割生成 $[0,t]^n$ 的 Borel $\sigma$-代数，因此条件于 $N_t=n$，到达点的无序配置与
$n$ 个独立均匀点相同；排序后，
\[
 (T_1,\ldots,T_n)\mid\{N_t=n\}
\]
在 $0<t_1<\cdots<t_n<t$ 上具有密度 $n!/t^n$。

为严格得到无条件密度，固定 $T>0$。条件于 $N_T=m\ge n$ 时，全部 $m$ 个到达点是 $m$ 个
$\operatorname{Unif}(0,T)$ 点的次序统计量。把后 $m-n$ 个有序坐标在
$t_n<u_{n+1}<\cdots<u_m<T$ 上积分，得到前 $n$ 个到达时刻的条件密度
\[
 \frac{m!}{(m-n)!T^m}(T-t_n)^{m-n}
 \1_{\{0<t_1<\cdots<t_n<T\}}.
\]
乘以 $\Prob(N_T=m)=e^{-\lambda T}(\lambda T)^m/m!$ 并对 $m\ge n$ 求和，Tonelli 定理给出
\begin{align*}
 &\sum_{m=n}^\infty
 e^{-\lambda T}\frac{(\lambda T)^m}{m!}
 \frac{m!}{(m-n)!T^m}(T-t_n)^{m-n}\\
 &\qquad=e^{-\lambda T}\lambda^n
 \sum_{r=0}^\infty\frac{(\lambda(T-t_n))^r}{r!}
 =\lambda^ne^{-\lambda t_n}.
\end{align*}
左边是 $(T_1,\ldots,T_n)$ 在事件 $\{T_n<T\}$ 上的密度。任给
$0<t_1<\cdots<t_n$ 都可选择 $T>t_n$，故这就证明了整个正单纯形上的无条件密度
\eqref{eq:6-3-3-arrival-joint-density}。作三角变量替换
\[
 w_1=t_1,\qquad w_k=t_k-t_{k-1}\quad(k\geq2),
\]
雅可比行列式为 $1$，且 $t_n=\sum_{k=1}^nw_k$。故 $(W_1,\ldots,W_n)$ 的密度为
\[
 \lambda^ne^{-\lambda\sum_kw_k}\prod_{k=1}^n\1_{(0,\infty)}(w_k)
 =\prod_{k=1}^n\lambda e^{-\lambda w_k}\1_{(0,\infty)}(w_k).
\]
这对每个 $n$ 都成立，因而整列等待时间独立同分布为指数分布。特别地，$n=1$ 时退化为
$\Prob(T_1>t)=\Prob(N_t=0)=e^{-\lambda t}$。
\end{proof}

只有边缘正确仍可能构造出错误的过程分布。

\begin{example}[正确边缘、错误增量]\label{ex:6-3-3-4}
取一个 $U\sim\operatorname{Unif}(0,1)$，令 $\widetilde N_t=F_t^{-1}(U)$，其中 $F_t$ 是
$\operatorname{Poisson}(\lambda t)$ 的分布函数，分位数取左连续约定。每个
$\widetilde N_t$ 都有正确边缘，并且可取成随 $t$ 非减；可是全部时刻由同一个 $U$ 驱动。

例如取 $\lambda=1,s=1,t=2$。事件
$\{\widetilde N_1=0,\widetilde N_2=1\}$ 对应
\[
 e^{-2}<U\leq\min\{e^{-1},3e^{-2}\}=e^{-1},
\]
所以其概率为 $e^{-1}-e^{-2}$。真正的 Poisson 过程中，独立增量给
\[
 \Prob(N_1=0,N_2=1)=\Prob(N_1=0)\Prob(N_2-N_1=1)=e^{-2}.
\]
两数不同。这一比较直接说明一维 Poisson 公式不能替代独立增量。
\end{example}

中心化计数揭示线性漂移背后的鞅。

\begin{theorem}[Poisson 补偿鞅]\label{theorem:6-3-3-3}
设 $N$ 是速率 $\lambda$ 的 Poisson 过程，$M_t=N_t-\lambda t$。相对于原始自然滤过
$\mathcal F_t^N=\sigma(N_s:0\le s\le t)$ 或它的完备化，以下三个过程都是平方可积鞅：
\[
 M_t,\qquad M_t^2-\lambda t,\qquad M_t^2-N_t.
\]
\end{theorem}

\begin{proof}
固定 $0\leq s<t$，置 $Z=N_t-N_s$。它独立于 $\mathcal F_s$，且
$\E Z=\Var(Z)=\lambda(t-s)$。因此
\[
 \E[M_t-M_s\mid\mathcal F_s]=\E[Z-\lambda(t-s)]=0,
\]
证明 $M$ 为鞅。再令 $Y=Z-\lambda(t-s)$；$Y$ 独立于过去、均值为零、二阶矩为
$\lambda(t-s)$，所以
\[
 \E[M_t^2\mid\mathcal F_s]
 =\E[(M_s+Y)^2\mid\mathcal F_s]
 =M_s^2+\lambda(t-s).
\]
这给第二个鞅。最后
\[
 \E[M_t^2-N_t\mid\mathcal F_s]
 =M_s^2+\lambda(t-s)-N_s-\lambda(t-s)
 =M_s^2-N_s.
\]
Poisson 变量具有全部有限阶矩：由其概率质量函数，任意整数 $r\ge1$ 都有
\[
 \E N_t^r=e^{-\lambda t}\sum_{k=0}^\infty
 k^r\frac{(\lambda t)^k}{k!}<\infty,
\]
因为相邻项之比趋于零。特别地，前两个含 $M_t^2$ 的过程也属于 $L^2$，而上面的条件期望计算
均合法。把滤过完备化只加入零集及其子集；条件期望等式因此仍成立。本结论不需要预先断言原始
自然滤过已经右连续。
\end{proof}

$\lambda t$ 是由过去可以预测的平方波动累计量，而 $N_t$ 是实际跳跃平方和；二者之差仍是鞅。
这组公式将在 Brownian 二次变差出现后得到连续路径对照。

\begin{example}[复合 Poisson 过程]\label{ex:6-3-3-2}
令 $Y_1,Y_2,\ldots$ 为独立同分布的 $\R^d$ 值随机向量，共同分布为 $\nu$，并与 $N$ 独立。置
\[
 X_t=\sum_{k=1}^{N_t}Y_k,
\]
空和取零。若 $\varphi_Y(u)=\E e^{i\langle u,Y_1\rangle}$，则对 $u\in\R^d$，
\begin{align*}
 \E e^{i\langle u,X_t\rangle}
 &=\sum_{n=0}^\infty\Prob(N_t=n)\varphi_Y(u)^n\\
 &=e^{-\lambda t}\sum_{n=0}^\infty\frac{(\lambda t\varphi_Y(u))^n}{n!}
 =\exp\{\lambda t(\varphi_Y(u)-1)\}.
\end{align*}
这里 $e^{i\langle u,x\rangle}$ 只是正弦、余弦两项的合写，计算不使用复分析。
\end{example}

\begin{proposition}[复合 Poisson 过程的生成元]\label{proposition:6-3-3-compound-generator}
对从 $x$ 出发的复合 Poisson 过程，
\[
 P_tf(x)=\E f\left(x+\sum_{k=1}^{N_t}Y_k\right).
\]
对有界可测 $f$，其点态生成元为
\[
 Lf(x)=\lambda\int_{\R^d}\bigl(f(x+y)-f(x)\bigr)\nu(\dd y).
\]
若 $f\in C_0(\R^d)$，则 $P_tC_0\subset C_0$、
$\|P_tf-f\|_\infty\to0$，并且上述差商在 $C_0$ 范数中收敛。
\end{proposition}

\begin{proof}
当 $h\downarrow0$ 时，
\[
 \Prob(N_h=0)=1-\lambda h+O(h^2),\quad
 \Prob(N_h=1)=\lambda h+O(h^2),\quad
 \Prob(N_h\geq2)=O(h^2).
\]
按这三个事件分割，并用 $|f|\leq\|f\|_\infty$ 控制最后一项，得到一致于 $x$ 的展开
\[
 P_hf(x)=f(x)+\lambda h\int(f(x+y)-f(x))\nu(\dd y)+O(h^2)\|f\|_\infty.
\]
这给生成元公式。若 $f\in C_0$，则
$x\mapsto\int f(x+y)\nu(\dd y)$ 连续：对 $x_n\to x$ 用 DCT；它也在无穷远消失：先取紧集
$K$ 使 $\nu(K)>1-\varepsilon$，在 $K$ 上利用 $f(x+y)$ 的一致尾小性，在 $K^c$ 上用
$\|f\|_\infty\varepsilon$ 控制。Poisson 级数遂说明 $P_tf\in C_0$，而上面的统一展开给强连续性
与范数生成元结论。取 $f(x)=e^{iu\cdot x}$ 时，
\[
 Lf(x)=\lambda(\varphi_Y(u)-1)e^{iu\cdot x},
\]
与前一个特征函数计算相符。
\end{proof}

\begin{definition}[$Q$-矩阵]\label{def:6-3-3-3}
在可数状态空间 $E$ 上，数族 $(q_{ij})$ 称为保守 $Q$-矩阵，若
\[
 q_{ij}\geq0\ (i\ne j),\qquad
 q_i:=\sum_{j\ne i}q_{ij}<\infty,\qquad q_{ii}=-q_i.
\]
其形式生成元为
\[
 Qf(i)=\sum_jq_{ij}f(j)
 =\sum_{j\ne i}q_{ij}(f(j)-f(i)).
\]
在状态 $i$ 停留 $\operatorname{Exp}(q_i)$ 时间后，以概率 $q_{ij}/q_i$ 跳到 $j$；若 $q_i=0$，
约定 $i$ 为吸收态。
\end{definition}

\begin{example}[有限状态跳链]\label{ex:6-3-3-3}
若 $E$ 有限，取 $\Lambda>0$ 且 $\Lambda\geq\max_iq_i$，并置
\[
 R=I+\Lambda^{-1}Q.
\]
$R$ 是转移矩阵。取速率 $\Lambda$ 的 Poisson 时钟，并在每次时钟响起时按 $R$ 更新状态；
$R(i,i)=1-q_i/\Lambda$ 对应一次虚跳。条件于时钟次数，转移矩阵为
\[
 P_t=e^{-\Lambda t}\sum_{n=0}^\infty\frac{(\Lambda t)^n}{n!}R^n
 =e^{-\Lambda t}e^{\Lambda tR}
 =e^{\Lambda t(R-I)}=e^{tQ}.
\]
这称为一致化（uniformization）。Poisson 时钟在有限时间内只响有限次，故构造不爆炸。
\end{example}

同一论证适用于 $\sup_iq_i<\infty$ 的可数状态链，此时 $Q$ 是 $\ell^\infty(E)$ 上的有界算子。
仅有每个 $q_i<\infty$ 不够：纯出生链 $q_{i,i+1}=2^i$ 的连续等待时间期望和为
$\sum_i2^{-i}<\infty$，而非负总等待时间的期望有限迫使其总和 a.s. 有限；链可在有限时间内完成
无穷多次跳跃。这就是爆炸，必须另加条件排除。

\begin{proposition}[Kolmogorov 前后向方程]\label{proposition:6-3-3-4}
在有限状态或一致有界率情形，一致化定义的半群在矩阵范数或
$\ell^\infty$ 算子范数中可微，并满足
\[
 \frac{\dd}{\dd t}P_t=QP_t=P_tQ,\qquad P_0=I.
\]
\end{proposition}

\begin{proof}
$Q$ 有界，算子幂级数
$P_t=e^{tQ}=\sum_{n\geq0}t^nQ^n/n!$ 在每个有界时间区间上一致绝对收敛；导数级数也一致绝对
收敛。因此可逐项求导，得到
\[
 P_t'=\sum_{n=1}^\infty\frac{nt^{n-1}Q^n}{n!}
 =Q\sum_{m=0}^\infty\frac{t^mQ^m}{m!}=QP_t.
\]
$Q$ 与自己的幂交换，故同样等于 $P_tQ$。在概率语言中，左乘 $Q$ 是从起点作局部展开，右乘
$Q$ 是在终点作局部展开；二者由半群律协调。
\end{proof}

\begin{remark}[Poisson 与复合 Poisson 模型的应用]
若把 $N_t$ 解释为到达数，定理~\ref{theorem:6-3-3-2} 给出具有指数间隔的事件流；在排队模型中
它描述顾客到达。若 $Y_k\ge0$ 是第 $k$ 件事故的赔款，则例~\ref{ex:6-3-3-2} 的
$X_t=\sum_{k\le N_t}Y_k$ 是保险风险的累计损失。这两种解释都直接调用已证的独立平稳增量，
而不是仅凭单时刻 Poisson 边缘。

复合 Poisson 过程是纯跳、独立平稳增量过程的基本原型。
定理~\ref{theorem:6-3-3-3} 的补偿鞅 $N_t-\lambda t$ 把可预测强度与实际计数分开，为计数过程的随机积分提供基本的中心化噪声。
\end{remark}

\begin{remark}[\cadlagPath{}的空间边界]
Poisson 与一般跳链的自然路径有跳点，所以不能把它们的分布放到连续路径空间
$C([0,T])$ 上。适合的容器是 \cadlagPathSpace{} $D([0,T])$，其上通常配 Skorokhod 拓扑。
对显式构造的阶梯路径，当前的有限维分布计算已经足够。若要研究一般 $D([0,T])$-值随机元的弱收敛，
还需明确 Skorokhod 拓扑、其 Borel $\sigma$-代数以及相应的紧性判据。
\end{remark}

\begin{figure}[htbp]
\centering
\begin{tikzpicture}[x=1cm,y=1cm]
 \draw[book line,-{Stealth[length=2.4mm]}] (0,0)--(9.4,0) node[right,book label] {$t$};
 \foreach \x/\k in {1.1/1,2.7/2,4.8/3,7.6/4}{
   \draw[BookAccent,semithick] (\x,-0.12)--(\x,0.9);
   \node[book label,below] at (\x,-0.12) {$T_{\k}$};
 }
 \draw[decorate,decoration={brace,amplitude=4pt},BookWarm]
   (1.1,1.05)--(2.7,1.05) node[midway,above=4pt,book label] {$W_2$};
 \draw[BookWarning,semithick] (0,-0.65)--(1.1,-0.65)--(1.1,-0.35)--(2.7,-0.35)--
   (2.7,-0.05)--(4.8,-0.05)--(4.8,0.25)--(7.6,0.25)--(7.6,0.55)--(9.1,0.55);
 \node[book label] at (8.6,0.95) {$N_t$};
 \node[book warm node] at (4.7,-1.45) {局部率 $Q$\ $\xrightarrow{\ e^{tQ}\ }$\ 转移半群 $P_t$};
\end{tikzpicture}
\caption{等待时间决定到达点，到达点累计成单位跳阶梯；一致有界跳率经一致化生成半群。}
\label{fig:6-3-3-arrivals-and-generator}
\end{figure}

\begin{exercise}
\begin{enumerate}
 \item 从密度卷积重新推导式~\eqref{eq:6-3-3-gamma-density} 与 Poisson 计数分布。
 \item 证明条件于 $N_t=n$ 的到达点是 $n$ 个均匀点的次序统计量。
 \item 直接核验 $N_t-\lambda t$、$(N_t-\lambda t)^2-N_t$ 的鞅等式。
 \item 对纯出生链例证明无穷等待时间之和 a.s. 有限，并说明这一步为什么不能用于一致有界率链。
\end{enumerate}
\end{exercise}

Poisson 过程把时间积累为跳跃平方和。下一节转向连续路径噪声：Brownian motion 没有跳跃，
却仍有非零二次变差；这正是随机积分不能按经典有限变差积分处理的原因。

\section{Brownian motion 与随机微积分基础}\label{sec:06-04}

\subsection{Gaussian 过程与 Brownian motion 的构造}\label{subsec:6-4-1}
\index{Gaussian 过程}\index{Brownian motion}\index{Wiener 测度}

把中心随机游走的时间加快 $n$ 倍、空间缩小 $\sqrt n$ 倍，折线路径的单步高度变成
$n^{-1/2}$，而单位时间内方差仍为一。这个缩放暗示一个连续极限，却没有提供存在性证明。
此处把缩放关系作为模型动机，而存在性问题将由有限维 Gaussian 分布和路径连续性判据直接解决。

Brown 对花粉微粒的观察、Bachelier 的价格模型和 Einstein 的扩散计算都早于严格的路径测度。
Wiener 从路径空间出发，Kolmogorov 则把相容有限维分布拼成过程。这里采用第二条路线：先由一个
协方差核构造全部有限维 Gaussian 分布，再用连续性定理选择连续版本。扩张与连续化是两个不同步骤；
只完成前一步，得到的仍只是所有函数构成的典范空间上的坐标过程。

\begin{definition}[Gaussian 过程]\label{def:6-4-1-1}
实值过程 $(X_t)_{t\in T}$ 称为 \emph{Gaussian 过程}，如果对任意
$n\ge1$、$t_1,\ldots,t_n\in T$ 和 $a_1,\ldots,a_n\in\R$，线性组合
$\sum_{j=1}^na_jX_{t_j}$ 都具有一维 Gaussian 分布；允许方差为零的退化 Gaussian 分布。
其均值函数与协方差核分别为
\[
 m(t)=\E X_t,
 \qquad K(s,t)=\Cov(X_s,X_t).
\]
\end{definition}

对有限时刻组，向量 $(X_{t_1},\ldots,X_{t_n})$ 的特征函数为
\begin{equation}\label{eq:6-4-1-gaussian-fdd-cf}
 \E e^{i\sum_ju_jX_{t_j}}
 =\exp\!\left(i\sum_ju_jm(t_j)
 -\frac12\sum_{j,k}u_ju_kK(t_j,t_k)\right).
\end{equation}
这是第五章 Gaussian 特征函数计算对线性组合的应用。特征函数唯一性定理
\ref{theorem:5-4-3-2} 因而说明：Gaussian 过程的均值和协方差唯一决定全部有限维分布。
这项识别只对联合 Gaussian 族成立；对一般过程，相同的前两阶矩远不足以确定分布。

\begin{example}[协方差核]\label{ex:6-4-1-1}
对 $s,t\ge0$ 令 $K(s,t)=s\wedge t$。给定 $t_1,\ldots,t_n\ge0$ 与
$a_1,\ldots,a_n\in\R$，Tonelli 定理给出
\begin{align*}
 \sum_{j,k=1}^na_ja_k(t_j\wedge t_k)
 &=\sum_{j,k}a_ja_k\int_0^\infty
   \1_{[0,t_j]}(u)\1_{[0,t_k]}(u)\dd u\\
 &=\int_0^\infty\left(\sum_{j=1}^na_j\1_{[0,t_j]}(u)\right)^2\dd u
 \ge0.
\end{align*}
所以每个有限矩阵 $(t_j\wedge t_k)_{j,k}$ 都半正定，可以作为 Gaussian 向量的协方差矩阵。
\end{example}

\begin{definition}[Brownian motion]\label{def:6-4-1-2}
实过程 $(B_t)_{t\ge0}$ 称为\emph{标准 Brownian motion}，如果：
\begin{enumerate}
  \item $B_0=0$ a.s.；
  \item 几乎每条路径 $t\mapsto B_t$ 连续；
  \item 对任意 $0\le t_0<t_1<\cdots<t_n$，增量
  $B_{t_k}-B_{t_{k-1}}$ 相互独立，且
  \[
    B_{t_k}-B_{t_{k-1}}\sim N(0,t_k-t_{k-1}).
  \]
\end{enumerate}
\end{definition}

\begin{definition}[Wiener 测度]\label{def:6-4-1-3}
令
\[
 C_0([0,\infty))=\{\omega\in C([0,\infty)):\omega(0)=0\},
\]
并赋予局部一致收敛拓扑。在这个路径空间上，使坐标过程
$e_t(\omega)=\omega(t)$ 成为标准 Brownian motion 的 Borel 概率测度称为 \emph{Wiener 测度}。
这里下标 $0$ 表示路径从零出发，不表示在无穷远消失。
\end{definition}

例如，下式给出该拓扑的一个完备可分度量：
\[
 d(\omega,\eta)=\sum_{m=1}^\infty2^{-m}
 \bigl(1\wedge\|\omega-\eta\|_{C[0,m]}\bigr).
\]
这个度量确实完备：若 $(\omega_n)$ 为 Cauchy 列，则对每个 $m$，限制
$\omega_n|_{[0,m]}$ 在 Banach 空间 $C[0,m]$ 中为 Cauchy 列；所得极限在重叠区间上一致，因而拼成
一条从零出发的连续路径。它也可分：先令 $m$ 遍历正整数，再取在 $[0,m]$ 上具有有理分点和有理
函数值的折线，并在 $m$ 以后作常值延拓；这些路径构成可数集，而局部一致逼近与度量尾项
$\sum_{k>m}2^{-k}$ 给出稠密性。因此 $C_0([0,\infty))$ 是 Polish 空间。对连续函数，上确界可以只在有理时刻取；逐个整数区间
使用命题~\ref{proposition:6-1-3-4} 的论证可得
\begin{equation}\label{eq:6-4-1-wiener-borel}
 \Borel(C_0([0,\infty)))
 =\sigma(e_q:q\in\Q_+).
\end{equation}
所以 Wiener 测度由 Brownian 的有限维分布唯一决定。

\begin{proposition}[联合 Gaussian 中不相关即独立]\label{proposition:6-4-1-2}
设 $Y=(Y_1,\ldots,Y_n)$ 是 Gaussian 向量。若把指标分成两块 $I,J$，并且
$\Cov(Y_i,Y_j)=0$ 对所有 $i\in I,j\in J$ 成立，则子向量 $Y_I,Y_J$ 独立。
特别地，均值零、协方差核为 $s\wedge t$ 的 Gaussian 过程具有独立平稳增量。
\end{proposition}

\begin{proof}
把均值和协方差矩阵按 $I,J$ 分块。交叉协方差为零使矩阵为块对角矩阵，故
式~\eqref{eq:6-4-1-gaussian-fdd-cf} 给出
\[
 \varphi_Y(u_I,u_J)
 =\varphi_{Y_I}(u_I)\varphi_{Y_J}(u_J).
\]
右端是边缘分布乘积测度的特征函数。特征函数唯一性定理说明联合分布等于边缘分布的乘积，
所以两个子向量独立。

现在设 $X$ 是中心 Gaussian 过程，协方差为 $s\wedge t$。对 $0\le a<b\le c<d$，
\begin{align*}
 &\Cov(X_b-X_a,X_d-X_c)\\
 &\quad=(b\wedge d)-(b\wedge c)-(a\wedge d)+(a\wedge c)
 =b-b-a+a=0.
\end{align*}
任意有限个两两不交区间增量仍是一个联合 Gaussian 向量；上面的块论证逐个分块，给出共同独立性。
此外
\[
 \Var(X_t-X_s)=t+s-2(s\wedge t)=|t-s|,
\]
故 $X_t-X_s\sim N(0,|t-s|)$。其分布只依赖时间差，因而增量平稳。
\end{proof}

\begin{theorem}[Brownian 存在性]\label{theorem:6-4-1-1}
存在标准 Brownian motion，因而存在唯一 Wiener 测度。所得过程对每个
$0<\gamma<1/2$ 都有局部 $\gamma$-\Holder{} 版本；只用四阶矩时已经足以得到连续版本，并可得到
$0<\gamma<1/4$。
\end{theorem}

\begin{proof}
第一步构造有限维分布。对每个有限集 $F=\{t_1,\ldots,t_n\}\subset[0,\infty)$，令
$\mu_F$ 为中心 Gaussian 分布，其协方差矩阵为 $(t_j\wedge t_k)_{j,k}$。例
\ref{ex:6-4-1-1} 保证该矩阵半正定，故这样的 Gaussian 分布存在。若 $F\subset G$，把
$u\in\R^F$ 在 $G\setminus F$ 上补零，式~\eqref{eq:6-4-1-gaussian-fdd-cf} 说明
$(\pi_F^G)_*\mu_G$ 与 $\mu_F$ 的特征函数相同。由唯一性定理，它们相等。因此
$(\mu_F)$ 投影相容。Kolmogorov 扩张定理~\ref{theorem:6-1-2-1} 给出典范空间上的坐标过程
$X=(X_t)_{t\ge0}$，其均值为零、协方差为 $s\wedge t$。

第二步选择连续版本。若 $Z\sim N(0,1)$，记
\[
 I_k=\frac1{\sqrt{2\pi}}\int_{\R}x^ke^{-x^2/2}\dd x\qquad(k=0,2,4).
\]
在 $[-R,R]$ 上用 $xe^{-x^2/2}=-(e^{-x^2/2})'$ 分部积分，再令 $R\to\infty$，边界项
$R^{k-1}e^{-R^2/2}$ 趋于零，故 $I_k=(k-1)I_{k-2}$。由 $I_0=1$ 得
$\E Z^4=I_4=3I_2=3$。因此
\[
 \E|X_t-X_s|^4=3|t-s|^2.
\]
在每个 $[0,m]$ 上把 Kolmogorov--Chentsov 定理~\ref{theorem:6-1-3-1} 用于
$\alpha=4$、参数维数 $1$、$1+\beta=2$，得到任意 $\gamma<1/4$ 的连续版本
$X^{(m)}$。当 $m<\ell$ 时，两个版本在 $[0,m]$ 上互为版本；连续版本唯一性命题
\ref{proposition:6-1-3-3} 说明它们在该区间不可分辨。先在每个整数区间的有理时刻协调，再对
可数个 $m$ 取零集之并，便可把 $X^{(m)}$ 拼成 $[0,\infty)$ 上的连续版本 $B$。

版本改变不改变有限维分布。命题~\ref{proposition:6-4-1-2} 因而说明 $B$ 有独立平稳增量，且
$B_t-B_s\sim N(0,t-s)$ 对 $s<t$ 成立；协方差在 $t=0$ 为零，故 $B_0=0$ a.s.。所以 $B$
满足定义~\ref{def:6-4-1-2}。

若希望在同一个概率为一的事件上同时得到每个 $\gamma<1/2$，先取偶整数 $p>2$。Gaussian 缩放给出
\[
 \E|X_t-X_s|^p=(\E|Z|^p)|t-s|^{p/2}.
\]
在连续性定理中取 $\alpha=p$、$1+\beta=p/2$，可得到任意
$\gamma<1/2-1/p$。对 $k\ge3$ 置 $\gamma_k=1/2-1/k$，并选偶整数 $p_k>k$；连续版本唯一性把所得
局部 $\gamma_k$-\Holder{} 版本与刚构造的 $B$ 识别。在可数个 $k$ 与整数时间区间上取共同零集后，
$B$ 在同一个概率为一的事件上对每个 $k$ 都局部 $\gamma_k$-\Holder{}。任给 $\gamma<1/2$，选
$k$ 使 $\gamma<\gamma_k$；较高指数在每个紧时间区间上蕴含较低指数，故 $B$ 同时具有所述全部局部
\Holder{} 正则性。端点 $1/2$ 不由这项估计得到。

最后把每条连续路径之外的共同零集改成零路径。式~\eqref{eq:6-4-1-wiener-borel} 说明路径映射
$\omega\mapsto B_\bullet(\omega)$ 可测，其分布就是 Wiener 测度。有限维分布唯一性与
式~\eqref{eq:6-4-1-wiener-borel} 又说明这样的路径空间概率至多一个。
\end{proof}

\begin{example}[Brownian bridge]\label{ex:6-4-1-2}
令 $B$ 为 Brownian motion，并在 $0\le t\le1$ 上定义
\[
 \beta_t=B_t-tB_1.
\]
每个有限向量 $(\beta_{t_1},\ldots,\beta_{t_n})$ 是 Brownian 有限向量的线性变换，故联合
Gaussian 且均值为零。对 $0\le s,t\le1$，
\begin{align*}
 \Cov(\beta_s,\beta_t)
 &=s\wedge t-t\Cov(B_s,B_1)-s\Cov(B_1,B_t)+st\Var(B_1)\\
 &=s\wedge t-st.
\end{align*}
并且 $\beta_0=\beta_1=0$。这个过程称为 Brownian bridge；它在两个端点被钉住，因而不再具有
Brownian 的独立增量。
\end{example}

\begin{example}[Brownian 缩放]\label{ex:6-4-1-3}
固定 $c>0$，令 $Y_t=c^{-1/2}B_{ct}$。过程 $Y$ 连续、均值为零，并且
\[
 \Cov(Y_s,Y_t)=c^{-1}(cs\wedge ct)=s\wedge t.
\]
$Y$ 与 $B$ 都是 Gaussian 过程，所以均值和协方差相同蕴含全部有限维分布相同。二者又都取值于
$C_0([0,\infty))$，而该空间的 Borel $\sigma$-代数由有理评价生成；故
\[
 (c^{-1/2}B_{ct})_{t\ge0}\overset d=(B_t)_{t\ge0}
\]
作为路径空间随机元成立。连续性在最后一步把有限维相等升级为过程分布相等。
\end{example}

\begin{example}[平稳 Ornstein--Uhlenbeck Gaussian 过程]
\label{ex:6-4-1-ou-gaussian}
固定 $\theta,\eta>0$，并在 $\R$ 上定义
\[
 K(s,t)=\frac{\eta^2}{2\theta}e^{-\theta|t-s|}.
\]
这个核半正定，因为
\[
 K(s,t)=\eta^2\int_{-\infty}^{s\wedge t}
 e^{-\theta(s-u)}e^{-\theta(t-u)}\dd u,
\]
从而对任意有限数列 $(a_j,t_j)$，
\[
 \sum_{j,k}a_ja_kK(t_j,t_k)
 =\eta^2\int_{-\infty}^{\infty}
 \left(\sum_j a_je^{-\theta(t_j-u)}\1_{\{u\le t_j\}}\right)^2\dd u\ge0.
\]
Kolmogorov 扩张定理因而给出中心 Gaussian 过程 $X$，其协方差为 $K$。由
$1-e^{-x}\le x$，
\[
 \E|X_t-X_s|^2
 =\frac{\eta^2}{\theta}\bigl(1-e^{-\theta|t-s|}\bigr)
 \le \eta^2|t-s|.
\]
若 $p>2$、$C_p=\E|Z|^p$ 且 $Z\sim N(0,1)$，则 Gaussian 尺度关系给出
\[
 \E|X_t-X_s|^p
 =C_p\bigl(\E|X_t-X_s|^2\bigr)^{p/2}
 \le C_p\eta^p|t-s|^{p/2}.
\]
在定理~\ref{theorem:6-1-3-1} 中取 $\alpha=p$、参数维数 $d=1$ 与
$\beta=p/2-1>0$，便得到连续版本。由于均值为零且 $K(s,t)$ 只依赖 $|t-s|$，
任意平移前后的有限维 Gaussian 向量具有相同均值与协方差，故该版本平稳。更具体地，对 $h>0$ 令
\[
 \varepsilon_{t,h}=X_{t+h}-e^{-\theta h}X_t.
\]
对每个 $s\le t$，有
\[
 \Cov(\varepsilon_{t,h},X_s)
 =K(t+h,s)-e^{-\theta h}K(t,s)=0.
\]
联合 Gaussian 性与单调类论证说明 $\varepsilon_{t,h}$ 独立于
$\sigma(X_s:s\le t)$，且
\[
 \varepsilon_{t,h}\sim
 N\!\left(0,\frac{\eta^2}{2\theta}(1-e^{-2\theta h})\right).
\]
这就是平稳 Ornstein--Uhlenbeck 过程的 Gaussian 构造。例~\ref{ex:6-4-3-ou-sde}
从随机微分方程得到同一协方差核，因而给出这个过程的另一种实现。
\end{example}

令 $\mathcal F_t^0=\sigma(B_s:0\le s\le t)$ 为 Brownian 的自然滤过。独立增量不仅与一个
过去时刻独立，还与全部过去信息独立：先在有限坐标柱集上使用命题
\ref{proposition:6-4-1-2}，再以 $\pi$--$\lambda$ 定理推广到 $\mathcal F_t^0$。

\begin{proposition}[Markov 与鞅性质]\label{proposition:6-4-1-3}
Brownian motion 关于其自然滤过是平方可积鞅。若
\[
 q_t(y)=(2\pi t)^{-1/2}e^{-y^2/(2t)},\qquad
 P_tf(x)=\int_\R f(x+y)q_t(y)\dd y
\]
（$t>0$，且 $P_0f=f$），则对每个有界 Borel 函数 $f$ 及 $s,t\ge0$，
\[
 \E[f(B_{s+t})\mid\mathcal F_s^0]=P_tf(B_s)\quad\text{a.s.}
\]
\end{proposition}

\begin{proof}
$B_t\in L^2$ 且适应。增量 $B_{s+t}-B_s$ 独立于 $\mathcal F_s^0$、均值为零，所以
\[
 \E[B_{s+t}\mid\mathcal F_s^0]
 =B_s+\E[B_{s+t}-B_s\mid\mathcal F_s^0]=B_s.
\]
这证明鞅性。

令 $Z=B_{s+t}-B_s$。它独立于 $\mathcal F_s^0$，其分布的密度为 $q_t$。先对
$f=\1_A$ 和 $C\in\mathcal F_s^0$ 计算
\begin{align*}
 \E[\1_C\1_A(B_s+Z)]
 &=\E\left[\1_C\int_\R\1_A(B_s+y)q_t(y)\dd y\right]\\
 &=\E[\1_C P_t\1_A(B_s)].
\end{align*}
第一式先在 $C$ 与 $B_s$ 的简单事件上由独立性得到，再由函数单调类推广；Tonelli 保证参数积分
可测。对 $f$ 作简单逼近和正负分解，便得到全部有界 Borel 函数的条件等式。
\end{proof}

固定时刻的 Markov 性不回答“第一次命中以后”会发生什么。随机时刻的结论需要连续时间停时及其停止信息。

\begin{definition}[连续时间停时与停止信息]\label{def:6-4-1-continuous-stopping}
设 $(\mathcal F_t)_{t\ge0}$ 是滤过。取值于 $[0,\infty]$ 的随机变量 $\tau$ 若
$\{\tau\le t\}\in\mathcal F_t$ 对每个 $t\ge0$ 成立，则称为连续时间停时。定义
\[
 \mathcal F_\infty=\sigma\!\left(\bigcup_{t\ge0}\mathcal F_t\right),
\]
以及
\[
 \mathcal F_\tau=
 \{A\in\mathcal F_\infty:
 A\cap\{\tau\le t\}\in\mathcal F_t\text{ 对每个 }t\ge0\}.
\]
\end{definition}

补集和可数并的分配律说明 $\mathcal F_\tau$ 是 $\sigma$-代数。若 $\sigma\le\tau$，则
$\mathcal F_\sigma\subset\mathcal F_\tau$，证明与引理~\ref{lemma:6-2-3-stopped-sigma} 相同，
只需把整数 $n$ 换成实数 $t$。

\begin{example}[Brownian 首次命中]
先在路径不连续的共同零集上把 $B$ 改为零路径；这保持每个固定时刻的版本关系，也不改变下文的
完备右连续化滤过。因而以下命中事件的路径连续性等式是逐点成立的，而不只是相差一个未必属于
原始自然滤过的零集。
固定 $a\in\R$，令
\[
 \tau_a=\inf\{t\ge0:B_t=a\},
\]
并约定空集下确界为 $\infty$。对 $t\ge0$，路径连续性给出
\[
 \{\tau_a\le t\}
 =\left\{\inf_{q\in\Q\cap[0,t]}|B_q-a|=0\right\}.
\]
“$\subset$”由命中时刻附近的有理数逼近得到；反向若下确界为零，取相应有理序列，再从紧区间
$[0,t]$ 中取收敛子列，连续性给出某个时刻的取值为 $a$。右端由可数多个自然滤过可测随机变量构成，
所以 $\tau_a$ 是停时。$\mathcal F_{\tau_a}$ 保存的是命中发生时已经看见的整段路径信息，而不只是
数值 $\tau_a$。
\end{example}

令 $\mathcal N$ 为 $\mathcal F_\infty^0=\sigma(\bigcup_{t\ge0}\mathcal F_t^0)$ 中所有概率零集的
子集所成的族，并明确置
\[
 \mathcal F_t^B=\bigcap_{u>t}(\mathcal F_u^0\vee\mathcal N).
\]
这是自然滤过的完备右连续化。从上方逼近随机时刻时，右连续性保证极限信息与停止信息相容；
完备化则允许在共同零集上修改连续版本。

\begin{theorem}[Brownian 强 Markov 性]\label{theorem:6-4-1-strong-markov}
设 $\tau<\infty$ a.s. 是关于 $(\mathcal F_t^B)$ 的停时。则
\[
 \widetilde B_t=B_{\tau+t}-B_\tau,\qquad t\ge0,
\]
是独立于 $\mathcal F_\tau^B$ 的标准 Brownian motion。
\end{theorem}

\begin{proof}
先说明确定时刻的独立增量对右连续完备化仍成立。固定 $s$、
$0<t_1<\cdots<t_m$ 和有界连续 $\Phi:\R^m\to\R$。若 $s_k\downarrow s$，则
$A\in\mathcal F_s^B$ 属于每个 $\mathcal F_{s_k}^0\vee\mathcal N$。把未来向量从 $s_k$ 起算，原始
独立增量与完备化不改变零集给出
\[
 \E\bigl[\1_A\Phi(B_{s_k+t_j}-B_{s_k}:1\le j\le m)\bigr]
 =\Prob(A)c_\Phi,
\]
其中 $c_\Phi$ 是相同 Brownian 增量向量上的期望。路径连续性和 DCT 令 $k\to\infty$，得到同一
等式从 $s$ 起成立。函数单调类再把 $\Phi$ 从有界连续函数推广到有界 Borel 函数。

先假设 $\tau\le T$。令
\[
 \tau_n=2^{-n}\lceil2^n\tau\rceil.
\]
这是从上趋于 $\tau$ 的可数值停时，且 $\tau_n\le T+1$。由
$\tau\le\tau_n$，有 $\mathcal F_\tau^B\subset\mathcal F_{\tau_n}^B$。固定
$A\in\mathcal F_\tau^B$，按 $\tau_n=k2^{-n}$ 分片，并应用上段确定时刻的独立增量，得到
\begin{align*}
 &\E\bigl[\1_A\Phi(B_{\tau_n+t_j}-B_{\tau_n}:1\le j\le m)\bigr]\\
 &\quad=\sum_k
 \E\bigl[\1_{A\cap\{\tau_n=k2^{-n}\}}
 \Phi(B_{k2^{-n}+t_j}-B_{k2^{-n}}:j)\bigr]\\
 &\quad=\sum_k\Prob(A\cap\{\tau_n=k2^{-n}\})c_\Phi
 =\Prob(A)c_\Phi.
\end{align*}
路径连续性使被积的未来向量 a.s. 收敛到从 $\tau$ 起算的向量；DCT 因而把等式传到极限。
所以每个有限未来向量具有 Brownian 分布，并与 $\mathcal F_\tau^B$ 独立。

若 $\tau$ 只假设 a.s. 有限，取 $\tau\wedge r$。对
$A_r=A\cap\{\tau\le r\}$，由停止信息定义可核验
$A_r\in\mathcal F_{\tau\wedge r}^B$。在 $A_r$ 上，从 $\tau\wedge r$ 起算的增量就是从 $\tau$
起算的增量，故有
\[
 \E[\1_{A_r}\Phi(B_{\tau+t_j}-B_\tau:j)]
 =\Prob(A_r)c_\Phi.
\]
令 $r\to\infty$ 并用 DCT，得到对 $A$ 的同一等式。

最后，$\widetilde B$ 逐路径连续且从零出发。上述有限维结论说明它具有 Brownian 有限维分布；
式~\eqref{eq:6-4-1-wiener-borel} 和柱事件单调类说明整个路径随机元独立于
$\mathcal F_\tau^B$。因此它是独立于停止信息的 Brownian motion。
\end{proof}

特别地，在 $\tau_a<\infty$ 的事件上，命中 $a$ 后的位移从零重新开始。若要无条件套用上一定理，
仍须先证明 $\tau_a<\infty$ a.s.，或先在 $\tau_a\wedge T$ 上工作；“首次命中”这个名字不自动保证
命中发生。

\begin{proposition}[连续时间可选抽样]\label{proposition:6-4-1-continuous-optional}
设 $M$ 是关于一个满足通常条件的滤过的连续鞅，$\sigma\le\tau\le T<\infty$ 为停时。则
\[
 \E[M_\tau\mid\mathcal F_\sigma]=M_\sigma\quad\text{a.s.}
\]
若 $\tau$ 无界，只有在停止族 $\{M_{t\wedge\tau}:t\ge0\}$ 一致可积，或另有足以交换极限与
期望的支配条件时，才可令 $T\to\infty$。
\end{proposition}

\begin{proof}
在 $[0,T]$ 上取网格
$D_n=\{k2^{-n}:0\le k2^{-n}<T\}\cup\{T\}$，并令
$\rho_n$ 是 $D_n$ 中不小于停时 $\rho$ 的最小点。分别应用于 $\sigma,\tau$，得到
$\sigma_n\le\tau_n$、$\sigma_n\downarrow\sigma$、$\tau_n\downarrow\tau$。沿有限网格抽样的
序列 $(M_t)_{t\in D_n}$ 是离散鞅。离散有界可选抽样定理~\ref{theorem:6-2-3-2} 给出
\[
 \E[M_{\tau_n}\mid\mathcal F_{\sigma_n}]=M_{\sigma_n}.
\]

每个确定 $t\le T$ 都满足 $M_t=\E[M_T\mid\mathcal F_t]$。在同一有限网格上再对
$\rho_n$ 与终点 $T$ 使用离散可选抽样，得到
\[
 M_{\rho_n}=\E[M_T\mid\mathcal F_{\rho_n}].
\]
固定 $Y=M_T\in L^1$ 时，所有条件期望 $\E[Y\mid\mathcal G]$ 组成一致可积族：对
$A_K=\{|\E(Y\mid\mathcal G)|>K\}$，条件模长不等式及
$\Prob(A_K)\le\E|Y|/K$ 给出
\[
 \E[|\E(Y\mid\mathcal G)|\1_{A_K}]
 \le\E[|Y|\1_{\{|Y|>R\}}]+R\frac{\E|Y|}{K},
\]
先令 $R\to\infty$，再令 $K\to\infty$，所得界与 $\mathcal G$ 无关。因此
$(M_{\sigma_n})$ 与 $(M_{\tau_n})$ 共同 UI。

路径连续性给出 a.s. 收敛
$M_{\sigma_n}\to M_\sigma$、$M_{\tau_n}\to M_\tau$；Vitali 定理把它们升级为 $L^1$ 收敛。
若 $A\in\mathcal F_\sigma$，则 $A\in\mathcal F_{\sigma_n}$，所以网格条件等式给
\[
 \E[\1_AM_{\tau_n}]=\E[\1_AM_{\sigma_n}].
\]
令 $n\to\infty$，$L^1$ 收敛给
$\E[\1_AM_\tau]=\E[\1_AM_\sigma]$。这对所有 $A\in\mathcal F_\sigma$ 成立，正是所需条件
期望等式。这一极限过程使用的是停止值的 $L^1$ 收敛，因而不需要比较随 $n$ 变化的条件 $\sigma$-代数。
\end{proof}

连续并不意味着平滑。Brownian 的二进增量虽逐项趋零，绝对值之和却发散，平方和则保持有限。
更棘手的是“处处不可微”：对每个固定 $t$ 证明不可微，只得到一族依赖 $t$ 的零集，不能对不可数个
$t$ 取交。下面的锥形估计把所有可能的随机位置压缩成每层有限多个格子。

\begin{theorem}[路径粗糙性（基础版）]\label{theorem:6-4-1-4}
几乎必然，Brownian 路径在每个非退化时间区间上的总变差为 $+\infty$，并且在每个
$t>0$ 都不可微；在 $t=0$ 也不存在有限右导数。
\end{theorem}

\begin{proof}
先在 $[0,1]$ 上处理总变差。令
\[
 V_n=\sum_{k=0}^{2^n-1}
 |B_{(k+1)2^{-n}}-B_{k2^{-n}}|.
\]
这些增量独立，且各自与 $2^{-n/2}Z$ 同分布，其中 $Z\sim N(0,1)$。记
$m_1=\E|Z|=\sqrt{2/\pi}$。则
\[
 \E V_n=m_1\,2^{n/2},
 \qquad
 \Var(V_n)=2^n2^{-n}\Var(|Z|)=1-\frac2\pi.
\]
Chebyshev 不等式给出
\[
 \Prob\!\left(V_n\le\frac12\E V_n\right)
 \le\frac{4\Var(V_n)}{(\E V_n)^2}
 =C2^{-n}.
\]
右端可和，Borel--Cantelli I 因而说明 $V_n\to\infty$ a.s.。路径在 $[0,1]$ 上的总变差
至少为每个 $V_n$，故为无穷。对每个有理端点区间 $[a,b]\subset[0,\infty)$ 作仿射二进
分割，同一计算只把常数乘以 $b-a$。对可数多个有理区间取交后，每个非退化区间都包含其中一个，
总变差的单调性给出第一项结论。

下面处理不可微性。固定有理闭区间 $J=[a,b]\subset(0,\infty)$、整数 $r\ge1$。令
$\delta_n=2^{-n}$，并在包含 $J$ 的二进格网上记
\[
 \Delta_{k,n}=B_{(k+1)\delta_n}-B_{k\delta_n}.
\]
若某个 $t\in J$ 满足锥形界
\begin{equation}\label{eq:6-4-1-cone-bound}
 |B_s-B_t|\le r|s-t|
 \quad\text{只要 }|s-t|\le\eta
\end{equation}
对某个 $\eta>0$ 成立，则当 $n$ 充分大、$2\delta_n<\eta$ 时，取唯一的 $k$ 使
$t\in[k\delta_n,(k+1)\delta_n)$。四个格点
$(k-1)\delta_n,k\delta_n,(k+1)\delta_n,(k+2)\delta_n$ 到 $t$ 的距离都不超过
$2\delta_n$。由三角不等式，包围该格的三个增量满足
\begin{equation}\label{eq:6-4-1-three-small-increments}
 |\Delta_{j,n}|\le4r\delta_n,
 \qquad j=k-1,k,k+1.
\end{equation}

固定 $k$ 时，这三个变量独立且服从 $N(0,\delta_n)$。标准正态密度不超过
$(2\pi)^{-1/2}$，所以
\[
 \Prob(|\Delta_{j,n}|\le4r\delta_n)
 =\Prob(|Z|\le4r\sqrt{\delta_n})
 \le C r\sqrt{\delta_n}.
\]
与区间 $J$ 相交的二进格至多 $C_J\delta_n^{-1}$ 个。对所有可能的 $k$ 取并，得到
\begin{align*}
 &\Prob\bigl(\text{某个与 $J$ 相交的格满足
 \eqref{eq:6-4-1-three-small-increments}}\bigr)\\
 &\qquad\le C_J\delta_n^{-1}(Cr\sqrt{\delta_n})^3
 \le C_{J,r}2^{-n/2}.
\end{align*}
该概率关于 $n$ 可和。Borel--Cantelli I 说明几乎必然只有有限多层出现这样的格。然而
式~\eqref{eq:6-4-1-cone-bound} 会由上面的确定性推导迫使每个充分细层都出现一个，故在概率一
事件上不存在满足该锥形界的 $t\in J$。

若路径在某点 $t\in J$ 有有限导数 $L$，由导数定义可取整数 $r>|L|+1$ 与 $\eta>0$，使
式~\eqref{eq:6-4-1-cone-bound} 成立。对可数个 $r$ 和可数个有理区间 $J$ 取交，便排除了所有
$t>0$；没有对随机不可数点逐点取零集。对 $t=0$，若存在有限右导数，同一锥形界在
$[0,\eta]$ 成立，并迫使前三个二进增量都满足 $O(r\delta_n)$；其概率为
$O_r(2^{-3n/2})$，仍可和。故有限右导数也不存在。
\end{proof}

可微性只给出以 $t$ 为中心的锥形界，不要求路径在 $t$ 的整个邻域上 Lipschitz。
三个相邻增量把这条点态信息转化为可数的格点事件。总变差与不可微也是两个不同的结论：连续有限变差函数可能在某些点不可微，
而无限变差本身不排除在其他点存在导数。

\begin{figure}[htbp]
\centering
\begin{tikzpicture}[node distance=9mm and 10mm]
  \begin{scope}[xshift=0cm]
    \draw[book line,-{Stealth[length=2mm]}] (0,0)--(4.2,0);
    \draw[BookWarm,semithick]
      plot coordinates {(0,0) (.45,.35) (.9,.05) (1.35,.65) (1.8,.3) (2.25,.8) (2.7,.45) (3.15,.95) (3.6,.7) (4.05,1.05)};
    \draw[BookAccent,thick,smooth]
      plot coordinates {(0,0) (.45,.22) (.9,.18) (1.35,.5) (1.8,.32) (2.25,.65) (2.7,.55) (3.15,.82) (3.6,.72) (4.05,.96)};
    \node[book label] at (2,-.42) {扩散缩放下的折线};
  \end{scope}
  \node[book strong node] (kernel) at (6.0,.55) {半正定核\\$s\wedge t$};
  \node[book strong node,right=of kernel] (extend) {Kolmogorov\\扩张};
  \node[book warm node,right=of extend] (version) {矩估计\\连续版本};
  \draw[book accent arrow] (kernel)--(extend);
  \draw[book accent arrow] (extend)--(version);
\end{tikzpicture}
\caption{随机游走只提供动机；Brownian 的严格存在依次使用 Gaussian 核、有限维扩张和连续版本。}
\label{fig:6-4-1-construction}
\end{figure}

\begin{exercise}
\begin{enumerate}
  \item 不引用例~\ref{ex:6-4-1-1}，把 $t_1<\cdots<t_n$ 时的矩阵
  $(t_i\wedge t_j)$ 写成一个下三角矩阵与其转置之积。
  \item 对两个不交区间逐项计算 Brownian 增量的协方差，再用特征函数说明联合 Gaussian 的
  零协方差为何给出独立，而一般变量不能这样推断。
  \item 从有限维均值和协方差出发证明 Brownian scaling，并使用
  式~\eqref{eq:6-4-1-wiener-borel} 完成路径空间分布的最后一步。
  \item 核验 Brownian bridge 的协方差，并计算 $\Law(\beta_t)$ 及
  $\Cov(\beta_s,\beta_t-\beta_s)$；据此说明它没有独立增量。
\end{enumerate}
\end{exercise}

Brownian 路径的无限变差使经典 Riemann--Stieltjes 理论失去直接入口。\Ito{} 积分先在可预测阶梯过程上
利用独立增量定义，再由 $L^2$ 完备性延拓到一般被积过程。


\subsection{\Ito{} 积分、等距与二次变差}\label{subsec:6-4-2}
\index{Ito@\Ito{} 积分}\index{可预测过程}\index{二次变差}

若把 $\int_0^TB_t\dd B_t$ 当作普通 Riemann--Stieltjes 积分，左端点与右端点会给出不同答案。
这不是选点技巧不够精致，而是 Brownian 路径没有有限变差。\Ito{} 的做法是改变问题：被积过程在每个
小区间开始时必须已经可知，积分先作为平方可积随机变量定义。独立增量随后把不同小区间变成
$L^2(\Omega)$ 中的正交方向。

\begin{definition}[向量 Brownian motion 与 Hilbert--Schmidt 范数]
\label{def:6-4-2-vector-brownian}
相对于同一滤过 $(\mathcal F_t)$，连续适应过程
$B=(B^1,\ldots,B^r)$ 称为 $\R^r$ 值 Brownian motion，如果 $B_0=0$，并且对
$0\le s<t$，增量 $B_t-B_s$ 独立于 $\mathcal F_s$，其分布是均值为零、协方差矩阵为
$(t-s)I_r$ 的 Gaussian 分布。于是各坐标是相互独立的一维 Brownian motion；向量随机积分
始终相对于这一共同滤过，并按有限和
\[
 \left(\int_0^t\sigma_s\dd B_s\right)^i
 =\sum_{a=1}^r\int_0^t\sigma_{ia}(s)\dd B_s^a
\]
逐行定义。对实 $d\times r$ 矩阵 $A=(a_{ij})$，记
\[
 \|A\|_{\mathrm{HS}}^2=\sum_{i=1}^d\sum_{j=1}^r|a_{ij}|^2
 =\operatorname{tr}(AA^T).
\]
\end{definition}

以下固定一个满足通常条件的滤过 $(\mathcal F_t)_{t\ge0}$，并假定 $B$ 相对于它是
上述向量 Brownian motion；若没有加入额外信息，可以取前一小节构造的
$(\mathcal F_t^B)_{t\ge0}$。本小节的可预测性、渐进可测性、适应性与停时均相对于
$(\mathcal F_t)$ 而言。

\begin{definition}[简单可预测过程及其积分]\label{def:6-4-2-1}
在固定时域 $[0,T]$ 上，若
\[
 H_t=\sum_{k=0}^{n-1}H_k\1_{(t_k,t_{k+1}]}(t),
 \qquad0=t_0<t_1<\cdots<t_n=T,
\]
其中每个 $H_k$ 有界且对 $\mathcal F_{t_k}$ 可测，则称 $H$ 为\emph{简单可预测过程}，并定义
\[
 \int_0^TH_t\dd B_t
 :=\sum_{k=0}^{n-1}H_k(B_{t_{k+1}}-B_{t_k}).
\]
对 $0\le t\le T$，$\int_0^tH_s\dd B_s$ 用截断后的同一有限和定义。
\end{definition}

若换一个时间分割表示同一个 $H$，取两个分割的共同加细即可看到积分值不变。系数必须在左端点可测；
它可以依赖过去，但不能查看即将相乘的未来增量。

\begin{definition}[$L^2$ 可预测类]\label{def:6-4-2-2}
令 $\mathcal P$ 是 $\Omega\times[0,T]$ 上由
\[
 A\times(s,t]\quad(A\in\mathcal F_s, 0\le s<t\le T),
 \qquad A_0\times\{0\}\quad(A_0\in\mathcal F_0)
\]
生成的\emph{可预测 $\sigma$-代数}（predictable $\sigma$-algebra）。定义
\[
 L^2_{\mathrm{pred}}([0,T])
 =L^2(\Omega\times[0,T],\mathcal P,\Prob\otimes\dd t),
\]
并记
\[
 \|H\|_{\mathcal H^2}^2=\E\int_0^T|H_t|^2\dd t.
\]
其中等价关系是 $\Prob\otimes\dd t$-a.e. 相等。
\end{definition}

有界简单可预测过程在该空间中稠密。事实上，生成 $\mathcal P$ 的矩形
$A\times(s,t]$ 的指标函数本身就是简单可预测过程，而
$A_0\times\{0\}$ 在 $\Prob\otimes\dd t$ 下为零；有限交、差和有限并经过时间端点共同加细后
仍有这种表示。由可测简单函数逼近和截断，
这些指标函数的有限线性组合在有限测度空间
$(\Omega\times[0,T],\Prob\otimes\dd t)$ 的 $L^2$ 中稠密。这也说明完成对象及其等价关系已经明确，
因而积分的完成域是具体的乘积测度 $L^2$ 空间。

\begin{definition}[渐进可测过程与局部鞅]\label{def:6-4-2-progressive-local}
适应过程 $X$ 称为\emph{渐进可测}（progressively measurable），如果对每个 $t\ge0$，限制
$(\omega,s)\mapsto X_s(\omega)$（$0\le s\le t$）关于
$\mathcal F_t\otimes\Borel([0,t])$ 可测。适应过程 $M$ 称为\emph{局部鞅}，如果存在停时
$\tau_n\uparrow\infty$ a.s.，使每个停止过程
$M^{\tau_n}_t=M_{t\wedge\tau_n}$ 是鞅。
\end{definition}

适应且路径连续的过程是渐进可测的：在 $[0,t]$ 上用左端点阶梯过程逼近，阶梯过程关于
$\mathcal F_t\otimes\Borel([0,t])$ 可测，路径连续性给出逐点极限。适应且左连续的过程是
可预测，同一逼近中的系数在每个区间左端以前已经可知。连续过程同时满足两项；左连续本身
并不表示路径连续。

\begin{remark}[渐进可测系数为何仍能进入 Brownian 积分]
\label{remark:6-4-2-progressive-predictable}
若 $H$ 渐进可测且满足
\[
 \E\int_0^T|H_s|^2\dd s<\infty,
\]
则把 $H$ 在负时间延为零，并令
\[
 H_t^{(n)}=n\int_{(t-1/n)^+}^{t}H_s\dd s.
\]
每条局部可积路径的过去平均关于 $t$ 连续；它在时刻 $t$ 只使用 $s\le t$ 的信息，所以
$H^{(n)}$ 适应，从而可预测。把 $H$ 视为
$L^2(\Omega\times\R)$ 元素，Minkowski 不等式给出
\[
 \|H^{(n)}-H\|_2
 \le n\int_0^{1/n}\|H_{\bullet-u}-H_\bullet\|_2\dd u.
\]
第四章的 $L^2$ 平移连续性定理~\ref{theorem:4-4-1-1} 作用于时间变量，说明右端趋于零。
于是某个子列 a.e. 收敛；其极限是
可预测函数，并与 $H$ 在 $\Prob\otimes\dd t$ 意义下相等。因此渐进可测的平方可积
等价类有可预测代表。后面 $H_t=\sigma(t,X_t)$ 的情形正通过这个平均化进入 \Ito{} 积分；
只说“联合可测”不足以替代这一步。
\end{remark}

\begin{definition}[二次变差]\label{def:6-4-2-3}
对连续过程 $X$、$[0,T]$ 的有限分割 $\pi$，定义
\[
 [X]^{\pi}_T=\sum_{[u,v]\in\pi}(X_v-X_u)^2.
\]
给定分割列 $\pi_n$ 且网格 $|\pi_n|\to0$，若
$[X]^{\pi_n}_T$ 依概率趋于某随机变量 $[X]_T$，则称它为 $X$ 沿该分割列的二次变差。
\end{definition}

对一般连续函数，上述极限可能依赖分割列。定理~\ref{theorem:6-4-2-3} 与
命题~\ref{proposition:6-4-2-4} 分别说明，Brownian motion 与 \Ito{} 积分对任意确定的网格趋零分割列都有同一极限。

\begin{theorem}[\Ito{} 等距]\label{theorem:6-4-2-1}
对每个实简单可预测过程 $H$，
\[
 \E\left|\int_0^TH_t\dd B_t\right|^2
 =\E\int_0^T|H_t|^2\dd t.
\]
特别地，简单过程积分映射是从 $L^2_{\mathrm{pred}}([0,T])$ 的稠密子空间到
$L^2(\Omega)$ 的线性等距。
\end{theorem}

\begin{proof}
记 $\Delta B_k=B_{t_{k+1}}-B_{t_k}$、$\Delta t_k=t_{k+1}-t_k$。若 $j<k$，则
$H_j\Delta B_jH_k$ 对 $\mathcal F_{t_k}$ 可测，而
\[
 \E[\Delta B_k\mid\mathcal F_{t_k}]=0.
\]
因此塔式性质给出
\[
 \E[H_j\Delta B_jH_k\Delta B_k]=0.
\]
对角项中，独立增量给出
\[
 \E[(\Delta B_k)^2\mid\mathcal F_{t_k}]=\Delta t_k,
\]
所以
\[
 \E[H_k^2(\Delta B_k)^2]=\E[H_k^2]\Delta t_k.
\]
展开有限和的平方，全部交叉项消失，得到
\begin{align*}
 \E\left|\sum_kH_k\Delta B_k\right|^2
 &=\sum_k\E[H_k^2]\Delta t_k\\
 &=\E\int_0^T|H_t|^2\dd t.
\end{align*}
这也证明若两个简单表示在乘积测度意义下相等，它们的积分在 $L^2(\Omega)$ 中相等。
\end{proof}

\begin{theorem}[\Ito{} 积分的完成]\label{theorem:6-4-2-2}
简单过程积分唯一延拓为线性等距
\[
 I:L^2_{\mathrm{pred}}([0,T])\longrightarrow L^2(\Omega).
\]
对 $H\in L^2_{\mathrm{pred}}([0,T])$ 定义
\[
 M_t=I(H\1_{(0,t]})=\int_0^tH_s\dd B_s.
\]
则 $M$ 有连续版本；取该版本后，$M$ 是从零出发的平方可积鞅。
\end{theorem}

\begin{proof}
$L^2(\Omega)$ 完备，定理~\ref{theorem:6-4-2-1} 的等距把稠密子空间上的积分唯一延拓到完成。
对固定 $t$，取简单 $H^{(n)}\to H$ 于 $L^2_{\mathrm{pred}}$，则
$I_t(H^{(n)})\to I_t(H)$ 于 $L^2$；每个前者对 $\mathcal F_t$ 可测，故 $M_t$ 也可测。

若 $s<t$，简单过程的增量积分对 $\mathcal F_s$ 条件均值为零，证明与等距中交叉项的计算相同。
条件期望是 $L^2$ 收缩，令简单逼近趋于 $H$ 得
\[
 \E[M_t-M_s\mid\mathcal F_s]=0.
\]
所以 $(M_t)$ 在每个固定时刻族上是平方可积鞅。

还需同时选择连续版本。简单过程的积分过程逐路径连续。对简单 $G$，在任意有限时间网格上把离散
Doob $L^2$ 不等式~\ref{theorem:6-2-2-2} 用于鞅 $I_t(G)$，得到
\[
 \E\max_{t\in D}|I_t(G)|^2\le4\E|I_T(G)|^2.
\]
让有限网格递增到 $[0,T]$ 的一个可数稠密集；连续性与单调收敛定理给出
\begin{equation}\label{eq:6-4-2-continuous-doob}
 \E\sup_{t\le T}|I_t(G)|^2
 \le4\E\int_0^T|G_s|^2\dd s.
\end{equation}
因此若 $H^{(n)}$ 是 $H$ 的简单逼近，则
\[
 \E\sup_{t\le T}|I_t(H^{(n)}-H^{(m)})|^2
 \le4\|H^{(n)}-H^{(m)}\|_{\mathcal H^2}^2.
\]
所以积分过程在 Banach 值空间 $L^2(\Omega;C[0,T])$ 中为 Cauchy 列。其极限 $\widetilde M$
有连续路径；对每个 $t$，上确界范数收敛蕴含
$I_t(H^{(n)})\to\widetilde M_t$ 于 $L^2$，而端点等距又给极限 $M_t$，故
$\widetilde M_t=M_t$ a.s.。因此 $\widetilde M$ 是一个在全部时刻同时相容的连续版本。

最后，连续版本在有理时刻满足鞅条件；取有理数从右逼近任意 $s<t$，用路径连续性、
式~\eqref{eq:6-4-2-continuous-doob} 给出的 $L^2$ 控制及条件期望收缩，把等式传到任意实时间。
\end{proof}

\begin{example}[确定被积函数]\label{ex:6-4-2-1}
若 $h\in L^2[0,T]$ 是确定函数，取确定简单函数 $h_n\to h$ 于 $L^2$。每个
$\int h_n\dd B$ 是 Brownian 增量的有限线性组合，因而中心 Gaussian，方差为
$\int h_n^2\dd t$。\Ito{} 等距给出积分在 $L^2$ 中收敛；对特征函数使用
$|e^{ix}-e^{iy}|\le|x-y|$，可把 Gaussian 特征函数传到极限，得到
\[
 \int_0^Th(t)\dd B_t\sim N\!\left(0,\int_0^Th(t)^2\dd t\right).
\]
取 $h=\1_{(a,b]}$ 时，定义恢复 $B_b-B_a$。
\end{example}

平方可积完成仍要求 $\E\int H^2<\infty$。SDE 中常见的被积过程只在每条路径上局部平方可积；
在逐层截断能量之前，先证明积分与停止相容。

\begin{lemma}[\Ito{} 积分的停止恒等式]\label{lemma:6-4-2-stopping-integral}
设 $G\in L^2_{\mathrm{pred}}([0,T])$，
$M_t=\int_0^tG_s\dd B_s$，且 $\tau\le T$ 是停时。则对每个 $0\le t\le T$，
\[
 \int_0^tG_s\1_{\{s\le\tau\}}\dd B_s=M_{t\wedge\tau}
 \quad\text{a.s.}
\]
\end{lemma}

\begin{proof}
先令
\[
 \tau_n=2^{-n}\lceil2^n\tau\rceil\wedge T,
\]
并把 $T$ 加入相应网格；则 $\tau_n$ 是有限值停时且 $\tau_n\downarrow\tau$。若 $G$ 是简单
可预测过程，把它的确定跳点与该网格共同加细。对任一相邻区间 $(r_k,r_{k+1}]$，
\[
 \{r_{k+1}\le\tau_n\}
 =\{\tau_n\le r_k\}^{\,c}\in\mathcal F_{r_k},
\]
其中等式使用了 $\tau_n$ 在共同网格的相邻端点之间不取值。因此
$G_s\1_{\{s\le\tau_n\}}$ 仍是简单可预测过程；逐项展开定义中的有限和便得到
\[
 \int_0^tG_s\1_{\{s\le\tau_n\}}\dd B_s=M_{t\wedge\tau_n}.
\]

对一般 $G$，取简单 $G^{(m)}\to G$ 于 $L^2_{\mathrm{pred}}([0,T])$。\Ito{} 等距控制上式左边的
极限，而连续 Doob 估计~\eqref{eq:6-4-2-continuous-doob} 给出
\[
 \E\sup_{t\le T}\left|
 \int_0^t(G_s^{(m)}-G_s)\dd B_s\right|^2
 \le4\|G^{(m)}-G\|_{\mathcal H^2}^2\longrightarrow0.
\]
所以停止于 $\tau_n$ 的右边也在 $L^2$ 中收敛，有限值停时的恒等式对一般 $G$ 成立。

最后令 $n\to\infty$。DCT 给
$G\1_{\{s\le\tau_n\}}\to G\1_{\{s\le\tau\}}$ 于
$L^2(\Prob\otimes\dd s)$，故左边由 \Ito{} 等距收敛。另一方面，$M$ 连续，且
$\E\sup_{t\le T}|M_t|^2<\infty$，所以
$M_{t\wedge\tau_n}\to M_{t\wedge\tau}$ 于 $L^2$。取极限即得结论。
\end{proof}

\begin{proposition}[局部平方可积被积过程]\label{proposition:6-4-2-localization}
设 $H$ 可预测，且对每个 $T<\infty$，
\[
 \int_0^T|H_s|^2\dd s<\infty\quad\text{a.s.}
\]
令
\[
 \tau_n=\inf\left\{t\ge0:\int_0^t|H_s|^2\dd s\ge n\right\}\wedge n.
\]
则 $\tau_n$ 是停时、$\tau_n\uparrow\infty$ a.s.，且
$H\1_{(0,\tau_n]}$ 在每个有限时域属于 $L^2_{\mathrm{pred}}$。这些停止积分在重叠时域上一致，
可以拼成连续局部鞅
\[
 \int_0^tH_s\dd B_s.
\]
\end{proposition}

\begin{proof}
能量过程 $A_t=\int_0^t|H_s|^2\dd s$ 适应、连续且非降；适应性由对
$\Omega\times[0,t]$ 的可测积分得到。因此
\[
 \{\tau_n\le t\}=
 \begin{cases}
   \{A_t\ge n\},&t<n,\\
   \Omega,&t\ge n,
 \end{cases}
\]
属于 $\mathcal F_t$，所以 $\tau_n$ 是停时。路径局部有限性说明对每个固定 $T$，最终有
$A_T<n$ 且 $n>T$，从而 $\tau_n>T$；故 $\tau_n\uparrow\infty$ a.s.（必要时把它替换为
$\min_{k\ge n}\tau_k$ 的单调版本，原定义本身也因阈值和截断同时增加而非降）。

过程 $t\mapsto\1_{[0,\tau_n]}(t)$ 左连续且适应，因而可预测；它与
$\1_{(0,\tau_n]}$ 只在 $t=0$ 不同，代表同一个
$L^2(\Prob\otimes\dd t)$ 等价类。连续性还给出
\[
 \int_0^T|H_s|^2\1_{\{s\le\tau_n\}}\dd s\le n
\]
对所有 $T$ 成立，故停止后的被积过程属于 $L^2_{\mathrm{pred}}$。记
\[
 M^{(n)}_t=\int_0^tH_s\1_{\{s\le\tau_n\}}\dd B_s.
\]
若 $m\ge n$，对平方可积的被积过程
$H\1_{\{s\le\tau_m\}}$ 在任意固定 $[0,T]$ 上应用
引理~\ref{lemma:6-4-2-stopping-integral} 于停时 $\tau_n\wedge T$，得到对 $t\le T$，
\[
 M^{(m)}_{t\wedge\tau_n}
 =\int_0^tH_s\1_{\{s\le\tau_m\}}\1_{\{s\le\tau_n\}}\dd B_s
 =M^{(n)}_t.
\]
上式先对每个固定 $t$ 几乎必然成立。对可数个
$t\in\Q_+$ 与可数个整数对 $m\ge n$ 取共同的概率一事件，再用所有 $M^{(n)}$ 的连续性，得到一个对
所有 $m\ge n$ 与所有 $t\ge0$ 同时成立的共同零集。在该零集上将各层都定义为零路径。
所以各层在 $\tau_n$ 以前逐路径相容，可以逐层拼接。第 $n$ 层停止过程是平方可积鞅，故拼接过程以
$(\tau_n)$ 为局部化列，是连续局部鞅。
\end{proof}

\begin{example}[局部化回收 $H_t=B_t$]
Brownian 路径在紧时间区间上连续，故
$\int_0^TB_s^2\dd s<\infty$ a.s.，命题~\ref{proposition:6-4-2-localization} 已经定义
$\int_0^tB_s\dd B_s$。事实上
\[
 \E\int_0^TB_s^2\dd s=\int_0^Ts\dd s=\frac{T^2}{2}<\infty,
\]
所以它在每个有限时域还属于未局部化的 $L^2$ 理论。
\end{example}

\begin{theorem}[Brownian 二次变差]\label{theorem:6-4-2-3}
对 $[0,T]$ 的任意确定分割列 $\pi_n$，若 $|\pi_n|\to0$，则
\[
 \sum_{[u,v]\in\pi_n}(B_v-B_u)^2\longrightarrow T
 \quad\text{于 }L^2.
\]
沿二进分割，收敛还 a.s. 成立。
\end{theorem}

\begin{proof}
写 $\Delta t=v-u$、$\Delta B=B_v-B_u$。各分割区间上的增量独立，并且
$\Delta B\sim N(0,\Delta t)$。由 $\E Z^4=3$，
\[
 \E(\Delta B)^2=\Delta t,
 \qquad
 \Var((\Delta B)^2)=3(\Delta t)^2-(\Delta t)^2=2(\Delta t)^2.
\]
故平方增量和的均值为 $T$，方差为
\[
 2\sum_{[u,v]\in\pi_n}(v-u)^2
 \le2T|\pi_n|\longrightarrow0.
\]
这就是 $L^2$ 收敛。

若 $\pi_n$ 是二进分割，则方差等于 $2T^2 2^{-n}$。对每个
$m\ge1$，Chebyshev 不等式说明事件
$\{|[B]^{\pi_n}_T-T|>1/m\}$ 的概率关于 $n$ 可和。先对 $n$ 用 Borel--Cantelli，再对可数个
$m$ 取交，得到 a.s. 收敛。
\end{proof}

若某条连续路径在 $[0,T]$ 上总变差有限，则沿网格趋零的分割
\[
 \sum(\Delta B)^2
 \le\left(\max|\Delta B|\right)\sum|\Delta B|\longrightarrow0,
\]
因为一致连续性使最大增量趋零。二进二次变差却 a.s. 趋于 $T>0$，所以有限变差路径构成零概率
事件。这从二次变差再次回收了定理~\ref{theorem:6-4-1-4} 的总变差结论。

\begin{example}[左端、右端与 Stratonovich 和]\label{ex:6-4-2-2}
令 $\pi=\{0=t_0<\cdots<t_n=T\}$。代数恒等式
\[
 B_{t_{k+1}}^2-B_{t_k}^2
 =2B_{t_k}\Delta B_k+(\Delta B_k)^2
\]
给出
\[
 \sum_kB_{t_k}\Delta B_k
 =\frac12\left(B_T^2-\sum_k(\Delta B_k)^2\right).
\]
左端阶梯过程 $H_t^\pi=B_{t_k}$（$t_k<t\le t_{k+1}$）可预测，并且
\[
 \E\int_0^T|H_t^\pi-B_t|^2\dd t
 =\sum_k\int_{t_k}^{t_{k+1}}(t-t_k)\dd t
 \le\frac T2|\pi|.
\]
所以左端和于 $L^2$ 趋于 $\int_0^TB_t\dd B_t$。结合 Brownian 二次变差，
\[
 \int_0^TB_t\dd B_t=\frac12(B_T^2-T).
\]
右端和满足
\[
 \sum_kB_{t_{k+1}}\Delta B_k
 =\sum_kB_{t_k}\Delta B_k+\sum_k(\Delta B_k)^2
 \longrightarrow\frac12(B_T^2+T).
\]
两者的平均，也就是对称端点的 Stratonovich 和，趋于 $B_T^2/2$。三个答案之差正是二次变差，
所以 \Ito{} 积分不是逐路径选点无关的 Riemann--Stieltjes 积分。
\end{example}

\begin{example}[非适应被积过程]\label{ex:6-4-2-3}
在分割 $\pi$ 上事后选择
\[
 H_k=\sgn(B_{t_{k+1}}-B_{t_k}).
\]
则形式和为
\[
 \sum_kH_k(B_{t_{k+1}}-B_{t_k})=\sum_k|B_{t_{k+1}}-B_{t_k}|.
\]
$H_k$ 依赖区间右端的未来增量，通常不对 $\mathcal F_{t_k}$ 可测。对 $n$ 等分割，形式和的
期望为
\[
 n\E|N(0,T/n)|=\sqrt{\frac{2Tn}{\pi}}\longrightarrow\infty.
\]
这不是 \Ito{} 等距的反例，而是一个不在定义域中的“预知未来”策略；可预测性正是排除这种事后选符号。
\end{example}

\Ito{} 积分自身也有二次变差。以下论证先处理简单被积过程，再用 \Ito{} 等距完成逼近。

\begin{proposition}[\Ito{} 积分的二次变差]\label{proposition:6-4-2-4}
若 $H\in L^2_{\mathrm{pred}}([0,T])$ 且
$M_t=\int_0^tH_s\dd B_s$，则对任意确定分割列 $|\pi_n|\to0$，
\[
 \sum_{[u,v]\in\pi_n}(M_v-M_u)^2
 \longrightarrow\int_0^TH_s^2\dd s
 \quad\text{依概率}.
\]
局部平方可积情形在每个局部化停时之前满足同一结论。
\end{proposition}

\begin{proof}
先设 $H$ 为有界简单可预测过程。若分割细化了 $H$ 的确定跳点，在每个常系数区间内有
$\Delta M=H_k\Delta B$。把细分割记为 $0=r_0<\cdots<r_N=T$，并令 $h_\ell$ 表示
$H$ 在 $(r_\ell,r_{\ell+1}]$ 上的系数。因为所有粗分割端点已被加入，$h_\ell$ 对
$\mathcal F_{r_\ell}$ 可测。置
\[
 Z_\ell=h_\ell^2\bigl((B_{r_{\ell+1}}-B_{r_\ell})^2-(r_{\ell+1}-r_\ell)\bigr).
\]
Gaussian 二阶与四阶矩给出
\[
 \E[Z_\ell\mid\mathcal F_{r_\ell}]=0,
 \qquad
 \E[Z_\ell^2\mid\mathcal F_{r_\ell}]
 =2h_\ell^4(r_{\ell+1}-r_\ell)^2.
\]
若 $\ell<j$，则 $Z_\ell$ 对 $\mathcal F_{r_j}$ 可测，因而
$\E(Z_\ell Z_j)=0$。于是
\begin{align*}
 &\E\left|
 \sum_{\ell=0}^{N-1}(M_{r_{\ell+1}}-M_{r_\ell})^2
 -\int_0^TH_s^2\dd s
 \right|^2\\
 &\quad=\sum_{\ell=0}^{N-1}\E Z_\ell^2
 \le2\|H\|_\infty^4\sum_{\ell=0}^{N-1}(r_{\ell+1}-r_\ell)^2
 \le2\|H\|_\infty^4T|\pi|\longrightarrow0.
\end{align*}

对一般确定分割，至多有有限个区间跨越 $H$ 的确定跳点。若在区间 $[u,v]$ 中加入跳点
$c$，则平方和的改变量精确为
\begin{align*}
 &(M_c-M_u)^2+(M_v-M_c)^2-(M_v-M_u)^2\\
 &\qquad=-2(M_c-M_u)(M_v-M_c).
\end{align*}
当原分割的网格趋零时，$u,c,v$ 之间的距离同时趋零；$M$ 的连续性使右端几乎必然趋零。
对有限个跳点求和，得到加入全部跳点前后的差趋零。因此简单过程结论对任意网格趋零的确定分割成立。

现在取简单过程 $H^{(m)}\to H$ 于 $L^2_{\mathrm{pred}}$，并记
$M^{(m)}=\int H^{(m)}\dd B$、$R^{(m)}=M-M^{(m)}$。对任意确定分割，鞅增量正交与 \Ito{} 等距给
\[
 \E\sum_{[u,v]\in\pi}(R_v^{(m)}-R_u^{(m)})^2
 =\E\int_0^T|H_s-H_s^{(m)}|^2\dd s.
\]
所以右侧趋零时，Markov 不等式给出差过程平方增量和对所有分割一致地依概率趋零。路径上的
Cauchy--Schwarz 不等式又给
\begin{align*}
 &\left|\sum(\Delta M)^2-\sum(\Delta M^{(m)})^2\right|\\
 &\quad\le
 2\left(\sum(\Delta M^{(m)})^2\right)^{1/2}
   \left(\sum(\Delta R^{(m)})^2\right)^{1/2}
 +\sum(\Delta R^{(m)})^2.
\end{align*}
对固定 $m$，第一因子的平方和随网格变细收敛，因而构成紧族；再令 $m\to\infty$，上式右端
一致地依概率变小。目标能量也满足
\[
 \E\int_0^T|H_s^2-(H_s^{(m)})^2|\dd s
 \le\|H-H^{(m)}\|_{\mathcal H^2}
    (\|H\|_{\mathcal H^2}+\|H^{(m)}\|_{\mathcal H^2})\to0.
\]
先固定 $m$ 令分割细化，再令 $m\to\infty$，得到所需依概率收敛。

局部情形把以上论证用于 $H\1_{(0,\tau_m]}$；对任意固定有限时刻，事件
$\{\tau_m>T\}$ 的概率趋于一，故停止结论推出未停止过程的局部结论。
\end{proof}

\begin{proposition}[协变差与多维 \Ito{} 积分]\label{proposition:6-4-2-5}
沿同一分割以 polarization 定义
\[
 [M,N]_t=\frac14\bigl([M+N]_t-[M-N]_t\bigr).
\]
若 $B=(B^1,\ldots,B^r)$ 是定义~\ref{def:6-4-2-vector-brownian} 的
$\R^r$ 值 Brownian motion，且 $H,K$ 都关于它的共同滤过可预测，
\[
 M_t=\int_0^tH_s\dd B_s^i,
 \qquad N_t=\int_0^tK_s\dd B_s^j,
\]
其中两个被积过程均平方可积，则
\[
 [M,N]_t=\delta_{ij}\int_0^tH_sK_s\dd s
\]
沿任意确定网格趋零分割列依概率成立。
\end{proposition}

\begin{proof}
若 $i=j$，则 $M\pm N=\int(H\pm K)\dd B^i$。命题
\ref{proposition:6-4-2-4} 与恒等式
$(H+K)^2-(H-K)^2=4HK$ 立即给出结论。

若 $i\ne j$，先取简单被积过程。在共同确定分割上，交叉和的每项含
$\Delta B^i\Delta B^j$。把分割加细到两个简单过程的确定跳点后，记相应系数为
$H_\ell,K_\ell$。若 $\ell<m$，对 $\mathcal F_{t_m}$ 条件化可消去第 $m$ 项；对角项则由共同滤过下
向量 Brownian 增量的条件分布得到
\[
 \E\left|\sum_\ell H_\ell K_\ell
       \Delta B_\ell^i\Delta B_\ell^j\right|^2
 =\sum_\ell\E(H_\ell^2K_\ell^2)(\Delta t_\ell)^2
 \le\|H\|_\infty^2\|K\|_\infty^2\,t|\pi|.
\]
故交叉和在 $L^2$ 中趋零。

对一般被积过程，取共同简单逼近 $H^{(m)},K^{(m)}$。例如
\begin{align*}
 \left|\sum\Delta M\,\Delta N-\sum\Delta M^{(m)}\,\Delta N^{(m)}\right|
 &\le
 \left(\sum(\Delta(M-M^{(m)}))^2\right)^{1/2}
 \left(\sum(\Delta N)^2\right)^{1/2}\\
 &\quad+
 \left(\sum(\Delta M^{(m)})^2\right)^{1/2}
 \left(\sum(\Delta(N-N^{(m)}))^2\right)^{1/2}.
\end{align*}
命题~\ref{proposition:6-4-2-4} 与 \Ito{} 等距说明误差因子可先随 $m\to\infty$ 一致地依概率变小，
其余因子构成紧族；因此简单情形的极限传到一般情形。于是不同噪声分量的协变差为零。
\end{proof}

例如二维 Brownian 的离散乘积恒等式为
\[
 B_T^1B_T^2
 =\sum_kB_{t_k}^1\Delta B_k^2
  +\sum_kB_{t_k}^2\Delta B_k^1
  +\sum_k\Delta B_k^1\Delta B_k^2.
\]
最后一项依概率趋零，前两项分别趋于相应 \Ito{} 积分，所以
\[
 B_T^1B_T^2
 =\int_0^TB_s^1\dd B_s^2+\int_0^TB_s^2\dd B_s^1.
\]
同一噪声分量则会留下 $\dd t$。这正是多维 \Ito{} 公式中
$\operatorname{tr}(\sigma\sigma^TD^2f)$ 的来源。

\begin{remark}[指数鞅与终端密度]
固定 $\theta\in\R$，令
\[
 Z_t=\exp\!\left(\theta B_t-\frac12\theta^2t\right)
\]。
对 $0\le s\le t$，Brownian 增量的独立性和 Gaussian 矩母函数给出
\begin{align*}
 \E[Z_t\mid\mathcal F_s]
 &=Z_s\,\E\!\left[
   \exp\!\left(\theta(B_t-B_s)-\frac12\theta^2(t-s)\right)
   \middle|\mathcal F_s\right]\\
 &=Z_s.
\end{align*}
因而 $Z$ 是均值为 $1$ 的正鞅，$Z_T$ 可用作终端时刻的概率密度。在新测度下识别 Brownian 漂移需要
Girsanov 定理及相应的可积性条件。
\end{remark}

\begin{figure}[htbp]
\centering
\begin{tikzpicture}[node distance=10mm and 12mm]
  \draw[book line,-{Stealth[length=2mm]}] (0,0)--(5.2,0) node[right,book label]{$t$};
  \draw[BookAccent,thick]
    (0,.45)--(1,.45)--(1,1.05)--(2.2,1.05)--(2.2,.7)--(3.5,.7)--(3.5,1.3)--(4.8,1.3);
  \foreach \x in {0,1,2.2,3.5,4.8}{\draw[BookInk!45] (\x,-.12)--(\x,.12);}
  \node[book label] at (2.5,1.65) {左端已知的 $H_k$};
  \node[book strong node] (orth) at (7.1,1.0) {$H_k\Delta B_k$\\在 $L^2$ 中正交};
  \node[book warm node,below=of orth] (sum) {$\displaystyle\sum\|H_k\Delta B_k\|_2^2$
    \\$=\E\int H^2\dd t$};
  \node[book strong node,right=of orth] (qv) {$\displaystyle\sum(\Delta B)^2$
    \\$\longrightarrow T$};
  \draw[book accent arrow] (4.8,1.0)--(orth);
  \draw[book arrow] (orth)--(sum);
  \draw[book accent arrow] (orth)--(qv);
\end{tikzpicture}
\caption{可预测阶梯过程把 Brownian 增量配成正交和；平方范数形成 \Ito{} 等距，增量平方形成时间。}
\label{fig:6-4-2-isometry-qv}
\end{figure}

\begin{exercise}
\begin{enumerate}
  \item 对两个时间区间的简单可预测过程完整展开平方，标出每个交叉项所条件化的
  $\sigma$-代数，并重证 \Ito{} 等距。
  \item 由 $\int_0^TB_s\dd B_s=(B_T^2-T)/2$ 与 Gaussian 四阶矩计算
  $\E(\int_0^TB_s\dd B_s)^2$，再与 \Ito{} 等距比较。
  \item 对二进分割逐项计算 Brownian 二次变差的均值、方差和可和尾界，完成 a.s. 收敛证明。
  \item 比较左端、右端和对称端点和；若改用中点
  $B_{(t_k+t_{k+1})/2}$，先把每段拆成两个独立半增量再计算极限。
\end{enumerate}
\end{exercise}

Brownian 增量的一阶量形成随机积分，二阶量形成时间积分。普通链式法则遗漏的正是后者；下一小节
的 Taylor 展开保留这个二阶项，由此建立 SDE、生成元和热方程之间的联系。


\subsection{\Ito{} 公式、随机微分方程与热方程应用}\label{subsec:6-4-3}
\index{Ito@\Ito{} 公式}\index{随机微分方程}\index{Feynman--Kac 公式}

对有限变差路径，二阶 Taylor 项沿细分分割消失。Brownian 增量的典型大小是
$\sqrt{\Delta t}$，所以一阶随机项不能按绝对值求和，而二阶项的和却留下确定的时间量。这是
\Ito{} 公式中额外二阶导数的全部来源，也是扩散生成元与普通一阶链式法则不同的原因。

\begin{definition}[\Ito{} 过程]\label{def:6-4-3-1}
设 $B$ 是 $\R^r$ 值 Brownian motion。连续适应过程 $X$ 若可写成
\[
 X_t=X_0+\int_0^tb_s\dd s+\int_0^t\sigma_s\dd B_s,
\]
其中 $X_t\in\R^d$，$b$ 为 $\R^d$ 值渐进可测过程，$\sigma$ 为 $d\times r$ 矩阵值
可预测过程，并且对每个 $T<\infty$，
\[
 \int_0^T|b_s|\dd s<\infty,
 \qquad
 \int_0^T\|\sigma_s\|_{\mathrm{HS}}^2\dd s<\infty
 \quad\text{a.s.},
\]
则称 $X$ 为 \Ito{} 过程。
\end{definition}

这里的 Lebesgue 积分逐路径定义；矩阵值被积过程的随机积分按
命题~\ref{proposition:6-4-2-localization} 逐列积分后拼成。

\begin{definition}[随机微分方程]\label{def:6-4-3-2}
符号
\[
 \dd X_t=b(t,X_t)\dd t+\sigma(t,X_t)\dd B_t,
 \qquad X_0=\xi,
\]
表示积分方程
\[
 X_t=\xi+\int_0^tb(s,X_s)\dd s
          +\int_0^t\sigma(s,X_s)\dd B_s.
\]
在给定概率空间、滤过和 Brownian motion 上，对该 Brownian 适应并满足积分方程的连续过程
称为\emph{强解}。路径唯一性是指任意两个使用同一初值和同一 Brownian motion 的强解不可分辨，
而不只是每个固定时刻分别 a.s. 相等。
\end{definition}

\begin{definition}[扩散生成元]\label{def:6-4-3-3}
对时间齐次系数 $b:\R^d\to\R^d$、
$\sigma:\R^d\to\R^{d\times r}$ 及 $f\in C^2(\R^d)$，定义形式生成元
\[
 Lf(x)=b(x)\cdot\nabla f(x)
 +\frac12\operatorname{tr}\!\bigl(\sigma(x)\sigma(x)^TD^2f(x)\bigr).
\]
\end{definition}

\begin{theorem}[\Ito{} 公式]\label{theorem:6-4-3-1}
设 $X$ 是定义~\ref{def:6-4-3-1} 的 $\R^d$ 值连续 \Ito{} 过程，
$f\in C^{1,2}([0,\infty)\times\R^d)$。则存在一个概率一事件，使在该事件上对每个 $t\ge0$ 同时有
\begin{align*}
 f(t,X_t)=f(0,X_0)
 &+\int_0^t\left(\partial_sf(s,X_s)+b_s\cdot\nabla f(s,X_s)\right)\dd s\\
 &+\frac12\int_0^t
 \operatorname{tr}\!\left(\sigma_s\sigma_s^TD^2f(s,X_s)\right)\dd s\\
 &+\int_0^t\nabla f(s,X_s)\sigma_s\dd B_s.
\end{align*}
等价的微分记号是
\[
 \dd f(t,X_t)=
 \left(\partial_tf+b\cdot\nabla f+\tfrac12
 \operatorname{tr}(\sigma\sigma^TD^2f)\right)\dd t
 +\nabla f\,\sigma\dd B_t,
\]
其中所有导数都在 $(t,X_t)$ 评价。
\end{theorem}

\begin{proof}
先证明有界情形：固定 $T<\infty$，并假设在 $[0,T]$ 上
$|X_t|$、$\int_0^t|b_s|\dd s$ 与
$\int_0^t\|\sigma_s\|_{\mathrm{HS}}^2\dd s$ 都有确定上界。于是 $(t,X_t)$ 留在某个紧集。
证明末段将明确定义局部化停时，并用已经证明的停止积分恒等式把一般情形约化到这里。

记 $X=A+M$，其中
\[
 A_t=X_0+\int_0^tb_s\dd s,
 \qquad M_t=\int_0^t\sigma_s\dd B_s.
\]
连续有限变差过程 $A$ 沿任意网格趋零分割满足
\[
 \sum|\Delta A|^2
 \le\max|\Delta A|\operatorname{Var}(A;[0,T])\longrightarrow0.
\]
同一 Cauchy--Schwarz 估计及命题~\ref{proposition:6-4-2-4} 说明
$\sum\Delta A^i\Delta M^j\to0$ 依概率。按向量积分的定义，
\[
 M_t^i=\sum_{a=1}^r\int_0^t\sigma_{ia}(s)\dd B_s^a,
\]
而离散交叉和的有限双线性展开与命题~\ref{proposition:6-4-2-5} 给出
\[
 [M^i,M^j]_t
 =\sum_{a,b=1}^r
 \left[\int_0^\cdot\sigma_{ia}(s)\dd B_s^a,
       \int_0^\cdot\sigma_{jb}(s)\dd B_s^b\right]_t
 =\sum_{a=1}^r\int_0^t\sigma_{ia}(s)\sigma_{ja}(s)\dd s.
\]
所以
\begin{equation}\label{eq:6-4-3-ito-process-covariation}
 [X^i,X^j]_t=[M^i,M^j]_t
 =\int_0^t(\sigma_s\sigma_s^T)_{ij}\dd s.
\end{equation}

取确定分割 $\pi_n=\{0=t_0<\cdots<t_{N_n}=t\}$，$|\pi_n|\to0$。把一个增量拆成时间变化和
空间变化：
\begin{align*}
 &f(t_{k+1},X_{t_{k+1}})-f(t_k,X_{t_k})\\
 &\quad=
 \bigl[f(t_{k+1},X_{t_{k+1}})-f(t_k,X_{t_{k+1}})\bigr]
 +\bigl[f(t_k,X_{t_{k+1}})-f(t_k,X_{t_k})\bigr].
\end{align*}
由 $\partial_tf$ 在相关紧集上一致连续、$X$ 连续，第一项求和趋于
$\int_0^t\partial_sf(s,X_s)\dd s$。

对第二项作二阶空间 Taylor 展开：
\begin{align*}
 f(t_k,X_{t_{k+1}})-f(t_k,X_{t_k})
 &=\nabla f(t_k,X_{t_k})\cdot\Delta X_k\\
 &\quad+\frac12\sum_{i,j}\partial_{ij}f(t_k,X_{t_k})
      \Delta X_k^i\Delta X_k^j+R_{k,n}.
\end{align*}
$D^2f$ 在紧集上一致连续，所以存在模量 $\omega(\delta)\downarrow0$，使
\[
 |R_{k,n}|\le\omega\!\left(\max_k|\Delta X_k|\right)|\Delta X_k|^2.
\]
路径连续性给 $\max_k|\Delta X_k|\to0$ a.s.，而
$\sum_k|\Delta X_k|^2$ 由
\eqref{eq:6-4-3-ito-process-covariation} 依概率趋于
$\int_0^t\|\sigma_s\|_{\mathrm{HS}}^2\dd s$，因而构成紧族。故
$\sum_kR_{k,n}\to0$ 依概率。

一阶和中的有限变差部分为
\[
 \sum_k\nabla f(t_k,X_{t_k})\cdot\int_{t_k}^{t_{k+1}}b_s\dd s
 \longrightarrow\int_0^t\nabla f(s,X_s)\cdot b_s\dd s.
\]
这是左端阶梯函数对连续函数 $s\mapsto\nabla f(s,X_s)$ 的一致逼近，再乘可积函数 $|b_s|$
后的 DCT。一阶随机部分的被积过程
\[
 \sum_k\nabla f(t_k,X_{t_k})\sigma_s\1_{(t_k,t_{k+1}]}(s)
\]
在局部 $L^2_{\mathrm{pred}}$ 中趋于 $\nabla f(s,X_s)\sigma_s$；原因同样是左端阶梯函数一致收敛，
并由停止后的 $\int\|\sigma\|^2$ 支配。\Ito{} 等距于是把其积分和传到
$\int_0^t\nabla f(s,X_s)\sigma_s\dd B_s$。

二阶和需要一个带权的协变差版本。若有界可预测权函数 $w$ 在确定粗分割的每一块上为常数，
则把细分割加入粗分割端点后，命题~\ref{proposition:6-4-2-5} 逐块给出下式；加入有限个端点
前后的差是有限个跨端点增量的乘积，由 $M$ 的连续性依概率趋零。因此
\[
 \sum_k w_{t_k}\Delta M_k^i\Delta M_k^j
 \longrightarrow
 \int_0^t w_s(\sigma_s\sigma_s^T)_{ij}\dd s
 \quad\text{依概率}.
\]
对连续有界适应权函数 $w$，令 $w^{(m)}$ 是确定粗分割上的左端阶梯逼近。路径连续性给
$\|w^{(m)}-w\|_{\infty,[0,t]}\to0$ a.s.，而对任意细分割
\[
 \left|\sum_k(w_{t_k}-w^{(m)}_{t_k})\Delta M_k^i\Delta M_k^j\right|
 \le \|w-w^{(m)}\|_{\infty,[0,t]}
 \left(\sum_k|\Delta M_k^i|^2\right)^{1/2}
 \left(\sum_k|\Delta M_k^j|^2\right)^{1/2}.
\]
两个平方增量和由命题~\ref{proposition:6-4-2-4} 构成紧族；相应时间积分之差则由
\[
 \left|\int_0^t(w_s-w_s^{(m)})(\sigma_s\sigma_s^T)_{ij}\dd s\right|
 \le \|w-w^{(m)}\|_{\infty,[0,t]}
      \int_0^t\|\sigma_s\|_{\mathrm{HS}}^2\dd s
\]
趋于零。因此先令细分割网格趋零、再令 $m\to\infty$，带权结论对
$w_s=\partial_{ij}f(s,X_s)$ 成立。
含有限变差增量的带权项仍然消失：停止后权函数由常数 $C$ 控制，而
\begin{align*}
 \left|\sum_kw_k\Delta A_k^i\Delta A_k^j\right|
 &\le C\left(\sum_k|\Delta A_k^i|^2\right)^{1/2}
          \left(\sum_k|\Delta A_k^j|^2\right)^{1/2},\\
 \left|\sum_kw_k\Delta A_k^i\Delta M_k^j\right|
 &\le C\left(\sum_k|\Delta A_k^i|^2\right)^{1/2}
          \left(\sum_k|\Delta M_k^j|^2\right)^{1/2}.
\end{align*}
第一类平方和趋零，$\sum|\Delta M^j|^2$ 构成紧族，故两行右端都依概率趋零。因此
\[
 \sum_k\partial_{ij}f(t_k,X_{t_k})\Delta X_k^i\Delta X_k^j
 \longrightarrow
 \int_0^t\partial_{ij}f(s,X_s)(\sigma_s\sigma_s^T)_{ij}\dd s
\]
依概率成立。对有限个 $i,j$ 求和，正是迹项。

上述有限项之和依概率收敛到陈述中的右端；而对每个 $n$，Taylor 恒等式求和后的该和恒等于固定的
$f(t,X_t)-f(0,X_0)$。常值随机变量列若依概率收敛到另一个随机变量，则二者 a.s. 相等：对任意
$\varepsilon>0$，二者之差大于 $\varepsilon$ 的概率必为零。因此停止后的 \Ito{} 公式成立。

最后令
\[
 \rho_m=\inf\left\{t:|X_t|+\int_0^t|b_s|\dd s
 +\int_0^t\|\sigma_s\|_{\mathrm{HS}}^2\dd s\ge m\right\}\wedge m.
\]
大括号中的过程连续、适应且非负，因此
$\{\rho_m\le t\}$ 可由它在 $[0,t]$ 上的有理时刻上确界写出，故 $\rho_m$ 是停时。
阈值与末端截断都随 $m$ 增加，所以 $(\rho_m)$ 非降；连续性和局部可积性又给
$\rho_m\uparrow\infty$ a.s.。把前面的证明用于停止后的漂移
$b_s\1_{\{s\le\rho_m\}}$、扩散系数
$\sigma_s\1_{\{s\le\rho_m\}}$ 以及一个覆盖
$[0,m]\times\overline B(0,m)$ 的紧集，得到公式到 $t\wedge\rho_m$。对每个固定 $t$，在
$\{\rho_m>t\}$ 上这就是未停止公式；这些事件递增到概率一。故一般 $C^{1,2}$ 与局部可积系数的
结论对每个固定 $t$ 成立。再对可数集 $t\in\Q_+$ 取零集之并。等式左端连续；右端的两个
Lebesgue 积分由局部可积性对上限连续，随机积分取定理~\ref{theorem:6-4-2-2} 与命题
\ref{proposition:6-4-2-localization} 构造的连续版本。因此有理时刻的等式由连续性推广到所有实时刻，得到
定理声明中的共同概率一事件。
\end{proof}

\begin{example}[平方与指数]\label{ex:6-4-3-1}
对 $f(x)=x^2$，\Ito{} 公式给出
\[
 B_t^2=2\int_0^tB_s\dd B_s+t,
\]
故 $B_t^2-t$ 是平方可积鞅。这与例~\ref{ex:6-4-2-2} 的左端和计算完全一致。

固定 $\theta\in\R$ 并令
\[
 Z_t=\exp\!\left(\theta B_t-\frac12\theta^2t\right).
\]
对 $f(t,x)=e^{\theta x-\theta^2t/2}$ 应用 \Ito{} 公式，漂移项
$(-\theta^2/2+\theta^2/2)f$ 抵消，得到
\[
 Z_t=1+\theta\int_0^tZ_s\dd B_s.
\]
又由 Gaussian 矩母函数
$\E Z_s^2=e^{\theta^2s}$，所以
$\E\int_0^T\theta^2Z_s^2\dd s<\infty$。随机积分是真平方可积鞅，故 $(Z_t)$ 是均值为一的
指数鞅，而不只是局部鞅。
\end{example}

\begin{example}[几何 Brownian motion]\label{ex:6-4-3-2}
设 $X_0$ 为给定初值，$\mu,\sigma\in\R$。令
\[
 X_t=X_0\exp\!\left((\mu-\tfrac12\sigma^2)t+\sigma B_t\right).
\]
对指数函数应用 \Ito{} 公式，得到
\[
 \dd X_t=\mu X_t\dd t+\sigma X_t\dd B_t.
\]
反过来，若 $X_0>0$ 且 $X$ 是该方程的正解，对 $\log X_t$ 局部化应用 \Ito{} 公式，
\[
 \dd\log X_t=(\mu-\tfrac12\sigma^2)\dd t+\sigma\dd B_t,
\]
积分后恢复同一表达式。修正项 $-\sigma^2/2$ 来自二次变差；按普通链式法则会漏掉它。
\end{example}

\begin{example}[Ornstein--Uhlenbeck SDE 的显式解]
\label{ex:6-4-3-ou-sde}
设 $\theta,\eta>0$，$X_0\in L^2(\mathcal F_0)$，且 $B$ 相对于给定滤过为 Brownian motion；因而 $X_0$ 与 Brownian 的未来增量独立。考虑
\[
 \dd X_t=-\theta X_t\dd t+\eta\dd B_t.
\]
对 $Y_t=e^{\theta t}X_t$ 使用 \Ito{} 公式；由于确定因子 $e^{\theta t}$ 有零二次变差，
\[
 \dd Y_t=\eta e^{\theta t}\dd B_t.
\]
因而
\begin{equation}\label{eq:6-4-3-ou-solution}
 X_t=e^{-\theta t}X_0
 +\eta\int_0^te^{-\theta(t-s)}\dd B_s.
\end{equation}
右端的确定被积函数属于 $L^2[0,t]$，所以式中 \Ito{} 积分已由
例~\ref{ex:6-4-2-1} 定义。这个随机积分是中心 Gaussian，方差为
\[
 \eta^2\int_0^te^{-2\theta(t-s)}\dd s
 =\frac{\eta^2}{2\theta}(1-e^{-2\theta t}).
\]
特别地，若 $X_0\sim N(0,\eta^2/(2\theta))$ 并与 $B$ 独立，则式
\eqref{eq:6-4-3-ou-solution} 表明每个有限向量
$(X_{t_1},\ldots,X_{t_n})$ 都是 $X_0$ 与有限个确定 \Ito{} 积分组成的联合 Gaussian 向量的线性变换。
此外，对 $s\le t$，
\[
 X_t=e^{-\theta(t-s)}X_s
 +\eta\int_s^te^{-\theta(t-r)}\dd B_r,
\]
而最后的积分独立于题设给定的 $\mathcal F_s$ 且均值为零。因此
$\Var(X_t)=\eta^2/(2\theta)$，并且
\[
 \Cov(X_s,X_t)=e^{-\theta(t-s)}\Var(X_s)
 =\frac{\eta^2}{2\theta}e^{-\theta(t-s)}.
\]
因此这个强解与例~\ref{ex:6-4-1-ou-gaussian} 的平稳 Gaussian 构造具有相同有限维分布。
若 $X_0=x$ 为确定值，则
\[
 \E X_t=e^{-\theta t}x,\qquad
 \Var(X_t)=\frac{\eta^2}{2\theta}(1-e^{-2\theta t}),
\]
这显示该分布向平稳 Gaussian 分布回复。
\end{example}

定义~\ref{def:6-4-3-3} 的生成元现在有了严格含义。对时间齐次 \Ito{} 过程，公式给出
\[
 f(X_t)-f(X_0)-\int_0^tLf(X_s)\dd s
 =\int_0^t\nabla f(X_s)\sigma(X_s)\dd B_s.
\]
右端总是局部鞅；若例如
$\E\int_0^T\|\sigma(X_s)^T\nabla f(X_s)\|^2\dd s<\infty$，则它在 $[0,T]$ 上是真鞅。
因而，要把右端的局部鞅提升为真鞅，还需要相应的可积性。

\begin{theorem}[Lipschitz SDE 的强存在唯一性]\label{theorem:6-4-3-2}
设 Borel 函数
$b:[0,\infty)\times\R^d\to\R^d$、
$\sigma:[0,\infty)\times\R^d\to\R^{d\times r}$ 在每个有限时域 $[0,T]$ 上满足
\[
 |b(t,x)-b(t,y)|+\|\sigma(t,x)-\sigma(t,y)\|_{\mathrm{HS}}
 \le L_T|x-y|,
\]
以及
\[
 |b(t,x)|^2+\|\sigma(t,x)\|_{\mathrm{HS}}^2
 \le C_T(1+|x|^2).
\]
若 $X_0\in L^2$ 对 $\mathcal F_0$ 可测，$B$ 相对于给定的满足通常条件的滤过是 Brownian motion，
则 SDE 有路径唯一的连续适应强解，并且
\[
 \E\sup_{t\le T}|X_t|^2
 \le C_T'(1+\E|X_0|^2).
\]
\end{theorem}

\begin{proof}
令 $S_T^2=L^2(\Omega;C([0,T];\R^d))$，范数为
$\|X\|_{S_T^2}^2=\E\sup_{t\le T}|X_t|^2$，并只取连续适应元素。它是一个闭子空间，
因而完备：$C([0,T];\R^d)$-值 $L^2$ 极限仍有连续路径，而对每个固定 $t$，评价映射连续且
$L^2(\mathcal F_t)$ 是 $L^2(\mathcal F)$ 的闭子空间，所以适应性也在极限下保持。对
$X\in S_T^2$ 定义
\[
 (\Phi X)_t=X_0+\int_0^tb(s,X_s)\dd s
 +\int_0^t\sigma(s,X_s)\dd B_s.
\]
连续适应的 $X$ 渐进可测，Borel 复合仍渐进可测。随机系数的可预测代表由
评注~\ref{remark:6-4-2-progressive-predictable} 给出；线性增长保证所需 $L^2$ 条件。因此
$\Phi$ 定义良好。

对 $0<\delta\le T$，Cauchy--Schwarz、连续 Doob $L^2$ 不等式
\eqref{eq:6-4-2-continuous-doob} 和 \Ito{} 等距给出
\begin{align*}
 \|\Phi X-\Phi Y\|_{S_\delta^2}^2
 &\le2\delta\E\int_0^\delta|b(s,X_s)-b(s,Y_s)|^2\dd s\\
 &\quad+8\E\int_0^\delta
 \|\sigma(s,X_s)-\sigma(s,Y_s)\|_{\mathrm{HS}}^2\dd s\\
 &\le(2L_T^2\delta^2+8L_T^2\delta)\|X-Y\|_{S_\delta^2}^2.
\end{align*}
取 $\delta$ 足够小，使括号中的常数小于一，$\Phi$ 是压缩映射。线性增长与同一估计给
\[
 \|\Phi X\|_{S_\delta^2}^2
 \le C\bigl(1+\E|X_0|^2+\delta\|X\|_{S_\delta^2}^2\bigr)<\infty,
\]
故 $\Phi$ 确实把该空间映入自身。Banach 不动点定理产生 $[0,\delta]$ 上唯一强解。

以 $X_\delta$ 为新初值，在 $[\delta,2\delta]$ 上把积分下限改为 $\delta$ 并重复同一构造；
$X_\delta\in L^2(\mathcal F_\delta)$，而 $(B_{\delta+t}-B_\delta)_{t\ge0}$ 相对于移位滤过仍是
Brownian motion。有限次拼接覆盖固定的 $[0,T]$，各段端点相同且路径连续。再在
$[T,T+1]$、$[T+1,T+2]$ 等有限时段递归使用相应的局部常数，便得到 $[0,\infty)$ 上的强解。

为给出不依赖分段、也不预先假设任意候选解属于 $S_T^2$ 的路径唯一性，设 $X,Y$ 是两个连续强解，
并对整数 $n\ge1$ 令
\[
 \tau_n=\inf\{t\in[0,T]:|X_t|+|Y_t|\ge n\}\wedge T,
\]
其中空集的下确界取 $T$。当 $t<T$ 时，路径连续性给
\[
 \{\tau_n\le t\}
 =\left\{\sup_{q\in\Q\cap[0,t]}(|X_q|+|Y_q|)\ge n\right\}\in\mathcal F_t,
\]
而 $t\ge T$ 时该事件为 $\Omega$，所以 $\tau_n$ 是停时。在 $\tau_n$ 以前两个解的系数
平方可积；由随机积分停止恒等式，
\begin{align*}
 X_{t\wedge\tau_n}-Y_{t\wedge\tau_n}
 &=\int_0^t\1_{\{s\le\tau_n\}}
   \bigl(b(s,X_s)-b(s,Y_s)\bigr)\dd s\\
 &\quad+\int_0^t\1_{\{s\le\tau_n\}}
   \bigl(\sigma(s,X_s)-\sigma(s,Y_s)\bigr)\dd B_s .
\end{align*}
记
\[
 F_n(t)=\E\sup_{r\le t}|X_{r\wedge\tau_n}-Y_{r\wedge\tau_n}|^2.
\]
停止后各项可积，Cauchy--Schwarz、Doob--\Ito{} 估计与 Lipschitz 条件逐项给出
\[
 F_n(t)\le C_T\int_0^tF_n(s)\dd s.
\]
令 $G_n(t)=\int_0^tF_n(s)\dd s$，则 $G_n'(t)\le C_TG_n(t)$ a.e. 且 $G_n(0)=0$；
对 $e^{-C_Tt}G_n(t)$ 积分得到 $G_n=0$，继而 $F_n=0$。对每个样本点，$X,Y$ 在
$[0,T]$ 上连续有界，故 $\tau_n=T$ 对所有充分大的 $n$ 成立。对可数个 $n$ 的概率一事件取交后，
得到 $\sup_{t\le T}|X_t-Y_t|=0$ a.s.；$T$ 任意整数，故两解不可分辨。

最后对单个解使用
$|x+y+z|^2\le3(|x|^2+|y|^2+|z|^2)$、Cauchy--Schwarz、Doob--\Ito{} 估计和线性增长，得到
\[
 F_X(t):=\E\sup_{s\le t}|X_s|^2
 \le C_T(1+\E|X_0|^2)+C_T\int_0^tF_X(s)\dd s.
\]
同一指数乘子论证给出所述矩界。
\end{proof}

\begin{example}[缺少 Lipschitz／增长时的两种失败]\label{ex:6-4-3-4}
先令噪声系数为零。方程
\[
 \dd X_t=2\sqrt{X_t^+}\dd t,
 \qquad X_0=0
\]
有零解；对每个 $c\ge0$，还有等待解
\[
 X_t^{(c)}=
 \begin{cases}
 0,&t\le c,\\
 (t-c)^2,&t>c.
 \end{cases}
\]
它在 $t>c$ 时满足
$(X_t^{(c)})'=2(t-c)=2\sqrt{X_t^{(c)}}$，在 $t=c$ 两侧导数都为零。
系数具有线性增长，却在原点不 Lipschitz，所以路径唯一性失败。

另一方面，
\[
 \dd X_t=X_t^2\dd t,
 \qquad X_0=x_0>0
\]
的显式解是
$X_t=x_0/(1-x_0t)$，只存在到 $t=1/x_0$。系数局部 Lipschitz，却没有全局线性增长；局部唯一性
没有阻止爆炸。两个例子分别只撤去一项核心假设。
\end{example}

\begin{proposition}[Dynkin 公式]\label{proposition:6-4-3-3}
设 $X$ 是从 $X_0=x$ 出发、在全部有限时域上有定义的时间齐次 SDE
\[
 \dd X_s=b(X_s)\dd s+\sigma(X_s)\dd B_s
\]
的强解，$L$ 如定义~\ref{def:6-4-3-3}，$f\in C^2(\R^d)$，$\tau$ 为停时。若
\[
 \E_x\int_0^{t\wedge\tau}|Lf(X_s)|\dd s<\infty
\]
且
\[
 N_s=\int_0^{s\wedge\tau}\nabla f(X_r)\sigma(X_r)\dd B_r,
 \qquad 0\le s\le t,
\]
是真鞅，则
\[
 \E_x f(X_{t\wedge\tau})
 =f(x)+\E_x\int_0^{t\wedge\tau}Lf(X_s)\dd s.
\]
\end{proposition}

\begin{proof}
把 \Ito{} 公式应用到停止过程，得到
\[
 f(X_{t\wedge\tau})=f(x)+\int_0^{t\wedge\tau}Lf(X_s)\dd s+N_t.
\]
第一项的绝对可积性由假设给出，停止随机积分 $N$ 是真鞅且从零出发，故
$\E N_t=0$。对等式取期望即得结论。若这些条件只在紧集退出前成立，应先取
$\tau\wedge\tau_m$ 局部化，再用 DCT、MCT 或 UI 中适用的一项放开 $m$。因此随机积分项取期望为零的根据是真鞅性，
而非局部鞅性。
\end{proof}

\begin{example}[调和函数与退出分布]\label{ex:6-4-3-3}
令 $W$ 是 $\R^d$ 值 Brownian motion，$D\subset\R^d$ 为有界开集，固定 $x\in D$，并令
$B_t^x=x+W_t$ 及
\[
 \tau_D=\inf\{t\ge0:B_t^x\notin D\}.
\]
这是停时。事实上，连续性与 $[0,t]$ 的紧致性给出
\[
 \{\tau_D\le t\}
 =\left\{\inf_{q\in\Q\cap[0,t]}
   \operatorname{dist}(B_q^x,D^c)=0\right\}\in\mathcal F_t^B.
\]
取嵌套开集耗尽
\[
 D_m=\{y\in D:|y|<m,\ \operatorname{dist}(y,D^c)>1/m\},
 \qquad D_m\uparrow D,
\]
则 $\overline{D_m}\subset D$ 为紧集。相应退出时刻
$\tau_{D_m}$ 递增到 $\tau_D$：若极限严格小于 $\tau_D$，则路径在一个稍长的紧时间段上的像是
$D$ 的紧子集，因而最终全部落在某个 $D_m$ 中，这与退出时刻的定义矛盾。

若 $u\in C^2(D)\cap C(\overline D)$ 且 $\Delta u=0$，在 $\tau_{D_m}$ 前应用 \Ito{} 公式。
$\nabla u$ 在 $\overline{D_m}$ 上有界，所以随机积分平方可积，得到
$u(B_{t\wedge\tau_{D_m}}^x)$ 是真鞅。令 $m\to\infty$，路径连续性与
$|u|\le\|u\|_{C(\overline D)}$ 给
\[
 u(B_{t\wedge\tau_{D_m}}^x)\longrightarrow u(B_{t\wedge\tau_D}^x)
 \quad\text{a.s. 且于 }L^1.
\]
固定 $0\le s\le t$。对每个 $m$ 的鞅等式取条件期望，再用条件期望的 $L^1$ 收缩性把上式传到
极限，得到
\[
 \E_x[u(B_{t\wedge\tau_D}^x)\mid\mathcal F_s^B]
 =u(B_{s\wedge\tau_D}^x).
\]
因此 $u(B_{t\wedge\tau_D}^x)$ 是有界真鞅，而不只是逐时刻期望相同的过程。

取 $R$ 使 $D\subset[-R,R]^d$。由于
\[
 \Prob_x(\tau_D>n)
 \le\Prob(|x_1+W_n^1|\le R)
 \le\frac{2R}{\sqrt{2\pi n}}\longrightarrow0,
\]
所以 $\tau_D<\infty$ a.s.。连续路径在退出时落在 $\partial D$。连续时间可选抽样命题
\ref{proposition:6-4-1-continuous-optional} 的有界情形可用于
$t\wedge\tau_D$；再由 DCT 放开 $t$，给出
\[
 u(x)=\E_xu(B_{\tau_D}^x).
\]
因此边界值通过 Brownian 退出分布恢复。函数 $u$ 的有界性使停止鞅族一致可积，从而保证放开时间截断后期望恒等式仍成立。

在一维区间 $D=(a,b)$ 中取调和函数 $u(x)=(x-a)/(b-a)$，便得到
\[
 \Prob_x(B_{\tau_D}^x=b)=\frac{x-a}{b-a},
 \qquad a<x<b.
\]
\end{example}

\begin{theorem}[热方程与 Feynman--Kac 基础式]\label{theorem:6-4-3-4}
令 $B=(B^1,\ldots,B^d)$ 的各坐标是独立的一维 Brownian motion，并定义
\[
 P_tg(x)=\E[g(x+B_t)],\qquad t\ge0.
\]
若 $g\in C_b^2(\R^d)$ 且其一、二阶导数有界，则
$u(t,x)=P_tg(x)$ 满足
\[
 \partial_tu=\frac12\Delta u,\qquad u(0,x)=g(x).
\]
对一般 $g\in C_b(\R^d)$，$P_tg$ 在 $t>0$ 时仍是光滑经典解，且
$P_tg(x)\to g(x)$ 对每个 $x$ 成立；若 $g\in C_0(\R^d)$，则还成立
\[
 \|P_tg-g\|_\infty\longrightarrow0\qquad(t\downarrow0).
\]
一般 $C_b$ 数据不保证这一范数收敛。

进一步，设 $V\in C_b(\R^d)$、$g\in C_0(\R^d)$。令从 $x$ 出发的 Brownian motion 仍记为
$B$，并置
\begin{equation}\label{eq:6-4-3-fk-expectation}
 u(t,x)=\E_x\!\left[
 \exp\!\left(-\int_0^tV(B_s)\dd s\right)g(B_t)\right].
\end{equation}
则 $u\in C([0,\infty);C_0(\R^d))$，并且它是 Volterra 方程
\begin{equation}\label{eq:6-4-3-fk-duhamel}
 u(t)=P_tg-\int_0^tP_{t-s}\bigl(Vu(s)\bigr)\dd s
\end{equation}
在 $C([0,\infty);C_0(\R^d))$ 中的唯一温和解（mild solution）。若仅假设
$g\in C_b(\R^d)$，式~\eqref{eq:6-4-3-fk-expectation} 与
\eqref{eq:6-4-3-fk-duhamel} 仍逐点成立，并且 $u(t,x)\to g(x)$；这里不声称
$C_b$ 范数下的强连续性。

若进一步 $g,V\in C_b^2(\R^d)$ 且它们的一、二阶导数均有界，则
\eqref{eq:6-4-3-fk-expectation} 是经典解
\[
 \partial_tu=\frac12\Delta u-Vu,\qquad u(0,x)=g(x).
\]
\end{theorem}

\begin{proof}
记
\[
 q_t(z)=(2\pi t)^{-d/2}\exp\!\left(-\frac{|z|^2}{2t}\right),
 \qquad t>0.
\]
独立坐标的联合密度给出
\begin{equation}\label{eq:6-4-3-brownian-heat-kernel}
 P_tg(x)=\int_{\R^d}q_t(z)g(x+z)\dd z
       =\int_{\R^d}q_t(x-y)g(y)\dd y.
\end{equation}
直接求导得到
\[
 \partial_tq_t(z)
 =q_t(z)\left(-\frac d{2t}+\frac{|z|^2}{2t^2}\right)
 =\frac12\Delta q_t(z).
\]
对每个固定 $t>0$，$q_t$ 的任意阶空间导数以及 $\partial_tq_t$ 都属于 $L^1(\R^d)$。
因此有界的 $g$ 可在式~\eqref{eq:6-4-3-brownian-heat-kernel} 中与核求导交换；这说明
$P_tg$ 在 $t>0$ 光滑，并满足 $\partial_tP_tg=\frac12\Delta P_tg$。

写 $B_t\overset{d}=\sqrt t\,Z$，其中 $Z$ 是标准 $d$ 维 Gaussian 向量。对固定 $x$，
$g(x+\sqrt t\,Z)\to g(x)$ a.s.，并且绝对值不超过 $\|g\|_\infty$，故 DCT 给出
$P_tg(x)\to g(x)$。若 $g\in C_0$，则 $(q_t)_{t>0}$ 是近似恒等，定理
\ref{theorem:4-4-2-1} 给出一致范数收敛；同一个卷积表示也说明 $P_tg\in C_0$。

当 $g\in C_b^2$ 且所列导数有界时，还可把差商移到 $g$ 上，得到
\[
 \partial_iP_tg=P_t(\partial_i g),\qquad
 \partial_{ij}P_tg=P_t(\partial_{ij}g).
\]
对 $g(x+B_s)$ 应用 \Ito{} 公式。由于 $\nabla g$ 有界，
$\int_0^t\nabla g(x+B_s)\dd B_s$ 是平方可积真鞅，故
\begin{equation}\label{eq:6-4-3-heat-dynkin}
 P_tg(x)-g(x)=\frac12\int_0^tP_s(\Delta g)(x)\dd s.
\end{equation}
在式~\eqref{eq:6-4-3-heat-dynkin} 中令 $t$ 作差商并用 DCT，便得到零时刻的右导数；
正时间的结论也已由核求导得到。这完成热方程的第一层结论。

范数收敛不能从 $C_0$ 无条件扩大到 $C_b$。事实上，对 $t>0$ 和 $h\in\R^d$，
\[
 \|P_tg(\,\cdot+h)-P_tg\|_\infty
 \le\|g\|_\infty\|q_t(\,\cdot-h)-q_t\|_1.
\]
$L^1$ 平移连续性说明 $P_tg$ 一致连续。若某个 $t_n\downarrow0$ 满足
$\|P_{t_n}g-g\|_\infty\to0$，则 $g$ 是一致连续函数的一致极限，因而一致连续。但在一维中
$g(x)=\sin(x^2)$ 不一致连续：令
\[
 x_n=\sqrt{2\pi n+\frac\pi2},\qquad
 y_n=\sqrt{2\pi n+\frac{3\pi}2},
\]
则 $|x_n-y_n|\to0$，而 $g(x_n)=1$、$g(y_n)=-1$。所以一般 $C_b$ 上不存在所述强连续性。

下面证明 Feynman--Kac 部分。记 $K=\|V\|_\infty$。对每条连续路径，微积分基本定理给出
\begin{equation}\label{eq:6-4-3-path-exponential-identity}
 \exp\!\left(-\int_0^tV(B_r)\dd r\right)
 =1-\int_0^tV(B_s)
 \exp\!\left(-\int_s^tV(B_r)\dd r\right)\dd s.
\end{equation}
两边乘以 $g(B_t)$。被积函数的绝对值不超过
$K e^{Kt}\|g\|_\infty$，故 Fubini 可用。再对 $\mathcal F_s$ 条件化；Brownian 的 Markov 性给
\[
 \E_x\!\left[
 e^{-\int_s^tV(B_r)\dd r}g(B_t)\,\middle|\,\mathcal F_s
 \right]
 =u(t-s,B_s).
\]
确切地，置 $W_q=B_{s+q}-B_s$（$0\le q\le t-s$）。向量 Brownian 的独立增量先说明
$W$ 的任意有限个有理时刻评价与 $\mathcal F_s$ 独立；连续路径空间的 Borel
$\sigma$-代数由这些评价生成，函数单调类定理因而说明整个连续路径 $W$ 独立于
$\mathcal F_s$ 且具有 $d$ 维 Wiener 分布。路径泛函
\[
 \Psi_y(w)=
 \exp\!\left(-\int_0^{t-s}V(y+w_q)\dd q\right)g(y+w_{t-s})
\]
关于 $(y,w)\in\R^d\times C_0([0,t-s];\R^d)$ Borel 可测：时间积分是连续
Riemann 和的逐点极限。它还满足
$|\Psi_y(w)|\le e^{K(t-s)}\|g\|_\infty$。独立随机元的条件积分公式于是给
\[
 \E_x[\Psi_{B_s}(W)\mid\mathcal F_s]
 =\int\Psi_{B_s}(w)\,\mathsf W(\dd w)
 =u(t-s,B_s),
\]
其中 $\mathsf W$ 是相应的 Wiener 测度。这就严格给出上式。
于是
\[
 u(t)=P_tg-\int_0^tP_s\bigl(Vu(t-s)\bigr)\dd s.
\]
代换 $r=t-s$ 就得到式~\eqref{eq:6-4-3-fk-duhamel}。因此概率表示在不假定时间导数或二阶空间导数存在时，
已经给出温和方程。

现在假设 $g\in C_0$。在任意固定 $T>0$ 上，从 $u_0(t)=P_tg$ 开始递归定义
\[
 u_{n+1}(t)=-\int_0^tP_{t-s}\bigl(Vu_n(s)\bigr)\dd s.
\]
热半群把 $C_0$ 映入自身并在其上强连续，且乘法算子
$f\mapsto Vf$ 把 $C_0$ 映入自身。归纳可知
$u_n\in C([0,T];C_0)$。又由 $\|P_t\|_{\infty\to\infty}\le1$，
\begin{equation}\label{eq:6-4-3-dyson-bound}
 \|u_n(t)\|_\infty
 \le\|g\|_\infty\frac{(Kt)^n}{n!},
 \qquad 0\le t\le T.
\end{equation}
因此 $\sum_{n\ge0}u_n$ 在 $C([0,T];C_0)$ 中绝对一致收敛；逐项代入可知其和满足
\eqref{eq:6-4-3-fk-duhamel}。

若 $v,w\in C([0,T];C_0)$ 都满足该方程，令
$D(t)=\sup_{s\le t}\|v(s)-w(s)\|_\infty$，则
\[
 D(t)\le K\int_0^tD(s)\dd s.
\]
对右端反复迭代，得到
$D(t)\le D(T)(Kt)^n/n!$ 对每个 $n$ 成立；令 $n\to\infty$ 得 $D(t)=0$。
故温和解唯一。同一迭代也适用于任意两个在 $[0,T]\times\R^d$ 上一致有界的逐点解，
因为此时 $D(T)<\infty$，而论证没有使用范数连续性。概率表示已经满足 Duhamel 方程，
所以等于上述级数，因而属于
$C([0,T];C_0)$。$T$ 任意，结论在全部有限时域成立。若 $g$ 只属于 $C_b$，前面的条件期望计算
仍成立；概率表示关于 $x$ 连续且有界，而初值的逐点收敛由
\[
 |u(t,x)-g(x)|
 \le e^{Kt}P_t\!\left(|g-g(x)|\right)(x)
    +(e^{Kt}-1)|g(x)|
\]
和 $P_t(|g-g(x)|)(x)\to0$ 得到。

最后假设 $g,V\in C_b^2$ 且所列导数有界。写
\[
 A_t(x)=\int_0^tV(x+B_s)\dd s,\quad
 A_{t,i}(x)=\int_0^t\partial_iV(x+B_s)\dd s,\quad
 A_{t,ij}(x)=\int_0^t\partial_{ij}V(x+B_s)\dd s.
\]
这些量在有限时域上一致有界。对式~\eqref{eq:6-4-3-fk-expectation} 的空间差商使用 DCT，
得到
\begin{align*}
 \partial_i u(t,x)
 &=\E\!\left[e^{-A_t(x)}
   \bigl(\partial_i g(x+B_t)-g(x+B_t)A_{t,i}(x)\bigr)\right],\\
 \partial_{ij}u(t,x)
 &=\E\!\left[e^{-A_t(x)}
 \begin{aligned}[t]
 \bigl(&\partial_{ij}g(x+B_t)
       -\partial_i g(x+B_t)A_{t,j}(x)
       -\partial_j g(x+B_t)A_{t,i}(x)\\
      &-g(x+B_t)A_{t,ij}(x)
       +g(x+B_t)A_{t,i}(x)A_{t,j}(x)\bigr)
 \end{aligned}\right].
\end{align*}
这些公式及 DCT 表明 $u(t,\cdot)\in C_b^2$，并且 $u$ 及其前两阶空间导数在
$(t,x)$ 上连续。

令 $h(t)=Vu(t)$。乘积法则说明 $h(t,\cdot)\in C_b^2$，其前两阶导数在每个有限时域上一致有界。
因此可在 Duhamel 方程中把空间导数移过 $P_{t-s}$，并写成
\[
 \Delta P_{t-s}h(s)=P_{t-s}\Delta h(s).
\]
对时间求导时，积分上端贡献 $h(t)$，其余部分由上述一致界支配。于是
\begin{align*}
 \partial_tu(t)
 &=\frac12\Delta P_tg-h(t)
   -\frac12\int_0^tP_{t-s}\Delta h(s)\dd s\\
 &=\frac12\Delta u(t)-V u(t).
\end{align*}
初值由前面已经证明的逐点收敛给出。由此完成经典层级的升级，也完成全部结论。
\end{proof}

\begin{remark}[复指数、Fourier 与生成元]\label{remark:6-4-3-fourier}
固定 $\xi\in\R^d$。分别把 \Ito{} 公式用于
$\cos(\xi\cdot x)$ 与 $\sin(\xi\cdot x)$，再把两条实公式合写，得到
\[
 P_tf_\xi(x)=e^{-t|\xi|^2/2}f_\xi(x),
 \qquad
 Lf_\xi=-\frac12|\xi|^2f_\xi,
 \qquad f_\xi(x)=e^{i\xi\cdot x}.
\]
这里的 $L$ 是 Brownian 生成元 $\frac12\Delta$。同一个符号可由鞅再核对一次：对实部和虚部分别
应用 \Ito{} 公式，
\[
 M_t=\exp\!\left(i\xi\cdot B_t+\frac12|\xi|^2t\right)
\]
的漂移项为零。对每个 $T<\infty$，$|M_t|\le e^{|\xi|^2T/2}$，所以它在 $[0,T]$ 上是真鞅，
不是只由形式微分得到的局部鞅。取期望便有
$\E e^{i\xi\cdot B_t}=e^{-t|\xi|^2/2}$。

这项计算说明 Brownian 平均在指数测试函数上乘以 Gaussian 因子，也就是热半群的 Fourier 乘子关系。
式中的复指数可通过实部和虚部分别验证。
\end{remark}

\begin{figure}[htbp]
\centering
\resizebox{0.98\textwidth}{!}{%
\begin{tikzpicture}[node distance=10mm and 14mm]
  \node[book strong node] (tay) {二阶 Taylor\\沿分割求和};
  \node[book warm node,above right=of tay] (first) {一阶随机项\\$\nabla f\,\sigma\Delta B$};
  \node[book warm node,below right=of tay] (second) {二阶随机项\\
    $\frac12D^2f:(\Delta B)(\Delta B)^T$};
  \node[book strong node,right=of first] (mart) {\Ito{} 积分\\局部鞅};
  \node[book strong node,right=of second] (time) {协变差\\时间积分};
  \node[book warm node,right=22mm of mart,yshift=-10mm] (gen) {生成元
    $L$\\Dynkin 公式};
  \node[book strong node,right=of gen] (pde) {热方程\\Feynman--Kac};
  \draw[book accent arrow] (tay)--(first);
  \draw[book accent arrow] (tay)--(second);
  \draw[book arrow] (first)--(mart);
  \draw[book arrow] (second)--(time);
  \draw[book accent arrow] (mart)--(gen);
  \draw[book accent arrow] (time)--(gen);
  \draw[book accent arrow] (gen)--(pde);
\end{tikzpicture}
}
\caption{一阶 Brownian 增量进入随机积分，二阶增量经协变差留下时间项；二者在生成元中汇合，
再通向 Dynkin 公式与抛物方程。高阶 Taylor 余项的消失由 \Ito{} 公式证明中的局部估计保证。}
\label{fig:6-4-3-ito-tree}
\end{figure}

\begin{remark}[主要结论的适用范围]\label{remark:6-4-3-boundaries}
符号 $\dd X_t=b\,\dd t+\sigma\,\dd B_t$ 始终表示定义~\ref{def:6-4-3-2} 的积分方程，不能按普通
微分代数约分。局部 Lipschitz 系数只在爆炸前提供局部解；例~\ref{ex:6-4-3-4} 说明增长条件为何
另有职责。停止 \Ito{} 公式只有在随机积分升级为真鞅后才能取期望，命题
\ref{proposition:6-4-3-3} 已把这一条件写进陈述。最后，Feynman--Kac 的概率表示、$C_0$ 温和解
与经典 PDE 解是三个不同层级；定理~\ref{theorem:6-4-3-4} 分别给出从概率表示到后两个层级所需的假设。
\end{remark}

\begin{exercise}
\begin{enumerate}
  \item 不直接引用例~\ref{ex:6-4-3-1}，分别用 \Ito{} 等距与 Gaussian 矩母函数证明
  $B_t^2-t$ 和
  $\exp(\theta B_t-\theta^2t/2)$ 在每个有限时域上是真鞅；指出哪一步排除了“仅为局部鞅”。
  \item 设 $\lambda>0$、$\sigma\in\R$。把
  $Y_t=e^{\lambda t}X_t$ 代入 \Ito{} 公式，求解
  $\dd X_t=-\lambda X_t\dd t+\sigma\dd B_t$，并计算 $\E X_t$ 与 $\Var(X_t)$。
  \item 对 $a<x<b$，用 $u(y)=(y-a)/(b-a)$ 和连续时间可选抽样重证
  $\Prob_x(B_{\tau_{(a,b)}}=b)=(x-a)/(b-a)$；证明中须先说明退出时刻几乎必然有限。
  \item 设 $X,Y$ 是同一 Brownian motion 驱动、初值相同的全局 Lipschitz SDE 解。
  从积分方程出发推导
  $\E\sup_{s\le t}|X_s-Y_s|^2\le C\int_0^t
  \E\sup_{r\le s}|X_r-Y_r|^2\dd s$，再完成 \Gronwall{} 论证。
  \item 逐段验证例~\ref{ex:6-4-3-4} 的等待解，并求出爆炸解的最大存在区间；
  对照定理~\ref{theorem:6-4-3-2}，分别指出两个方程缺少哪一项假设。
\end{enumerate}
\end{exercise}

本节建立了 Brownian 驱动、全局 Lipschitz 系数下的基础随机微积分。一般半鞅积分、
弱解与分布唯一性、非 Lipschitz 扩散以及更广的 Feynman--Kac 正则性需要额外理论。在当前框架中，指数测试函数已给出
\[
 \E e^{i\xi\cdot B_t}=e^{-t|\xi|^2/2},
 \qquad
 P_tf_\xi=e^{-t|\xi|^2/2}f_\xi;
\]
这两个等式以特征函数和半群两种语言表示同一 Gaussian 因子。

